%%%%%%%%%%%%%%%%%%%%%%%%%%%%%%%%%%%%%%%%%%%%%%%%%%%%%%%%%%%%%%%%%%%%%%
% Lớp 11 | Các số đặc trưng đo xu thế trung tâm cho mẫu số liệu | Trung vị và tứ phân vị của mẫu số liệu ghép nhóm
% Dạng 1D5?2-2: Trung vị 
% Số câu: 158  (sinh tự động bởi tach_cau.py; từ nay thầy sửa trực tiếp file này)
%%%%%%%%%%%%%%%%%%%%%%%%%%%%%%%%%%%%%%%%%%%%%%%%%%%%%%%%%%%%%%%%%%%%%%

%%%%%%%%%%%%%%%% TRẮC NGHIỆM 4 PHƯƠNG ÁN %%%%%%%%%%%%%%%%

% nguồn: Ok_Toan11_TracNghiem_Lan2
\begin{ex}%[1D5N2-2]%[ma:324ce26b]
	Kết quả khảo sát cân nặng của $25$ quả bơ ở một lô  hàng cho trong bảng sau
	\begin{center}
		\begin{tabular}{|c|c|c|c|c|c|}
			\hline 
			 Cân nặng (g) & $[150; 155)$ & $[155; 160)$ & $[160; 165)$ & $[165; 170)$ & $[170; 175)$ \\
			\hline 
			Số quả bơ & $1$ & $7$ & $12$ & $3$ &  $2$ \\
			\hline 
		\end{tabular}
	\end{center}
	Trung vị của mẫu số liệu trên thuộc khoảng nào trong các khoảng  dưới đây?
	 \choice
		{$[170; 175)$} 
		{$[155;  160)$}
		{$[165; 170)$} 
		{\True $[160; 165)$} 
	 \loigiai{
		 Gọi $x_{1}$, $x_{2}$,$\ldots$, $x_{25}$ là cân nặng của $25$ quả bơ xếp theo thứ tự không giảm.\\
		 Do $x_{1} \in[150; 155)$; $x_{2}$,$\ldots$, $x_{8} \in[155; 160)$; $x_{9}$,$\ldots$, $x_{20} \in[160; 165)$ nên trung vị của mẫu số liệu $x_{1}$, $x_{2}$,$\ldots$, $x_{25}$ là $x_{13} \in[160; 165)$.
	} 
\end{ex}

% nguồn: Ok_Toan11_TracNghiem_Lan2
\begin{ex}%[1D5N2-2]%[ma:9e487d04]
	 Thời gian truy cập Internet mỗi buổi tối của một số học sinh được cho trong bảng sau
	\begin{center}
		\begin{tabular}{|c|c|c|c|c|c|}
			 \hline
			Thời gian (phút) & $[9{,}5; 12{,}5)$ & $[12{,}5; 15{,}5)$ & $[15{,}5; 18{,}5)$ & $[18{,}5; 21{,}5)$ &  $[21{,}5; 24{,}5)$ \\
			 \hline
			Số học sinh &  $3$ & $12$ & $15$ & $24$ & $2$ \\
			 \hline
		\end{tabular}
	\end{center}
	Tính trung vị của mẫu số liệu  ghép nhóm này.
	 \choice
		{\True $18{,}1$} 
		 {$18{,}5$}
		{$17{,}2$} 
		 {$15{,}6$}
	 \loigiai{
		Cỡ  mẫu là $n=3+12+15+24+2=56$.\\
		Gọi $x_{1}$,$\ldots$, $x_{56}$ là  thời gian vào internet của $56$ học sinh và giả sử dãy này được sắp xếp theo thứ tự tăng dần. Khi đó, trung vị là $\dfrac{x_{28}+x_{29}}{2}$.\\
		Do $2$ giá trị $x_{28}$, $x_{29}$ thuộc  nhóm $[15{,}5; 18{,}5)$ nên nhóm này chứa trung vị.\\
		Do  đó $p=3$, $a_{3}=15{,}5$, $m_{3}=15$, $m_{1}+m_{2}=3+12=15$, $a_{4}-a_{3}=3$ và ta có\\ \centerline{$M_{e}=15{,}5+\dfrac{\dfrac{56}{2}-15}{15} \cdot 3=18{,}1$.}
	 }
\end{ex}

% nguồn: Ok_Toan11_TracNghiem_Lan2
\begin{ex}%[1D5N2-2]%[ma:a6b250b8]
	Khảo sát thời gian chạy  bộ trong một ngày của một số học sinh khối 11 thu được mẫu số liệu ghép nhóm sau
	\begin{center}
		\begin{tabular}{|c|c|c|c|c|c|}
			\hline 
			Thời gian (phút) & $[0; 20)$ & $[20; 40)$ & $[40; 60)$ & $[60; 80)$ &  $[80; 100)$ \\
			 \hline
			Số học sinh & $5$ & $9$ &  $12$ & $10$ & $6$ \\
			 \hline
		\end{tabular}
	\end{center}
	 Nhóm chứa trung vị là
	\choice 
		{$[0; 20)$} 
		{$[20; 40)$} 
		{\True $[40;  60)$}
		{$[60;  80)$}
	 \loigiai{
		 $\dfrac{n}{2}=\dfrac{42}{2}=21$; $5+9+12=26$.\\
		Nhóm chứa  trung vị $[40; 60)$.
	} 
\end{ex}

% nguồn: Ok_Toan11_TracNghiem_Lan2
\begin{ex}%[1D5N2-2]%[ma:f477ac69]
	Trong tuẫn lễ bảo vệ môi trường, các học sinh khối $11$ tiến hành thu nhặt  vỏ chai nhựa để tái chế. Nhà trường thống kê kết quả thu nhặt vỏ chai của học sinh khối $11$ ở bảng sau
	\begin{center}
		\begin{tabular}{|c|c|c|c|c|c|}
			\hline 
			Số vỏ chai nhựa & $[11; 15]$ & $[16; 20]$ & $[21; 25]$ & $[26; 30]$ & $[31;  35]$ \\
			 \hline
			Số  học sinh & $53$ & $82$ & $48$ & $39$ & $18$ \\
			 \hline
		\end{tabular}
	\end{center}
	Hãy tìm trung vị  của mẫu số liệu ghép nhóm trên.
	 \choice
		 {$19{,}51$}
		 {\True $19{,}59$}
		 {$20{,}2$}
		{$18{,}6$} 
	 \loigiai{
		Số học sinh tham gia thu  nhặt vỏ chai nhựa là $n=53+82+48+39+18=240$.\\
		Gọi $x_{1}$, $x_{2}$,$\ldots$, $x_{240}$ lần lượt là số  vỏ chai nhựa $240$ học sinh khối $11$ thu nhặt được xếp theo thứ tự không giảm.\\
		Do $x_{1}$,$\ldots$, $x_{53} \in[10{,}5; 15{,}5)$; $x_{54}$,$\ldots$, $x_{135} \in[15,5; 20,5)$ nên trung vị của mẫu số liệu $x_{1}$, $x_{2}$,$\ldots$, $x_{240}$ là $\dfrac{1}{2}(x_{120}+x_{121})  \in[15{,}5; 20{,}5)$.\\
		Ta xác định được $n=240$, $n_{m}=82$, $C=53$, $u_{m}=15{,}5$,  $u_{m+1}=20{,}5$.\\
		 Vậy trung vị của mẫu số liệu ghép nhóm là $M_{e}=15{,}5+\dfrac{\dfrac{240}{2}-53}{82} \cdot(20{,}5-15{,}5)=\dfrac{803}{41} \approx 19{,}59$.
	} 
\end{ex}

% nguồn: Ok_Toan11_TracNghiem_Lan2
\begin{ex}%[1D5N2-2]%[ma:a2d33cc9]
	Khối lượng của $30$ củ khoai lang thu hoạch ở  $1$ hộ gia đình
	\begin{center}
		\begin{tabular}{|c|c|}
			\hline 
			Lớp  khối lượng (gam) & Tần số \\
			\hline 
			$[70;  80)$ & $3$ \\
			\hline 
			$[80;  90)$ & $6$ \\
			\hline 
			$[90; 100)$  & $12$ \\
			 \hline
			$[100; 110)$ & $6$  \\
			 \hline
			$[110;  120)$ & $3$ \\
			 \hline
		\end{tabular}
	\end{center}
	Trung vị của mẫu số liệu  ghép nhóm trên là
	\choice 
		{$90$} 
		 {\True $95$}
		 {$98$}
		 {$97$}
	\loigiai{ 
		Ta  có $Q_2=90+\dfrac{15-9}{12} \cdot 10=95$.
	} 
\end{ex}

% nguồn: Ok_Toan11_TracNghiem_Lan2
\begin{ex}%[1D5N2-2]%[ma:7860c22a]
	Doanh thu bán hàng trong $20$ ngày được  lựa chọn ngẫu nhiên của một cửa hàng được ghi lại ở bảng sau (đơn vị: triệu đồng).
	\begin{center}
		\begin{tabular}{|c|c|c|c|c|c|}
			 \hline
			Doanh thu & $[5; 7)$ & $[7; 9)$ & $[9; 11)$ & $[11; 13)$ & $[13; 15)$  \\
			 \hline
			Số ngày & $2$ & $7$  & $7$ & $3$ & $1$ \\
			\hline 
		\end{tabular}
	\end{center}
	Trung vị  của mẫu số liệu trên thuộc khoảng nào trong các khoảng dưới đây?
	\choice 
		 {$[7; 9)$}
		{\True $[9; 11)$} 
		{$[11; 13)$} 
		 {$[13; 15)$}
	\loigiai{ 
		 Gọi $x_{1}$, $x_{2}$,$\ldots$, $x_{20}$ là doanh thu bán hàng trong $20$ ngày xếp theo thứ tự không giảm.\\
		Khi đó: $x_{1}$, $x_{2} \in [5; 7)$; $x_{3}$,$\ldots$, $x_{9} \in[7; 9)$; $x_{10}$,$\ldots$, $x_{16} \in[9; 11)$; $x_{17}$,$\ldots$,  $x_{19} \in[11; 13)$; $x_{20} \in[13; 15)$.\\
		Do đó, trung vị của mẫu số liệu  thuộc nhóm $[9; 11)$.
	} 
\end{ex}

% nguồn: Ok_Toan11_TracNghiem_Lan2
\begin{ex}%[1D5N2-2]%[ma:d62653f3]
	Người ta tiến hành phỏng vấn $40$ người về một mẫu áo khoác. Người điều  tra yêu cầu cho điểm mẫu áo đó theo thang điểm là $100$. Kết quả được trình bày trong bảng ghép nhóm sau
	\begin{center}
		\begin{tabular}{|c|c|c|c|c|c|c|}
			\hline 
			Nhóm & $[50; 60)$ &  $[60; 70)$ & $[70; 80)$ & $[80; 90)$ & $[90; 100)$ & \\
			\hline 
			Tần số & $4$ & $5$ & $23$ & $6$ & $2$  & $n=40$ \\
			\hline 
		\end{tabular}
	\end{center}
	Trung vị của mẫu số liệu ghép nhóm trên gần nhất với giá  trị
	\choice 
		 {$74$}
		{\True $75$} 
		{$76$} 
		 {$77$}
	\loigiai{ 
		 $M_e=70+\dfrac{20-9}{23} \cdot 10=75$.
	} 
\end{ex}

% nguồn: Ok_TracNghiemLop11_Lan1
\begin{ex}%[1D5N2-2]%[ma:b81b4c52]
	Thống kê điểm học kì  môn toán của các học sinh lớp $11$A của một trường THPT, người ta thu được số liệu sau
	\begin{center}
		\begin{tabular}{|c|c|c|c|c|c|c|c|c|c|c|c|c|c|c|}
			 \hline
			$3$ & $5{,}5$ & $5$ & $4$ & $4{,}5$ & $4{,}5$ & $3$ & $5$ & $4$  & $4{,}5$ & $4{,}5$ & $6{,}5$ & $6{,}5$ & $7{,}5$ & $3{,}5$ \\
			 \hline
			$5$ & $6$ & $7$ & $8$ &  $8$ & $7$ & $4{,}5$ & $6$ & $5$ & $7$ & $4$ & $5{,}5$ & $7{,}5$ & $8{,}5$ & $9{,}5$ \\
			\hline 
			$4$ & $3{,}5$ & $5$ & $8{,}5$ & $6{,}5$ & $4{,}5$ & $7{,}5$ & $7$ & $4{,}5$  & $3$ & $7$ & $5{,}5$ & $5{,}5$ & $6{,}5$ & $9$ \\
			\hline 
		\end{tabular}
	\end{center}
	Tìm số trung vị của  mẫu số liệu khi ta ghép lớp thành các nhóm có độ dài là $1$ như sau
	\[\left[3; 4\right), \left[4; 5\right),\ldots, \left[9; 10\right)\] 
	 \choice
		{$6{,}7$} 
		 {$9{,}3$}
		{$5{,}8$} 
		 {\True $5{,}7$}
	 \loigiai{
		Ta có bảng tần số ghép lớp,  tần số tích lũy sau
		\begin{center}
			\begin{tabular}{|c|c|c|}
				 \hline
				Lớp & Tần số & Tần số tích lũy \\ 
				\hline 
				$\left[3; 4\right)$ & $5$ &  $5$ \\
				 \hline
				 $\left[4; 5\right)$ & $11$ & $16$ \\
				\hline 
				$\left[5; 6\right)$ & $9$ & $25$  \\
				 \hline
				$\left[6;  7\right)$ & $6$ & $31$ \\
				 \hline
				$\left[7; 8\right)$ & $8$  & $39$ \\
				 \hline
				$\left[8; 9\right)$ & $4$  & $43$ \\
				 \hline
				$\left[9; 10\right)$  & $2$ & $45$ \\
				\hline 
				& $n=45$  & \\
				 \hline
			\end{tabular}
		\end{center}
		Ta có số phần tử của mẫu  là\\
		$n=45\Rightarrow \dfrac{n}{2}=22{,}5$.\\ 
		Mà  $cf_2=16< 22{,}5< cf_3=25$ suy ra nhóm $3$ là nhóm đầu tiên có tần số tích lũy lớn hơn hoặc bằng $22{,}5$.\\
		Xét nhóm $3$ là nhóm $\left[5; 6\right)$ có $r=5$; $d=1$; $n_3=9$ và nhóm $2$ là nhóm $\left[4; 5\right)$ có $cf_2=16$.\\ 
		Áp dụng công thức ta có trung  vị của mẫu số liệu là
		$$M_e=5+\left(\dfrac{22{,}5-16}{9} \right)\cdot 1\approx  5{,}7.$$
	 }
\end{ex}

% nguồn: Ok_TracNghiemLop11_Lan1
\begin{ex}%[1D5N2-2]%[ma:b58d02cb]
	Cho bảng phân  phối tần số thống kê số tiền (đơn vị: nghìn đồng) của $60$ khách hàng trong $1$ ngày ở một cửa hàng như sau
	\begin{center}
		\begin{tabular}{|c|c|c|c|c|c|}
			 \hline
			Khoảng & $[40;49]$ & $[50;59]$  & $[60;69]$ & $[70;79]$ & $[80;89]$\\
			\hline 
			Tần số &  $4$ & $7$ & $17$ & $22$ & $10$\\
			 \hline
		\end{tabular}
	\end{center}
	Trung vị của mẫu số  liệu ghép nhóm này nằm ở khoảng nào?
	\choice 
		{$[60;69)$} 
		 {$[50;59)$}
		{\True  $[70;79)$}
		{$[80;89)$} 
	 \loigiai{
		Cỡ mẫu là  $60$.\\
		Do vậy số trung vị sẽ ở phần tử khoảng thứ $[30;31]$.\\ 
		Khi  đó, xét theo bảng số liệu ghép nhóm trên, khoảng $[70;79)$ chứa các phần tử từ $[29;40]$, do đó trung vị nằm ở khoảng $[70;79)$.
	 }
\end{ex}

% nguồn: Ok_TracNghiemLop11_Lan1
\begin{ex}%[1D5N2-2]%[ma:173ac693]
	Một học viện bóng đá điều tra về lứa tuổi của $85$  học viên trẻ đăng kí đầu tiên để tham gia khóa học mới và thu được bảng sau:
	\begin{longtable}{|c|c|c|c|c|}
		\hline 
		Nhóm  tuổi&$[8;10)$&$[10;12)$&$[12;14)$&$[14;16)$\\
		 \hline
		Số  học viên&$14$&$20$&$33$&$18$\\
		\hline 
	\end{longtable}
	Nhóm chứa trung vị của mẫu số liệu ghép nhóm trên là 
	 \choice
		{$[8;10)$} 
		 {$[10;12)$}
		 {\True $[12;14)$}
		 {$[14;16)$}
	 \loigiai{
		Cỡ mẫu  $n=85$.\\
		Gọi $x_1,x_2,\ldots,x_{85}$ lần  lượt là số tuổi của $85$ học viên xếp theo thứ tự không giảm.\\
		Do đó trung vị của mẫu  số liệu $x_1,x_2,\ldots,x_{85}$ là $M_e=x_{43}\in [12;14)$.
	 }
\end{ex}

% nguồn: Ok_TracNghiemLop11_Lan1
\begin{ex}%[1D5N2-2]%[ma:7165e8f7]
	Cho dãy số liệu thống kê: $50$, $48$, $34$, $36$, $56$, $35$, $43$, $38$, $55$.  Số trung vị là
	\choice 
		 {$38$}
		{$35$} 
		 {\True $43$}
		{$55$} 
	\loigiai{ 
		Dãy số liệu thống kê xếp thành  dãy không giảm là $34$, $35$, $36$, $38$, $43$, $48$, $50$, $55$, $56$.\\
		Số trung vị là $M_e=43$. 
	} 
\end{ex}

% nguồn: Ok_TracNghiemLop11_Lan1
\begin{ex}%[1D5N2-2]%[ma:773cb63c]
	Doanh thu bán hàng trong $20$ ngày được lựa chọn ngẫu nhiên của một  của hàng được ghi lại ở bảng sau (đơn vị: triệu đồng):
	\begin{center}
		\begin{tabular}{|c|c|c|c|c|c|}
			 \hline
			 Doanh thu & $\left[5;7\right)$ & $\left[7;9\right)$ &$\left[9;11\right)$ &$\left[11;13\right)$ &$\left[13;15\right)$ \\
			\hline 
			Số ngày & $2$ &$7$ &$7$  &$3$ &$1$ \\
			 \hline
		\end{tabular}
	\end{center}
	Trung vị của mẫu số liệu trên thuộc khoảng nào trong các  khoảng dưới đây?
	 \choice
		{$\left[7;9\right)$} 
		 {\True $\left[9;11\right)$}
		{$\left[11;13\right)$} 
		 {$\left[13;15\right)$}
	\loigiai{ 
		Gọi $x_1,x_2,\ldots,x_{20}$ là  doanh thu bán hàng trong $20$ ngày xếp theo thứ tự không giảm.\\
		Khi đó $x_1,x_2\in \left[5;7\right)$; $x_3,\ldots,  x_9\in \left[7;9\right)$; $x_{10},\ldots, x_{16}\in \left[9;11\right)$; $x_{17},\ldots, x_{19}\in \left[11;13\right)$; $x_{20}\in \left[13;15\right)$.\\
		Do đó, trung vị của mẫu số  liệu thuộc nhóm $\left[9;11\right)$.
	 }
\end{ex}

% nguồn: Ok_TracNghiemLop11_Lan1
\begin{ex}%[1D5N2-2]%[ma:56cb7f16]
	Khảo sát thời gian chạy bộ trong một ngày của một số học sinh khối $11$ thu được mẫu số liệu ghép  nhóm sau
	\begin{center}
		{\tabcolsep = 2mm 
		\begin{tabular}{|c|c|c|c|c|c|}
			\hline 
			Thời gian (phút) & $[0; 20)$ & $[20; 40)$  & $[40; 60)$ & $[60; 80)$ & $[80; 100)$ \\
			 \hline
			Số  học sinh & $5$ & $9$ & $12$ & $10$ & $6$\\
			 \hline
		\end{tabular}
		 }
	\end{center}
	Nhóm  chứa trung vị của mẫu số liệu trên là
	\choice 
		{\True $[40; 60)$} 
		 {$[20; 40)$}
		{$[60; 80)$} 
		{$[80; 100)$} 
	 \loigiai{
		Ta có $n=42$  nên trung vị của mẫu số liệu trên là $\mathrm{M}_e = \mathrm{Q}_2 = \dfrac{x_{21}+ x_{22}}{2}$.\\
		Mà  $x_{21}, x_{22} \in [40; 60)$.\\
		Vậy nhóm chứa  trung vị của mẫu số liệu trên là nhóm $[40;60)$.
	} 
\end{ex}

% nguồn: Ok_TracNghiemLop11_Lan1
\begin{ex}%[1D5N2-2]%[ma:6593f1a9]
	Doanh thu bán hàng trong $20$ ngày được lựa chọn ngẫu nhiên của một của hàng được ghi lại ở bảng sau  (đơn vị: triệu đồng)
	\begin{center}
		{\tabcolsep = 2mm 
		\begin{tabular}{|c|c|c|c|c|c|}
			 \hline
			Doanh thu & $[5; 7)$ & $[7; 9)$ & $[9;  11)$ & $[11; 13)$ & $[13; 15)$ \\
			 \hline
			Số ngày & $2$ & $7$ & $7$ & $3$ & $1$\\ 
			\hline 
		\end{tabular}
		 }
	\end{center}
	Trung vị của mẫu số  liệu trên thuộc khoảng nào trong các khoảng dưới đây?
	 \choice
		{$[7;  9)$}
		{\True $[9; 11)$} 
		{$[11; 13)$} 
		{$[13; 15)$} 
	\loigiai{ 
	Gọi $x_1$, $x_2$, \ldots, $x_{20}$  là doanh thu bán hàng trong $20$ ngày xếp theo thứ tự không giảm.\\
	 Khi đó
	$x_1, x_2 \in [5;  7)$;
	$x_3, \cdots,  x_9 \in [7; 9)$;
	$x_{10}, \cdots, x_{16} \in [9; 11)$; 
	$x_{17}, x_{18}, x_{19} \in [11;  13)$;
	$x_{20}  \in [13; 15)$.\\
	Do đó, trung vị của mẫu số liệu  thuộc nhóm $[9; 11)$.
	 }
\end{ex}

% nguồn: Ok_TracNghiemLop11_Lan1
\begin{ex}%[1D5N2-2]%[ma:d60d60fa]
	Thời gian đọc sách mỗi ngày của một số học sinh được  cho trong bảng sau
	\begin{center}
		\begin{tabular}{|c|c|c|c|c|c|}
			\hline  Thời gian (phút) & {$[10; 15)$} & {$[15; 20)$} & {$[20; 25)$} & {$[25; 30)$} & {$[30; 35)$} \\
			\hline Số học sinh  & $3$ & $10$ & $12$ & $15$ & $20$ \\
			\hline 
		\end{tabular}
	\end{center}
	Xác định trung vị của mẫu số liệu  ghép nhóm
	 \choice
		{$25$} 
		{$26$} 
		 {$25{,}56$}
		{\True $26{,}67$} 
	\loigiai{ 
		 Cỡ mẫu là $n=3+10+12+15+20=60$.\\
		 Ta có $\dfrac{n}{2}=30$.\\
		Mà $cf_3=25< 30< cf_4=40$. Suy ra nhóm $4$ là nhóm đầu tiên  có tần số tích lũy lớn hơn hoặc bằng $30$.\\
		Xét nhóm $4$ là nhóm $\left[25;30\right)$ có $r=25$; $d=5$, $n_4=15$; và nhóm $3$  có $cf_3=25$.\\
		Vậy trung vị của mẫu số  liệu là
		\[M_e=25+\dfrac{30-25}{15}\cdot 5\approx  26{,}67.\]
	} 
\end{ex}

% nguồn: Ok_TracNghiemLop11_Lan1
\begin{ex}%[1D5N2-2]%[ma:9c6856cb]
	Dưới đây là một mẫu số liệu cho ở  dạng bảng tần số ghép nhóm
	\begin{center}
		\begin{tabular}{|c|c|c|c|}
			 \hline
			Nhóm ghép & Tần số & Nhóm ghép &  Tần số \\
			 \hline
			$\left[62{,}5;67{,}5\right)$ & $4$ &  $\left[82{,}5;87{,}5\right)$ & $22$ \\
			\hline 
			$\left[67{,}5;72{,}5\right)$ & $7$ &  $\left[87{,}5;92{,}5\right)$ & $5$ \\
			\hline 
			$\left[72{,}5;77{,}5\right)$ & $10$ & $\left[92{,}5;97{,}5\right)$ &  $10$ \\
			\hline 
			$\left[77{,}5;82{,}5\right)$ & $26$ &  $\left[97{,}5;102{,}5\right)$ & $16$ \\
			\hline 
		\end{tabular}
	\end{center}
	 Khẳng định nào sau đây là đúng
	 \choice
		{$77{,}5\le  M_e < 82{,}5$}
		{\True $82{,}5\le M_e  < 87{,}5$}
		 {$87{,}5\le M_e < 92{,}5$}
		{$92{,}5\le  M_e < 97{,}5$}
	\loigiai{ 
		Ta có $\dfrac{N}{2}=\dfrac{100}{2}=50$.\\ 
		Mặt khác $4+7+10+26+22=69$. Vậy nhóm chứa trung vị là $\left[82{,}5;87{,}5\right)$. 
	 }
\end{ex}

% nguồn: Ok_TracNghiemLop11_Lan1
\begin{ex}%[1D5N2-2]%[ma:e098fd1a]
	Thống kê điểm kiểm tra một tiết môn Toán của lớp 11A của trường THPT  Nguyễn Huệ được ghi lại như sau
	\begin{center}
		\begin{tabular}{|c|c|c|c|c|c|c|c|}
			\hline Giá trị $(x)$ & $4$ & $5$ & $6$ & $7$ & $8$ & $9$ & $10$  \\
			\hline Tần số $(n)$ & $2$ & $3$ & $7$ & $6$ &  $12$ & $4$ & $2$ \\
			\hline 
		\end{tabular}
	\end{center}
	Số trung vị của  mẫu số liệu trên là
	\choice 
		{$8{,}0$} 
		{\True $7{,}5$} 
		{$7{,}8$} 
		{$8{,}5$} 
	 \loigiai{
		Các số liệu đã được xếp theo thứ tự tăng dần.\\ 
		 Cỡ mẫu là $n=2+3+7+8+10+4+2=36$.\\
	Khi  đó trung vị là $$M_e=\dfrac{x_{18}+x_{19}}{2}=\dfrac{7+8}{2}=7{,}5.$$}
\end{ex}

% nguồn: Ok_TracNghiemLop11_Lan1
\begin{ex}%[1D5N2-2]%[ma:10e1bb05]
	Khảo sát thời gian tập thể dục trong ngày của  một số học sinh khối 11 thu được mẫu số liệu ghép nhóm sau
	\begin{center}
		\begin{tabular}{|c|c|c|c|c|c|}
			\hline Thời gian (phút) & {$[0; 15)$} & {$[15; 30)$} & {$[30;  45)$} & {$[45; 60)$} & {$[60; 75)$} \\
			\hline Số học sinh & 9 & 5 & 15 & 14 & 7 \\ 
			 \hline
		\end{tabular}
	\end{center}
	Nhóm  chứa trung vị là
	 \choice
		{$[15;  30)$}
		{\True $[30; 45)$} 
		{$[60;  75)$}
		{$[45;  60)$}
	\loigiai{ 
		 Cỡ mẫu là $n=9+5+15+14+7=50$.\\
		Gọi $x_1, \ldots, x_{50}$ là mẫu số liệu được sắp xếp theo thứ tự không giảm.\\ Khi đó, trung  vị là $\dfrac{1}{2}\left(x_{25}+x_{26}\right)$ thuộc nhóm $[30; 45)$.
	 }
\end{ex}

% nguồn: Ok_TracNghiemLop11_Lan1
\begin{ex}%[1D5N2-2]%[ma:427aafb7]
	Thời gian  (phút) xem tivi mỗi buổi tối của một số học sinh được cho trong bảng sau:
	\begin{center}
		\begin{tabular}{|c|c|c|c|c|c|c|c|}
			\hline \begin{tabular}{c}
			TG  \\
			 (phút)
		\end{tabular} & {$[6{,}5; 9{,}5)$} & {$[9{,}5; 12{,}5)$} & {$[12{,}5; 15{,}5)$} & {$[15{,}5; 18{,}5)$} & {$[18{,}5; 21{,}5)$} & {$[21{,}5; 24{,}5)$} & {$[24{,}5; 27{,}5)$} \\
		\hline Số HS &  $2$ & $3$ & $12$ & $15$ & $24$ & $2$ & $2$ \\
		\hline 
	\end{tabular}
\end{center}
Số trung vị của mẫu số liệu ghép  nhóm này là
\choice  {\True $18{,}1$}{$15{,}1$}{$21{,}1$ }{$15$}
\loigiai{Cỡ  mẫu: $n=2+3+12+15+24+2+2=60$.\\
	Nhóm chứa trung vị:  $[15{,}5; 18{,}5)$.\\
	Suy  ra: $u_m=15{,}5$ và $u_{m+1}=18{,}5$.\\
	Tần số của nhóm chứa trung  vị: $n_m=15$.
	$$C=n_1+n_2+n_3=2+3+12=17$$. 
	Vậy trung vị của mẫu số  liệu ghép nhóm là $$M_e=15{,}5+\dfrac{\dfrac{60}{2}-17}{15} \cdot(18{,}5-15{,}5)=18{,}1.$$
 }
\end{ex}

% nguồn: Ok_Toan11_TracNghiem_Lan2
\begin{ex}%[1D5H2-2]%[ma:269b0ac2]
	Cho bảng số liệu đo chiều  cao của $40$ học sinh nam tại một trường THPT như sau:
	\begin{center}
		\begin{tabular}{|c|c|c|c|c|c|c|}
			 \hline
			Chiều cao (cm) & $[141;147)$ & $[147;153)$ & $[153;159)$ & $[159;165)$ & $[165;171)$ & $[171;177)$  \\
			 \hline
			Số học sinh  & $6$ & $4$ & $2$ & $6$ & $10$ & $12$ \\
			\hline 
		\end{tabular}
	\end{center}
	Số trung vị của mẫu  số liệu thuộc nhóm nào?
	\choice 
		{$\left[153;159\right)$} 
		{$\left[159;165\right)$} 
		{$\left[147;153\right)$} 
		 {\True $\left[165;171\right)$}
	 \loigiai{
		Gọi $x_1; x_2;\ldots;x_{40}$ là chiều của của 40 học sinh  nam xếp theo thứ tự không giảm.\\
		Từ bảng số liệu trên ta thấy $x_1$; $x_2$; $\ldots$; $x_6 \in \left[141;147\right)$; $x_7$; $x_8$; $\ldots$; $x_{10} \in \left[147;153\right)$; $x_{11}$; $x_{12} \in \left[153;159\right)$; $x_{13}$; $x_{14}$; $\ldots$; $x_{18} \in \left[159;165\right)$; $x_{19}$; $ x_{20}$; $\ldots$; $x_{28} \in \left[165;171\right)$; $x_{29}$; $x_{30}$; $\ldots$; $x_{40} \in \left[171;177\right)$ nên số trung vị của mẫu số liệu $x_1$; $x_2$; $\ldots$; $x_{40}$ là $\dfrac{1}{2} \left(x_{20}+x_{21} \right)$ thuộc  nhóm $\left[165;171\right)$.
	 }
\end{ex}

% nguồn: Ok_Toan11_TracNghiem_Lan2
\begin{ex}%[1D5H2-2]%[ma:ada0e71f]
	Cân nặng của các em học sinh nam lớp $11A$  được thống kê ở bảng sau
	\begin{center}
		\begin{tabular}{|c|c|c|c|c|c|c|}
			 \hline
			Cân nặng & $[45;49)$  & $[49;53)$ & $[53;57)$ & $[57;61)$ & $[61;65)$ \\
			\hline 
			Số học sinh & $4$ & $5$  & $7$ & $7$ & $5$ \\
			\hline 
		\end{tabular}
	\end{center}
	Trung vị của mẫu số  liệu trên là
	\choice 
		{$55{,}85$} 
		 {$55{,}87$}
		 {$53{,}86$}
		{\True  $55{,}86$}
	 \loigiai{
		Gọi $x_1$, $x_2$, {\dots}, $x_{28}$ là cân nặng của  các em học sinh lớp $11A$ xếp theo thứ tự không giảm.\\
		Do $x_1$, $x_2$, $x_3$, $x_4 \in \left[45;49\right)$; $x_5$, $\ldots$, $x_9 \in \left[49;53\right)$; $x_{10}$, $\ldots$, $x_{16} \in \left[53;57\right)$; $x_{17}$, $\ldots$, $x_{23} \in \left[57;61\right)$; $x_{24}$, $\ldots$, $x_{28} \in \left[61;65\right)$ nên trung vị của mẫu số liệu $x_1$, $x_2$, $\ldots$, $x_{28}$ là $\dfrac{1}{2}  \left(x_{14}+x_{15} \right)\in \left[53;57\right)$.\\
		Ta xác định  được $n=28$, $n_m=7$, $C=4+5=9$, $u_m=53$, $u_{m+1}=57$. Vậy trung vị của mẫu số liệu ghép nhóm là
		\[M_e=53+\dfrac{\dfrac{28}{2}-9}{7} \left(57-53\right)\approx 55{,}86.\] 
	 }
\end{ex}

% nguồn: Ok_Toan11_TracNghiem_Lan2
\begin{ex}%[1D5H2-2]%[ma:820ef564]
	Số lượng huy chương vàng  tại Sea Games $31$ được thống kê trong bảng sau:
	\begin{center}
		\begin{tabular}{|c|c|c|c|c|}
			\hline 
			Số huy chương & $\left[0;10\right)$ & $\left[10;50\right)$ & $\left[50;100\right)$ &  $\left[100;210\right)$ \\
			 \hline
			Quốc gia & $5$ & $2$  & $3$ & $1$ \\
			\hline 
		\end{tabular}
	\end{center}
	Xác định trung vị của mẫu  số liệu ghép nhóm trên
	\choice 
		 {$10$}
		{$15$} 
		{\True $20$} 
		 {$30$}
	 \loigiai{
		\begin{itemize}
			\item  Cỡ mẫu $n=11$.
			\item Gọi $x_1$, $x_2$, $\cdots$, $x_{11}$ lần lượt là số huy chương  nhận được của các quốc gia được sắp xếp theo thứ tự không giảm.
			\item  $M_e=Q_2=x_6\in\left[10; 20\right)$. Suy ra
			\[M_e=Q_2=10+\dfrac{2\cdot\dfrac{11}{4}-5}{2}\cdot\left(50-10\right)=20.\] 
		\end{itemize}
	} 
\end{ex}

% nguồn: Ok_Toan11_TracNghiem_Lan2
\begin{ex}%[1D5H2-2]%[ma:60a7a8f3]
	Thống kê về thời lượng mỗi trận đấu bi-a trong vòng tứ kết giải đấu European Open người ta được mẫu số liệu ghép nhóm như  sau:
	\begin{center}
		\begin{tabular}{|c|c|c|c|c|c|c|}
			 \hline
			Thời gian & $[9,5;12,5)$ & $[12,5;15,5)$ & $[15,5;18,5)$ & $[18,5;21,5)$ & $[21,5;24,5)$  \\
			 \hline
			Số trận & $3$ & $12$ & $15$ &  $24$ & $2$ \\
			 \hline
		\end{tabular}
	\end{center}
	Số trung vị của  mẫu số liệu trên là
	\choice 
		 {$18$}
		 {\True $18{,}1$}
		{$18{,}2$} 
		 {$18{,}3$}
	 \loigiai{
		Ta có: Số phần tử của  mẫu là $n=56\Rightarrow \dfrac{n}{2}=28$.\\
		Suy ra  nhóm $\left[15{,}5; 18{,}5\right)$ chứa số trung vị.\\
		Do đó số trung vị  của mẫu số liệu trên là
		\[M_e=15{,}5+\dfrac{28-15}{15}\cdot  3=18{,}1.\]
	} 
\end{ex}

% nguồn: Ok_Toan11_TracNghiem_Lan2
\begin{ex}%[1D5H2-2]%[ma:6ca85e82]
	Cho mẫu số liệu ghép nhóm về chiều cao của $25$ cây dừa giống như sau: 
	\begin{center}
		\begin{tabular}{|c|c|c|c|c|c|c|}
			 \hline
			Chiều cao (cm) &  $[0;10)$ & $[10;20)$ & $[20;30)$ & $[30;40)$ & $[40;50)$ \\
			\hline 
			Số cây & $4$ & $6$ & $7$ &  $5$ & $3$ \\
			 \hline
		\end{tabular}
	\end{center}
	Trung  vị của mẫu số liệu ghép nhóm này là
	 \choice
		{$M_e=\dfrac{175}{7}$} 
		{$M_e=\dfrac{165}{5}$} 
		{\True  $M_e=\dfrac{165}{7}$}
		 {$M_e=\dfrac{165}{3}$}
	 \loigiai{
		 Cỡ mẫu: $n=4+6+7+5+3=25$.
		$x_1$, $x_2$, $\ldots$, $x_{25}$ là chiều cao của $25$ cây dừa giống được sắp xếp theo thứ tự không giảm. Khi  đó, trung vị là $x_{13}$. Do $x_{13}$ thuộc nhóm $\left[20;30\right)$ nên nhóm này chứa trung vị. Do đó:
		$p=3$, $a_3=20$,  $m_3=7$, $m_1+m_2=10$, $a_4-a_3=10$. Do đó:
		 \[M_e=20+\dfrac{\dfrac{25}{2}-10}{7}\cdot 10=\dfrac{165}{7}.\]
	 }
\end{ex}

% nguồn: Ok_Toan11_TracNghiem_Lan2
\begin{ex}%[1D5H2-2]%[ma:507957bb]
	Doanh thu bán  hàng trong $20$ ngày được lựa chọn ngẫu nhiên của một của hàng được ghi lại ở bảng sau:
	\begin{center}
		\begin{tabular}{|c|c|c|c|c|c|c|}
			\hline 
			Doanh thu & $[5;7)$ & $[7;9)$ &  $[9;11)$ & $[11;13)$ & $[13;15)$ \\
			 \hline
			Số ngày & $2$ &  $7$ & $7$ & $3$ & $1$ \\
			 \hline
		\end{tabular}
	\end{center}
	Trung vị của mẫu  số liệu trên là?
	\choice 
		 {$6{,}286$}
		{\True  $9{,}286$}
		 {$9{,}386$}
		{$6{,}386$} 
	\loigiai{ 
		Goi $x_1$, $ x_2$, $\ldots$, $ x_{20}$ là doanh  thu bán hàng trong 20 ngày xếp theo thứ tự không giảm.\\
		Khi đó: $x_1$, $x_2 \in \left[5;7\right)$, $x_3$, $\ldots$, $x_9 \in \left[7; \; 9\right)$, $x_9$, $\ldots$, $x_{16} \in \left[9; \;  11\right)$ $x_{17}$, $\ldots$, $x_{19} \in \left[11; \; 13\right)$, $x_{20} \in \left[13; \; 15\right)$.\\
		Do đó, trung  vị của mẫu số liệu thuộc nhóm $\left[9; \; 11\right)$.\\
		Ta xác định được $n=20$, $n_m=7$, $C=2+7=9$, $U_m=9$,  $U_{m+1}=11$.\\
		Vậy trung  vị của mẫu ghép nhóm là
		 \[M_e=9+\dfrac{\dfrac{20}{2}-9}{7}\cdot(11-9)=9{,}286\]
	} 
\end{ex}

% nguồn: Ok_Toan11_TracNghiem_Lan2
\begin{ex}%[1D5H2-2]%[ma:5a8584f8]
	Thời gian đọc sách mỗi ngày của một số học sinh được cho trong bảng sau 
	\begin{center}
		\begin{tabular}{|c|c|c|c|c|c|c|}
			 \hline
			Thời gian (phút)  & $[10;15)$ & $[15;20)$ & $[20;25)$ & $[25;30)$ & $[30;35)$ \\
			\hline 
			Số học sinh & $3$ & $10$ & $12$ & $15$ &  $20$ \\
			 \hline
		\end{tabular}
	\end{center}
	Xác định trung  vị của mẫu số liệu ghép nhóm.
	 \choice
		{$25$} 
		 {$26$}
		 {$25{,}56$}
		{\True $26{,}67$} 
	 \loigiai{
		Cỡ mẫu là  $n=3+10+12+15+20=60$.\\
		Ta  có: $\dfrac{n}{2}=30$.\\
		Mà $cf_3=25< 30< cf_4=40$. Suy ra nhóm $4$ là nhóm  đầu tiên có tần số tích lũy lớn hơn hoặc bằng $30$.\\
		Xét nhóm $4$ là nhóm $\left[25; 30\right)$ có $r=25$; $d=5$; $n_4=15$; và nhóm $3$ có $cf_3=25$.\\ 
		Vậy  trung vị của mẫu số liệu là:
		\[M_e=25+\dfrac{30-25}{15}\cdot5\approx 26{,}67.\] 
	 }
\end{ex}

% nguồn: Ok_Toan11_TracNghiem_Lan2
\begin{ex}%[1D5H2-2]%[ma:ef111679]
	Cân nặng của học sinh khối $10$ của trường Ngô Quyền vào đầu năm được thống  kê như bảng dưới đây:
	\begin{center}
		\begin{tabular}{|c|c|c|c|c|c|c|}
			 \hline
			Cân nặng & $[37,5;42,5)$ & $[42,5;47,5)$ & $[47,5;52,5)$ & $[52,5;57,5)$ &  $[57,5;62,5)$ & $[62,5;67,5)$ \\
			\hline 
			Số  học sinh & $27$ & $123$ & $150$ & $272$ & $45$ & $10$ \\
			 \hline
		\end{tabular}
	\end{center}
	Trung vị của mẫu số liệu trên gần với kết quả  nào sau dưới đây
	 \choice
		{\True  $52{,}75$}
		{$53{,}75$} 
		 {$54{,}24$}
		{$53{,}18$} 
	\loigiai{ 
		Số học sinh  khối $10$ của trường Ngô Quyền là $n=27+123+150+272+45+10=627$.\\
		Gọi $x_1$; $x_2$; $\ldots$; $x_{627}$ là cân nặng của $627$ học sinh khối $10$ xếp theo thứ tự không giảm. Do $x_1$, $\ldots$, $x_{27} \in \left[37{,}5;42{,}5\right)$; $x_{28}$; $\ldots$; $x_{150} \in \left[42{,}5;47{,}5\right)$; $x_{151}$; $\ldots$; $x_{300} \in \left[47{,}5;52{,}5\right)$; $x_{301}$; $\ldots$; $x_{572} \in \left[52{,}5;57{,}5\right)$; $x_{573}$; $\ldots$; $x_{617} \in \left[57{,}5;62{,}5\right)$; $x_{618}$;  $\ldots$; $x_{627} \in \left[62{,}5;67{,}5\right)$ nên trung vị của mẫu số liệu $x_1;x_2;\ldots;x_{627}$ là $x_{314} \in \left[52,5;57,5\right)$.\\
		Ta xác định được $n=627$, $n_m=272$, $C=27+123+150=300$, $u_m=52{,}5$,  $u_{m+1}=57{,}5$.\\
		Vậy trung vị của mẫu số liệu  ghép nhóm là
		\[M_e=u_m+\dfrac{\dfrac{n}{2}-C}{n_m}\cdot \left(u_{m+1}-u_m \right)=52{,}5+\dfrac{\dfrac{627}{2}-300}{272}\cdot  \left(57{,}5-52{,}5\right)=\dfrac{28695}{544} \approx 52{,}75.\]
	} 
\end{ex}

% nguồn: Ok_Toan11_TracNghiem_Lan2
\begin{ex}%[1D5H2-2]%[ma:96eaff36]
	Khảo sát thời gian tập thể dục trong ngày của một số học sinh khối 11 thu được mẫu số liệu ghép nhóm sau: 
	\begin{center}
		\begin{tabular}{|c|c|c|c|c|c|c|}
			\hline 
			Thời gian & $[0;20)$ & $[20;40)$ & $[40;60)$ &  $[60;80)$ & $[80;100)$ \\
			\hline 
			Số học sinh & $5$ & $9$ & $12$ & $10$ & $6$  \\
			\hline 
		\end{tabular}
	\end{center}
	Tần số của  nhóm chứa trung vị là
	 \choice
		{$5$} 
		{$9$} 
		 {\True $12$}
		 {$6$}
	 \loigiai{
		 Cỡ mẫu $n=42$.\\
		Trung vị $M_e$ là $\dfrac{x_{21}+x_{22}}{2}$. Do $x_{21}$; $x_{22}$ đều thuộc nhóm $[40;60)$ nên nhóm này chứa trung  vị.\\
		Tần số  của nhóm chứa trung vị là $12$
	} 
\end{ex}

% nguồn: Ok_Toan11_TracNghiem_Lan2
\begin{ex}%[1D5H2-2]%[ma:f9db93a0]
	Cho mẫu số liệu ghép nhóm về thống kê điểm số của $100$  học sinh tham dự kỳ thi học sinh giỏi toán, ta có bảng số liệu sau:
	\begin{center}
		\begin{tabular}{|c|c|c|c|c|c|c|c|}
			 \hline
			Điểm &  $[8;10)$ & $[10;12)$ & $[12;14)$ & $[14;16)$ & $[16;18)$ & $[18;20)$ \\
			 \hline
			Số học sinh & $6$ & $21$ & $30$ & $25$ & $14$ &  $4$ \\
			\hline 
		\end{tabular}
	\end{center}
	Tìm  trung vị của mẫu số liệu ghép nhóm trên.
	 \choice
		{$12{,}18$} 
		 {$12{,}81$}
		{$13{,}35$} 
		{\True $13{,}53$} 
	 \loigiai{
		 Cở mẫu là $n=6+21+30+25+14+4=100$.\\
		Gọi $x_1$; $x_2$; $\ldots$; $x_{100}$ là số điểm của $100$ học sinh và được sắp xếp theo thứ tự tăng dần. Khi đó, trung vị là $\dfrac{x_{50}+x_{51}}{2}$, do hai giá trị $x_{50}$, $x_{51}$ thuộc  nhóm $\left[12;14\right)$ nên nhóm này chứa trung vị.\\
		Ta có $p=3$; $a_3=12$; $m_3=30$; $m_1+m_2=25$; $a_4-a_3=2$ khi đó trung  vị của mẫu số liệu là $M_e=12+\dfrac{\dfrac{100}{2}-27}{30}\cdot2=13{,}53$.
	} 
\end{ex}

% nguồn: Ok_Toan11_TracNghiem_Lan2
\begin{ex}%[1D5H2-2]%[ma:038ae3da]
	Cho mẫu số liệu ghép nhóm  về thống kê huyết áp của $20$ người, ta có bảng số liệu sau:
	\begin{center}
		\begin{tabular}{|c|c|c|c|c|c|c|c|}
			\hline 
			Huyết áp (mmHg) & $[70;80)$ & $[80;90)$ & $[90;100)$ & $[100;110)$  & $[110;120)$ & $[120;130)$ \\
			\hline 
			Số người  & $4$ & $2$ & $3$ & $6$ & $3$ & $2$ \\
			\hline 
		\end{tabular}
	\end{center}
	Tìm trung vị  của mẫu số liệu ghép nhóm trên.
	 \choice
		 {$100{,}67$}
		{\True $101{,}67$} 
		 {$101{,}76$}
		{$100{,}76$} 
	\loigiai{ 
		Cỡ mẫu là  $n=4+2+3+6+3+2=20$.\\
		Gọi $x_1$; $x_2$; $\ldots$; $x_{20}$ là giá trị huyết áp của $20$ người và được sắp xếp theo thứ  tự tăng dần. Khi đó, trung vị là $\dfrac{x_{10}+x_{11}}{2}$, do hai giá trị $x_{10}$, $x_{11}$ thuộc nhóm $\left[100;110\right)$ nên nhóm này chứa trung vị.\\
		Ta có  $p=4$; $a_4=100$; $m_4=6$; $m_1+m_2+m_3=9$; $a_5-a_4=10$ khi đó trung vị của mẫu số liệu là $M_e=100+\dfrac{\dfrac{20}{2}-9}{6}\cdot 10=101{,}67$.
	 }
\end{ex}

% nguồn: Ok_Toan11_TracNghiem_Lan2
\begin{ex}%[1D5H2-2]%[ma:cf75bd8f]
	Khảo sát về cân nặng của  các học sinh lớp 11D3 người ta được một mẫu dữ liệu ghép nhóm như sau:
	\begin{center}
		\begin{tabular}{|c|c|c|c|c|c|c|c|}
			 \hline
			 Cân nặng & $[30;40)$ & $[40;50)$ & $[50;60)$ & $[60;70)$ & $[70;80)$ & $[80;90)$ \\
			 \hline
			Số học sinh & $2$ & $10$ & $16$ & $8$ & $2$ & $2$  \\
			 \hline
		\end{tabular}
	\end{center}
	Số trung vị của mẫu số  liệu trên là
	\choice 
		{$50$} 
		 {$60$}
		{\True  $55$}
		{$48$} 
	\loigiai{ 
		Ta có:  Số phần tử của mẫu là $n=40\Rightarrow \dfrac{n}{2}=20$.\\
		Suy  ra nhóm $\left[50; 60\right)$ chứa số trung vị.\\
		Do đó số trung vị của mẫu số  liệu trên là
		 \[M_e=50+\dfrac{20-12}{16} \cdot 10=55.\]
	} 
\end{ex}

% nguồn: Ok_Toan11_TracNghiem_Lan2
\begin{ex}%[1D5H2-2]%[ma:1db219c7]
	Doanh thu bán hàng trong $20$ ngày được lựa chọn ngẫu nhiên của một của  hàng được ghi lại ở bảng sau:
	\begin{center}
		\begin{tabular}{|c|c|c|c|c|c|c|}
			\hline 
			Doanh thu & $[5;7)$ & $[7;9)$ & $[9;11)$ & $[11;13)$ &  $[13;15)$ \\
			 \hline
			Số ngày & $2$ & $7$ & $7$ & $3$ & $1$ \\ 
			\hline 
		\end{tabular}
	\end{center}
	Trung  vị của mẫu số liệu trên thuộc khoảng nào trong các khoảng dưới đây?
	 \choice
		{$\left[7;  \; 9\right)$}
		{\True $\left[9; \; 11\right)$} 
		{$\left[11; \; 13\right)$} 
		{$\left[13; \; 15\right)$} 
	\loigiai{ 
		Goi $x_1$, $x_2$, $\ldots$, $x_{20}$ là  doanh thu bán hàng trong 20 ngày xếp theo thứ tự không giảm.\\
		Khi đó: $x_1$,  $x_2 \in \left[5;7\right)$, $x_3$, $\ldots$, $ x_9 \in \left[7; \; 9\right)$, $x_9$, $\ldots$, $x_{16} \in \left[9; \; 11\right)$ $x_{17}$, $\ldots$, $x_{19} \in \left[11; \; 13\right)$, $x_{20} \in \left[13; \; 15\right)$.\\
		Do đó, trung vị của mẫu số liệu thuộc nhóm $\left[9; \;  11\right)$.
	} 
\end{ex}

% nguồn: Ok_Toan11_TracNghiem_Lan2
\begin{ex}%[1D5H2-2]%[ma:7b60ea8f]
	Khảo sát thời gian tập thể dục của một số học sinh khối 11 thu  được mẫu số liệu ghép nhóm sau:
	\begin{center}
		\begin{tabular}{|c|c|c|c|c|c|c|}
			 \hline
			Thời gian (phút) & $[0;20)$ & $[20;40)$ & $[40;60)$ & $[60;80)$  & $[80;100)$ \\
			 \hline
			 Số học sinh & $5$ & $9$ & $12$ & $10$ & $6$ \\
			 \hline
		\end{tabular}
	\end{center}
	Nhóm chứa trung vị của mẫu số liệu trên là 
	 \choice
		 {\True $[40;60)$}
		 {$[20;40)$}
		{$[60;80)$} 
		 {$[80;100)$}
	\loigiai{ 
		Ta có: $n=42$.\\ 
		Nên trung vị của mẫu  số liệu trên là $Q_2=\dfrac{x_{21}+x_{22}}{2}$.\\
		Mà $x_{21}$,  $x_{22} \in \left[40;60\right)$.\\
		Vậy nhóm chứa trung vị của mẫu số liệu  trên là nhóm $[40;60)$.
	} 
\end{ex}

% nguồn: Ok_Toan11_TracNghiem_Lan2
\begin{ex}%[1D5H2-2]%[ma:9e20c9dc]
	Khảo sát về số giờ mượn sách thư viện của học sinh khối $11$ trường Y ta  được một mẫu số liệu ghép nhóm như sau:
	\begin{center}
		\begin{tabular}{|c|c|c|c|c|c|c|}
			\hline 
			Số  giờ mượn & $[0;4)$ & $[4;8)$ & $[8;12)$ & $[12;16)$ & $[16;20)$ \\
			 \hline
			Số học sinh  & $54$ & $78$ & $120$ & $45$ & $12$ \\
			\hline 
		\end{tabular}
	\end{center}
	Số trung vị của mẫu số liệu trên là 
	\choice 
		 {\True $8{,}75$}
		{$8$} 
		{$10$} 
		 {$9$}
	 \loigiai{
		Ta có: Cỡ mẫu $n=309\Rightarrow \dfrac{n}{2}=154{,}5$.\\ 
		Suy ra nhóm $\left[8; 12\right)$ chứa số  trung vị.\\
		Do đó số trung vị của mẫu số liệu  trên là
		 \[M_e=8+\dfrac{154{,}5-132}{120}\cdot4=8{,}75.\]
	 }
\end{ex}

% nguồn: Ok_Toan11_TracNghiem_Lan2
\begin{ex}%[1D5H2-2]%[ma:196e390d]
	Thống kê điểm thi đánh giá năng lực môn Toán của $120$ học sinh một  trường THPT được cho trong bảng sau:
	\begin{center}
		\begin{tabular}{|c|c|c|c|c|c|}
			 \hline
			 Điểm & $\left[0;20\right)$ & $\left[20;40\right)$ & $\left[40;60\right)$ & $\left[60;80\right)$ & $\left[80;100\right]$ \\
			\hline 
			Số học sinh & $25$ & $35$ & $37$ & $15$ &  $8$ \\
			 \hline
		\end{tabular}
	\end{center}
	Trung vị của mẫu số liệu ghép nhóm thuộc  khoảng nào sau đây?
	 \choice
		 {\True $\left[40;45\right]$}
		 {$\left[45;50\right]$}
		{$\left[50;55\right]$} 
		{$\left[55;60\right]$} 
	 \loigiai{
		\begin{itemize}
			\item Cỡ mẫu  $n=120$.
			\item Gọi $x_1$, $x_2$, $\cdots$, $x_{120}$ lần lượt là số điểm thi đánh giá năng lực  môn Toán được sắp xếp theo thứ tự không giảm.
			\item Do  $x_{60}\in\left[20;40\right)$, $x_{61}\in\left[40;60\right)$ nên $M_e=Q_2=40$.
		\end{itemize}
	} 
\end{ex}

% nguồn: Ok_Toan11_TracNghiem_Lan2
\begin{ex}%[1D5H2-2]%[ma:6eadb64b]
	Cho mẫu số liệu ghép nhóm về tuổi thọ (đơn vị tính là năm) của một loại bóng  đèn mới như sau:
	\begin{center}
		\begin{tabular}{|c|c|c|c|c|}
			\hline 
			 Tuổi thọ & $\left[2;3{,}5\right)$ & $\left[3{,}5;5\right)$ & $\left[5;6{,}5\right)$ & $\left[6{,}5;8\right)$ \\
			 \hline
			 Số bóng đèn & $8$ & $22$ & $35$ & $15$ \\
			 \hline
		\end{tabular}
	\end{center}
	Nhóm chứa  trung vị của mẫu số liệu là
	\choice 
		 {$\left[2;3{,}5\right)$}
		{$\left[3{,}5;5\right)$} 
		{\True $\left[5;6{,}5\right)$} 
		 {$\left[6{,}5;8\right)$}
	 \loigiai{
		\begin{itemize}
			\item Cỡ  mẫu $n=80$.
			\item Gọi $x_1$, $x_2$, $\cdots$, $x_{80}$ lần lượt là tuổi thọ (năm) của các bóng đèn được sắp  xếp theo thứ tự không giảm.
			 \item $M_e=Q_2=\dfrac{1}{2}\left(x_{40}+x_{41}\right)\in\left[5; 6{,}5\right)$.
		\end{itemize}
	 }
\end{ex}

% nguồn: Ok_TracNghiemLop11_Lan1
\begin{ex}%[1D5H2-2]%[ma:92566d3b]
	Trong tuẫn lễ bảo vệ môi trường, các học sinh khối  $11$ tiến hành thu nhặt vỏ chai nhựa để tái chế. Nhà trường thống kê kết quả thu nhặt vỏ chai của học sinh khối $11$ ở bảng sau
	\begin{center}
		\begin{tabular}{|l|c|c|c|c|c|}
			\hline Số vỏ chai nhựa & {$ [11; 15] $} & {$ [16; 20]  $} & {$ [21; 25] $} & {$ [26; 30] $} & {$ [31; 35] $} \\
			\hline Số học sinh & $53$ &  $82$ & $48$ & $39$ & $18$ \\
			\hline 
		\end{tabular}
	\end{center}
	Hãy tìm trung vị của mẫu số liệu ghép nhóm trên. 
	 \choice
		 {$19{,}51$}
		{\True $19{,}59$} 
		{$20{,}2$} 
		 {$18{,}6$}
	\loigiai{ 
		Số học sinh tham  gia thu nhặt vỏ chai nhựa là $n=53+82+48+39+18=240$.\\
		Gọi $x_{1}; x_{2}; \ldots; x_{240} $ lần lượt là số vỏ chai nhựa $240$ học sinh khối $11$ thu nhặt được xếp  theo thứ tự không giảm.\\
		Do $ x_{1}; \ldots; x_{53} \in[10{,}5; 15{,}5); x_{54}; \ldots; x_{135} \in[15{,}5; 20{,}5) $ nên trung vị của mẫu số liệu $ x_{1}; x_{2};  \ldots; x_{240} $ là $ \dfrac{1}{2}\left(x_{120}+x_{121}\right) \in[15{,}5; 20{,}5) $.\\
		Ta xác  định được $ n=240, n_{m}=82, C=53, u_{m}=15{,}5; u_{m+1}=20{,}5 $.\\
		Vậy trung vị của  mẫu số liệu ghép nhóm là \\ $\mathrm{M}_{e}=15{,}5+\dfrac{\dfrac{240}{2}-53}{82} \cdot (20{,}5-15{,}5)=\dfrac{803}{41} \approx 19{,}59$.
	} 
\end{ex}

% nguồn: Ok_TracNghiemLop11_Lan1
\begin{ex}%[1D5H2-2]%[ma:cabffca0]
	Thời gian truy cập Internet mỗi  buối tối của một số học sinh được cho trong bảng sau
	\begin{center}
		\begin{tabular}{|l|c|c|c|c|c|}
			\hline Thời gian (phút) & {$ [9{,}5; 12{,}5)  $} & {$ [12{,}5; 15{,}5) $} & {$ [15{,}5; 18{,}5) $} & {$ [18{,}5; 21{,}5) $} & {$ [21{,}5; 24{,}5) $} \\
			\hline Số học sinh & $3$ & $12$ & $15$ &  $24$ & $2$ \\
			\hline 
		\end{tabular}
	\end{center}
	Tính trung vị của mẫu  số liệu ghép nhóm này.
	 \choice
		{\True $18{,}1$} 
		{$18{,}5$} 
		 {$17{,}2$}
		 {$15{,}6$}
	 \loigiai{
		Cỡ mẫu là $n=3+12+15+24+2=56$.\\ 
		Gọi $x_{1}, \ldots, x_{56} $ là thời gian vào internet của $56$ học sinh và giả sử dãy này được sắp xếp theo thứ tự không giảm. \\ 
		Khi đó, trung vị là $ \dfrac{x_{28}+x_{29}}{2} $.\\ 
		Do $2$ giá trị  $ x_{28}, x_{29} $ thuộc nhóm $ [15{,}5; 18{,}5) $ nên nhóm này chứa trung vị.\\
		Do đó, $p=3; a_{3}=15{,}5; m_{3}=15; m_{1}+m_{2}=3+12=15; a_{4}-a_{3}=3 $ và ta có $\mathrm{M}_{e}=15{,}5+\dfrac{\dfrac{56}{2}-15}{15} \cdot  3=18{,}1$.
	 }
\end{ex}

% nguồn: Ok_TracNghiemLop11_Lan1
\begin{ex}%[1D5H2-2]%[ma:af1f3af4]
	Kết quả khảo sát cân nặng của $ 25 $ quả bơ ở một lô hàng cho trong bảng  sau:
	\begin{center}
		\begin{tabular}{|c|c|c|c|c|c|}
			 \hline
			\textbf{Cân nặng}\textbf{ (g)} & $ \left[150; 155 \right) $ & $ \left[155; 160\right) $  & $ \left[160; 165\right) $ & $ \left[165; 170\right) $ & $ \left[170; 175 \right) $ \\
			 \hline
			\textbf{Số quả bơ} & $ 1 $  & $ 7 $ & $ 12 $ & $ 3 $ & $ 2 $ \\
			 \hline
		\end{tabular}
	\end{center}
	Hãy tìm  trung vị của mẫu số liệu ghép nhóm trên.
	 \choice{\True$161{,}875$}{$161{,}875$}{$161{,}875$}{$161{,}875$}
	 \loigiai{
		Gọi $ x_1; x_2; \ldots; x_{25} $ là cân nặng của $ 25$ quả bơ xếp theo thứ tự không giảm.\\ 
		Do $ x_1\in \left[150; 155 \right) $; $ x_2, \ldots, x_8 \in \left[155; 160\right) $;  $ x_9, \ldots, x_{20} \in \left[160; 165\right) $ nên trung vị của mẫu số liệu $ x_1; x_2; \ldots; x_{25} $ là $ x_{13}\in \left[160; 165\right)$.\\
		Ta xác định được $ n=25 $, $p=3 $, $ a_3=160$, $  m_3=12$, $m_1+m_2=1+7=8$, $a_4-a_3=165-160=5$.\\
		Vậy trung vị của mẫu số liệu ghép nhóm là $$ M_e=160+\dfrac{\dfrac{25}{2}-8}{12}\cdot  5 =161{,}875.$$
	} 
\end{ex}

% nguồn: Ok_TracNghiemLop11_Lan1
\begin{ex}%[1D5H2-2]%[ma:8d25cceb]
	Cho mẫu số liệu ghép nhóm về thống kê  huyết áp (mmHG) của $20$ người, ta có bảng số liệu sau
	\begin{center}
		\begin{tabular}{|l|c|c|c|c|c|c|}
			\hline 
			Tuổi thọ &$[7;80)$ & $[80;90)$& $[90;100)$& $[100;110)$& $[110;120)$ &  $[120;130)$ \\
			\hline 
			Số lượng & $4$& $2$ &  $3$& $6$ & $3$ &$2$ \\ \hline
		\end{tabular}
	\end{center}
	Tìm trung vị của  mẫu số liệu ghép nhóm trên (kết quả làm tròn đến hàng phần trăm).
	\choice 
		{$101{,}76  $}
		{$100{,}67  $}
		{$100{,}76  $}
		 {\True $101{,}67 $}
	 \loigiai{
		\begin{center}
			\begin{tabular}{|l|c|c|c|c|c|c|}
				 \hline
				 Tuổi thọ &$[7;80)$ & $[80;90)$& $[90;100)$& $[100;110)$& $[110;120)$ & $[120;130)$ \\
				\hline 
				Số lượng  & $4$& $2$ & $3$& $6$ & $3$ &$2$ \\ \hline
				Tần số tích lũy & $4$ & $6$& $9$ & $15$ & $18$ & $20  $ \\ \hline
			\end{tabular}
		\end{center}
		Ta có  $\dfrac{n}{2}=10$.\\
		Nhóm chứa  $M_e$ là $[100;110)$.\\
		Suy  ra trung vị $M_e = 100 + \dfrac{10-9}{6} \cdot (110-100)=101{,}67$.
	} 
\end{ex}

% nguồn: Ok_TracNghiemLop11_Lan1
\begin{ex}%[1D5H2-2]%[ma:cdab351b]
	Khảo sát thời gian tập thể dục của một số học sinh  khối 11 thu được mẫu số liệu ghép nhóm sau:
	\begin{center}
		\begin{tabular}{|l|c|c|c|c|c|}
			\hline Thời gian (phút) &  {$[0; 20)$} & {$[20; 40)$} & {$[40; 60)$} & {$[60; 80)$} & {$[80; 100)$} \\
			\hline Số học sinh & $5$ &  $9$ & $12$ & $10$ & $6$ \\
			\hline 
		\end{tabular}
	\end{center}
	Nhóm chứa trung vị của mẫu số liệu trên  là
	 \choice
		{\True $[40;  60)$}
		{$[20; 40)$} 
		{$[60;  80)$}
		 {$[80; 100)$}
	 \loigiai{
		Ta có  $n=42$.\\
		Nên trung vị của mẫu số liệu trên  là $Q_2=\dfrac{x_{21}+x_{22}}{2}$.\\
		Mà $x_{21}, x_{22} \in[40; 60)$ nên  nhóm chứa trung vị của mẫu số liệu trên là nhóm $[40; 60)$.
	 }
\end{ex}

% nguồn: Ok_TracNghiemLop11_Lan1
\begin{ex}%[1D5H2-2]%[ma:34dafbc3]
	Số lượng huy chương vàng tại Sea Games $31$ được  thống kê trong bảng sau:
	\begin{center}
		\begin{tabular}{|c|c|c|c|c|}
			 \hline
			Số huy chương  & $[0;10)$ & $[10;50)$ & $[50;100)$ & $[100;210)$ \\
			 \hline
			Quốc gia & $5$  & $2$ & $3$ & $1$ \\
			 \hline
		\end{tabular}
	\end{center}
	Xác định trung vị của mẫu số liệu  ghép nhóm trên
	\choice 
		 {$10$}
		 {$15$}
		{\True  $20$}
		{$30$} 
	 \loigiai{
		\begin{itemize}
			 \item Cỡ mẫu $n=11$.
			\item Gọi $x_1$, $x_2$, $\cdots$, $x_{11}$ lần lượt là số huy chương nhận được của các quốc  gia được sắp xếp theo thứ tự không giảm.
			\item $M_e=x_6\in\left[10; 20\right)$. Suy  ra
			\[M_e=10+\dfrac{2\cdot\dfrac{11}{4}-5}{2}\cdot\left(50-10\right)=20.\] 
		\end{itemize}
	 }
\end{ex}

% nguồn: Ok_TracNghiemLop11_Lan1
\begin{ex}%[1D5H2-2]%[ma:96f16cfe]
	Khảo sát thời gian chạy  bộ trong một ngày của một số học sinh khối $11$ thu được mẫu số liệu ghép nhóm sau
	\begin{center}
		\begin{tabular}{|c|c|c|c|c|c|}
			\hline 
			Thời gian (phút) & $[0;20)$ & $[20;40)$ & $[40; 60)$ & $[60;80)$  & $[80;100)$ \\
			\hline 
			Số học sinh & $5$ &  $9$ & $12$ & $10$ & $6$ \\
			\hline 
		\end{tabular}
	\end{center}
	Nhóm chứa trung vị  là
	 \choice
		 {$[0;20)$}
		{$[20;40)$} 
		{\True $[40;60)$} 
		 {$[60;80)$}
	\loigiai{ 
		Ta có $\dfrac{n}{2}=\dfrac{42}{2}=21$.\\ 
		Nhóm chứa trung vị  $[40;60)$.
	} 
\end{ex}

% nguồn: Ok_TracNghiemLop11_Lan1
\begin{ex}%[1D5H2-2]%[ma:70770411]
	Cho  bảng số liệu ghép nhóm về nhiệt độ $\left({ }^{\circ} \mathrm{C}\right)$ tại một thành phố
	\begin{center}
		\begin{tabular}{|c|c|c|c|c|}
			\hline Nhiệt độ $\left( ^{\circ}\mathrm{C}\right)$ &  {$[18; 22)$} & {$[22; 26)$} & {$[26; 30)$} & {$[30; 34)$} \\
			\hline Số ngày & $5$  & $15$ & $6$ & $3$ \\
			 \hline
		\end{tabular}
	\end{center}
	Tính trung vị của mẫu số liệu trên (kết quả làm tròn đến hàng  phần mười).
	 \choice
		 {$24{,}3$}
		 {$24{,}6$}
		{\True $24{,}5$} 
		{$24{,}4$} 
	 \loigiai{
		Gọi $x_1$; $x_2$; $\ldots$; $x_{29}$ là  nhiệt độ của $29$ ngày xếp theo thứ tự không giảm.\\
		Do $x_1,\ldots,x_5 \in[18; 22)$; $x_6, \ldots, x_{20}\in[22; 26)$ nên trung vị  của mẫu số liệu $x_1; x_2; \ldots; x_{29}$ là $x_{15} \in[22; 26)$.\\
		 Ta xác định được $n=29$, $n_m=15$, $C=5$, $u_m=22$, $u_{m+1}=26$.\\
		Vậy trung vị của mẫu số liệu  ghép nhóm là
		$$ 
		 M_e=22+\dfrac{\dfrac{29}{2}-5}{15} \cdot(26-22)=\dfrac{368}{15}\approx 24{,}5.
		$$ 
	} 
\end{ex}

% nguồn: Ok_TracNghiemLop11_Lan1
\begin{ex}%[1D5H2-2]%[ma:2c70639f]
	Thời gian (phút) xem tivi mỗi buổi tối của một số học sinh được cho trong bảng  sau:
	\begin{center}
		\begin{tabular}{|l|c|c|c|c|c|c|c|}
			 \hline
			Thời gian & $[6{,}5; 9{,}5)$ & $[9{,}5; 12{,}5)$ & $[12{,}5; 15{,}5)$ & $[15{,}5; 18{,}5)$ & $[18{,}5; 21{,}5)$  & $[21{,}5; 24{,}5)$ & $[24{,}5; 27{,}5)$ \\
			\hline 
			Số HS & $2$ & $3$ & $12$ & $15$ &  $24$ & $2$ & $2$ \\
			 \hline
		\end{tabular}
	\end{center}
	Số trung vị của mẫu số liệu  ghép nhóm này là
	\choice 
		{\True $18{,}1$} 
		{$15{,}1$} 
		{$21{,}1$} 
		{$15$} 
	\loigiai{ 
		 Cỡ mẫu $n=2+3+12+15+24+2+2=60$.\\
		Nhóm chứa trung vị $\left[15{,}5; 18{,}5\right)$. Suy ra: $u_m=15{,}5$ và  $u_{m+1}=18{,}5$.\\
		Tần số của nhóm  chứa trung vị: $n_m=15$; $C=n_1+n_2+n_3=2+3+12=17$.\\
		 Vậy trung vị của mẫu số liệu ghép nhóm là $M_e=15{,}5+\dfrac{\dfrac{60}{2}-17}{15}\cdot \left(18{,}5-15{,}5\right)=18{,}1$.
	} 
\end{ex}

% nguồn: Ok_TracNghiemLop11_Lan1
\begin{ex}%[1D5H2-2]%[ma:9f3bca45]
	Cho mẫu số liệu ghép nhóm về  thống kê điểm số của học sinh tham dự kỳ thi học sinh giỏi
	 toán, ta có bảng số liệu sau
	\begin{center}
		\begin{tabular}{|c|c|c|c|c|c|c|c|}
			\hline 
			Điểm&$[8;10)$&$[10;12)$&$[12;14)$&$[14;16)$&$[16;18)$&$[18;20)$\\ 
			 \hline
			 Số học sinh &$6$&$21$&$30$&$25$&$14$&$4$\\
			 \hline
		\end{tabular}
	\end{center}
	Tìm  nhóm chứa trung vị của mẫu số liệu ghép nhóm trên.
	\choice 
		{$[10;12)$} 
		 {$[14;16)$}
		 {\True$[12;14)$}
		{$[16;18)$} 
	 \loigiai{
		Số phần tử của  mẫu là $n=6+21+30=25+14+4=100$, ta có $\dfrac{n}{2}=\dfrac{100}{2}=50$.\\
		Do $cf_2=27<50<56=cf_3$ nên trung vị của mẫu  nằm ở nhóm $3$ là nhóm $[12;14)$.
	 }
\end{ex}

% nguồn: Ok_TracNghiemLop11_Lan1
\begin{ex}%[1D5H2-2]%[ma:b4fb1055]
	Khảo  sát về số giờ mượn sách thư viện của học sinh khối 11 trường Y ta được một mẫu số liệu ghép nhóm như sau:
	\begin{center}
		\begin{tabular}{|c|c|c|c|c|c|}
			\hline Số giờ mượn & {$[0; 4)$} & {$[4; 8)$}  & {$[8; 12)$} & {$[12; 16)$} & {$[16; 20)$} \\
			\hline Số học sinh & $54$ & $78$ & $120$ &  $45$ & $12$\\
			\hline 
		\end{tabular}
	\end{center}
	Số trung vị của mẫu số liệu  trên là
	\choice 
		{\True $8,75$} 
		{$8$} 
		 {$10$}
		 {$9$}
	\loigiai{ 
		Cỡ mẫu $n=309\Rightarrow\dfrac{n}{2}=154,5$.\\ 
		Suy ra nhóm $\left[8\,;\,12\right)$ chứa  số trung vị.\\
		Do đó số trung vị của mẫu  số liệu trên là\\
	$M_{\mathrm  e}=8+\dfrac{154{,}5-132}{120}\cdot4=8{,}75$.}
\end{ex}

% nguồn: Ok_TracNghiemLop11_Lan1
\begin{ex}%[1D5H2-2]%[ma:1a9b9547]
	Kết quả khảo sát cân nặng của $25$ quả bơ ở  một lô hàng cho trong bảng sau
	\begin{center}
		\begin{tabular}{|l|c|c|c|c|c|}
			\hline Cân nặng $(\mathrm{g})$ & {$[150;  155)$} & {$ [155; 160)$} & {$[160; 165) $} & {$[165; 170)$} & {$[170; 175)$} \\
			\hline Số quả bơ & $1$ & $7$ & $12$ & $3$ & $2$  \\
			\hline 
		\end{tabular}
	\end{center}
	Trung vị của mẫu số liệu trên  thuộc khoảng nào trong các khoảng dưới đây?
	\choice 
		 {$[170; 175)$}
		{$[155; 160)$} 
		{$[165;  170)$}
		{\True $[160;  165)$}
	 \loigiai{
		Gọi $x_{1}; x_{2}; \ldots; x_{25}$ là cân nặng của $25$ quả bơ xếp theo thứ tự  không giảm.\\
		Do $x_{1} \in[150; 155) $; $ x_{2}; \ldots; x_{8} \in[155; 160); x_{9}; \ldots; x_{20} \in[160; 165) $ nên trung vị của mẫu số liệu  $ x_{1}; x_{2}; \ldots; x_{25} $ là $ x_{13} \in[160; 165)$.
	} 
\end{ex}

% nguồn: Ok_TracNghiemLop11_Lan1
\begin{ex}%[1D5H2-2]%[ma:1115d0ba]
	Điều tra $42$ học sinh của một lớp $11$ về số giờ tự học ở  nhà, người ta có bảng sau đây
	\begin{center}
		\begin{tabular}{|c|c|c|}
			\hline 
			Lớp & Tần số & Tần số tích lũy  \\
			\hline 
			 $\left[1; 2\right)$ & $8$ & $8$ \\
			 \hline
			$\left[2; 3\right)$ & $10$ &  $18$ \\
			\hline 
			$\left[3; 4\right)$ & $12$ & $30$  \\
			 \hline
			$\left[4;  5\right)$ & $9$ & $39$ \\
			 \hline
			$\left[5; 6\right)$  & $3$ & $42$ \\
			\hline 
		\end{tabular}
	\end{center}
	Số trung  vị của mẫu số liệu là.
	\choice 
		{$4{,}25$} 
		{$3{,}75$} 
		{$4{,}75$} 
		{\True $3{,}25$} 
	\loigiai{ 
		Ta có số phần  tử của mẫu là\\
		$n=42\Rightarrow \dfrac{n}{2}=21$.\\ 
		Mà $cf_2=18< 21< cf_3=30$ suy  ra nhóm $3$ là nhóm đầu tiên có tần số tích lũy lớn hơn hoặc bằng $21$.\\
		Xét nhóm $3$ là nhóm $\left[3; 4\right)$  có $r=3$; $d=1$; $n_3=12$ và nhóm $2$ là nhóm $\left[2; 3\right)$ có $cf_2=18$.\\
		Áp  dụng công thức ta có trung vị của mẫu số liệu là\\
		$$M_e=3+\left(\dfrac{21-18}{12} \right)\cdot  1=3{,}25.$$
	 }
\end{ex}

% nguồn: Ok_TracNghiemLop11_Lan1
\begin{ex}%[1D5H2-2]%[ma:dfadd31f]
	Thống kê điểm thi đánh giá năng lực môn Toán của $120$ học sinh một trường THPT  được cho trong bảng sau:
	\begin{center}
		\begin{tabular}{|c|c|c|c|c|c|}
			 \hline
			Điểm  & $[0;20)$ & $[20;40)$ & $[40;60)$ & $[60;80)$ & $[80;100]$ \\
			\hline 
			Số học sinh & $25$ & $35$ &  $37$ & $15$ & $8$ \\
			\hline 
		\end{tabular}
	\end{center}
	Trung vị của mẫu số liệu ghép nhóm thuộc khoảng nào sau  đây?
	 \choice
		{\True  $[40;45]$}
		{$[45;50]$} 
		 {$[50;55]$}
		 {$[55;60]$}
	 \loigiai{
		\begin{itemize}
			\item Cỡ  mẫu $n=120$.
			\item Gọi $x_1$, $x_2$, $\cdots$, $x_{120}$  lần lượt là số điểm thi đánh giá năng lực môn Toán được sắp xếp theo thứ tự không giảm.
			\item Do $x_{60}\in\left[20;40\right)$, $x_{61}\in [40;60)$ nên  $M_e=Q_2=40$.
		\end{itemize}
	} 
\end{ex}

% nguồn: Ok_TracNghiemLop11_Lan1
\begin{ex}%[1D5V2-2]%[ma:acbf942b]
	Cho mẫu số liệu ghép nhóm về thống kê điểm số (thang điểm $20$) của 100 học sinh  tham dự kỳ thi học sinh giỏi toán, ta có bảng số liệu sau:
	\begin{center}
		\begin{tabular}{|l|c|c|c|c|c|c|}
			\hline Điểm & {$[8; 10)$} & {$[10; 12)$} & {$[12; 14)$} & {$[14; 16)$} &  {$[16; 18)$} & {$[18; 20)$} \\
			\hline Số học sinh & $6$ & $21$ & $30$ & $25$  & $14$ & $4$\\
			 \hline
		\end{tabular}
	\end{center}
	Tìm trung vị của mẫu số liệu ghép nhóm trên. (Kết  quả làm tròn đến hàng phần trăm).
	 \choice
		{$12{,}18$} 
		 {$12{,}81$}
		 {$13{,}35$}
		{\True  $13{,}53$}
	\loigiai{ 
		 Cỡ mẫu là $n=6+21+30+25+14+4=100$.\\
		Gọi $x_1; x_2; \ldots; x_{100}$ là số điểm của  100 học sinh và được sắp xếp theo thứ tự tăng dần. Khi đó, trung vị là $\dfrac{x_{50}+x_{51}}{2}$, do hai giá trị $x_{50}, x_{51}$ thuộc nhóm $[12; 14)$ nên nhóm này chứa trung vị.\\
		Ta có  $p=3$; $a_3=12$; $m_3=30$; $m_1+m_2=25$; $a_4-a_3=2$ khi đó trung vị của mẫu số liệu là
		$$M_e=12+\dfrac{\dfrac{100}{2}-27}{30}\cdot 2=13{,}53.$$ 
	 }
\end{ex}

% nguồn: Ok_TracNghiemLop11_Lan1
\begin{ex}%[1D5V2-2]%[ma:7f3bc6a2]
	Cho mẫu số liệu ghép nhóm về chiều  cao của $25$ cây dừa giống như sau:
	\begin{center}
		\begin{tabular}{|l|c|c|c|c|c|}
			\hline Chiều cao $(\mathrm{cm})$ & {$[0; 10)$}  & {$[10; 20)$} & {$[20; 30)$} & {$[30; 40)$} & {$[40; 50)$} \\
			 \hline Số cây & $4$ & $6$ & $7$ & $5$ & $3$ \\
			\hline 
		\end{tabular}
	\end{center}
	Trung vị  của mẫu số liệu ghép nhóm này là
	\choice 
		{$M_e=\dfrac{175}{7}$} 
		{$M_e=\dfrac{165}{5}$} 
		{$M_e=\dfrac{165}{7}$} 
		{\True $M_e=\dfrac{165}{3}$} 
	 \loigiai{
		Cỡ mẫu $n=4+6+7+5+3=25$.\\ 
		$x_1, x_2, \ldots, x_{25}$ là chiều cao của $25$ cây  dừa giống được sắp xếp theo thứ tự không giảm. Khi đó, trung vị là $x_{13}$. Do $x_{13}$ thuộc nhóm $[20; 30)$ nên nhóm này chứa trung vị.\\
		Do đó $p=3$,  $a_3=20$, $m_3=7$, $m_1+m_2=10$, $a_4-a_3=10$.\\
		Vậy $M_e=20+\dfrac{\dfrac{25}{2}-10}{7} \cdot  10=\dfrac{165}{7}$.
	 }
\end{ex}

% nguồn: Ok_TracNghiemLop11_Lan1
\begin{ex}%[1D5B2-2]%[ma:88a63ec3]
	Cho hàm số $y =\dfrac{x^2 + x}{x - 2}$ có đồ thị $(C)$. Phương trình tiếp tuyến tại điểm $A(1; -2)$ của $(C)$  là
	\choice 
		{$y = -3x + 5$} 
		{$y  = -5x + 7$}
		{\True $y  = -5x + 3$}
		{$y = -4x  + 6$}
	 \loigiai{
		Ta có $y'=\dfrac{(2x+1)(x-2)-(x^2+x)}{(x-2)^2}=\dfrac{x^2-4x-2}{(x-2)^2}$.\\ 
		Phương trình tiếp tuyến của  $(C)$ tại $A(1;-2)$ là $$y=y'(1)\cdot (x-1)-2\Leftrightarrow y=-5(x-1)-2\Leftrightarrow y=-5x+3.$$
	} 
\end{ex}

% nguồn: Ok_TracNghiemLop11_Lan1
\begin{ex}%[1D5B2-2]%[ma:bdbbe8e5]
	Cho hàm số $y=x^3-3x^2+3$ có đồ thị là $(C)$. Viết phương  trình tiếp tuyến của $(C)$ tại điểm có hoành độ $x=1$.
	\choice 
		 {$y=2x-1$}
		{$y=-x+2$} 
		 {$y=-3x+3$}
		{\True  $y=-3x+4$}
	\loigiai{Ta có $x_0 =  1$ nên $y_0 = 1$ và $y'(1) = -3$.\\
		Phương trình tiếp tuyến cần tìm là $y = -3(x-1) + 1 = -3x +  4$.
	 }
\end{ex}

% nguồn: Ok_TracNghiemLop11_Lan1
\begin{ex}%[1D5Y2-2]%[ma:04b465e3]
	Cho đường cong $y=f(x)$ có đồ thị là $(C)$. Tiếp tuyến của đường cong $(C)$ tại điểm $M(x_0;  y_0)$ có hệ số góc là $k$, khi đó $k$ bằng
	\choice 
		{$k=f\left(x_0\right)$} 
		{\True  $k=f'\left(x_0\right)$}
		 {$k=f(x)$}
		{$k=f'(x)$} 
	\loigiai{ 
		Tiếp tuyến  của đường cong $(C)$ tại điểm $M(x_0; y_0)$ có hệ số góc là $k=f'\left(x_0\right)$.
	} 
\end{ex}

% nguồn: Ok_TracNghiemLop11_Lan1
\begin{ex}%[1D5B2-2]%[ma:0c04d4e6]
	Cho hàm số $y=\dfrac {1}{3} x^3-2x^2+x+2$ có đồ thị  $(C)$. Phương trình tiếp tuyến của đồ thị $(C)$ biết tiếp tuyến song song với đường thẳng
	 $d\colon y=-2x+\dfrac {10}{3}$ là
	\choice 
		{\True $y=-2x+2$} 
		 {$y=-2x-2$}
		{$y=-2x+10, y=-2x-\dfrac  {2}{3}$}
		{$y=-2x+10,  y=-2x+\dfrac {2}{3}$}
	\loigiai{ 
		Ta có $y'=x^2-4x+1$.  Vì tiếp tuyến song song với đường thẳng $d\colon y=-2x+\dfrac {10}{3}$ nên \\
		$y'(x_0)=-2\Leftrightarrow x_0^2-4x_0+1=-2  \Leftrightarrow x_0=1$ hoặc $x_0=3$.\\
		Phương trình tiếp tuyến là $y=-2x+2 $ hoặc $y=-2x+\dfrac  {10}{3}$ (loại).
	} 
\end{ex}

% nguồn: Ok_TracNghiemLop11_Lan1
\begin{ex}%[1D5B2-2]%[ma:22243092]
	Số  điểm một cầu thủ bóng rổ ghi được trong $20$ trận đấu được cho ở bảng sau
	\begin{center}
		\begin{tabular}{|l|c|c|c|c|c|}
			\hline Điểm số &{$[5{,}5;10{,}5)$}&{$[10{,}5;15{,}5)$}&{$[15{,}5;20{,}5)$}&{$[20{,}5;25{,}5)$}\\ 
			 \hline Số trận & $3$ & $9$ & $2$ & $6$ \\
			\hline 
		\end{tabular}
	\end{center}
	Tìm trung vị của mẫu số  liệu đã cho.
	\choice 
		 {\True $M_e\approx 11{,}6$}
		 {$M_e\approx 14{,}4$}
		{$M_e\approx 15{,}6$} 
		{$M_e\approx  21{,}3$}
	\loigiai{ 
		Cỡ mẫu là $n=20$.\\ 
		Gọi  $x_1, \ldots, x_{20}$ là mẫu số liệu được sắp xếp theo thứ tự không giảm.\\
		Khi đó,  trung vị là $\dfrac{x_{10}+x_{11}}{2}\in [10{,}5;15{,}5)$ \\Do đó, $p=2; \;a_2=10{,}5;\; m_2=9; \;m_1=3, \;a_3-a_2=5$ và ta có
		 $$
		M_e=10{,}5+\dfrac{\dfrac{20}{2}-3}{9}\cdot 5\approx 14{,}4. 
		$$ 
	 }
\end{ex}

% nguồn: Ok_TracNghiemLop11_Lan1
\begin{ex}%[1D5Y2-2]%[ma:5d7cb405]
	Hệ số góc $k$ của tiếp tuyến  đồ thị hàm số $y = x^3 + 1$ tại điểm $M(1; 2)$ là
	 \choice
		 {$k = 12$}
		 {\True $k = 3$}
		 {$k = 5$}
		{$k  = 4$}
	\loigiai{Ta có $y' =  3x^2$. Tại điểm $M$ có hoành độ $x = 1$ thì $k = y'(1) = 3\cdot 1^2 = 3$.
	 }
\end{ex}

% nguồn: Ok_TracNghiemLop11_Lan1
\begin{ex}%[1D5G2-2]%[ma:c3f24e1a]
	Cho các hàm số $y = f(x)$; $y = f \left( f(x) \right)$; $y = f(x^2 + 4)$ có đồ thị lần lượt là $(C_1)$; $(C_2)$; $(C_3)$. Đường thẳng $x=1$ cắt $(C_1)$, $(C_2)$, $(C_3)$ lần lượt tại $M$, $N$, $P$. Biết phương trình tiếp tuyến của $(C_1)$ tại $M$ và của $(C_2)$ tại $N$ lần lượt là $y  = 3x + 2$ và $y = 12x - 5$. Phương trình tiếp tuyến của $(C_3)$ tại $P$ có dạng là $y = ax + b$. Tính $a + b$.
	\choice 
		{\True $7$} 
		{$9$} 
		{$8$} 
		 {$6$}
	 \loigiai{
		Chú ý các công thức đạo hàm $(f(f(x))'=f'(f(x))\cdot f'(x)$; $(f(x^2+4))'=2x\cdot f'(x^2+4)$.\\ 
		Tiếp tuyến của $(C_1)$ tại $M$ có phương trình là $y = f'(1)(x - 1) +  f(1)$. \\
		Tiếp tuyến của $(C_2)$ tại $N$ có phương trình là $y = f'(1) \cdot f' \left( f(1) \right) (x - 1) +  f \left( f(1)\right)$. \\
		Tiếp tuyến của $(C_3)$ tại $P$ có phương trình là $y = 2f'(5)(x  - 1)+f(5)$. \\
		Khi đó ta có hệ $\heva{& f'(1) = 3 & \\ & -f'(1) + f(1) = 2 & \\ & f'(1) \cdot f' \left( f(1) \right) = 12 & \\ & - f'(1) \cdot f' \left( f(1) \right) + f \left( f(1) \right) = - 5} \Leftrightarrow \heva{& f'(1) = 3 & \\ & f(1) = 5 & \\ & f'(5) = 4 & \\ & - 3 \cdot f'(5) + f(5) = - 5} \Leftrightarrow \heva{& f'(1) = 3 & \\ & f(1) = 5 & \\ & f'(5)  = 4 & \\ & f(5) = 7.} $ \\
	Khi đó phương trình tiếp tuyến của $(C_3)$ tại $P$  là $y = 2\cdot 4 (x - 1) + 7 = 8x-1$. Vậy $a + b = 7$. }
\end{ex}

% nguồn: Ok_TracNghiemLop11_Lan1
\begin{ex}%[1D5Y2-2]%[ma:46ebce40]
	Nếu hàm số $y=f\left( {x} \right)$ có đạo hàm tại  $x_0$ thì phương trình tiếp tuyến của đồ thị hàm số tại điểm $M_0\left( {x_0;f\left( {x_0} \right)} \right)$ là
	 \choice
		 {$y=f'\left( {x} \right)\left( {x-x_0} \right)+f\left( {x_0} \right)$}
		{$y=f'\left( {x}  \right)\left( {x-x_0} \right)-f\left( {x_0} \right)$}
		{\True $y=f'\left(  {x_0} \right)\left( {x-x_0} \right)+f\left( {x_0} \right)$}
		{$y=f'\left( {x_0} \right)\left(  {x-x_0} \right)-f\left( {x_0} \right)$}
	\loigiai{Theo ý nghĩa hình học của đạo hàm, phương  trình tiếp tuyến của đồ thị hàm số tại điểm $M\left( {x_0;f\left( {x_0} \right)} \right)$ là $y=f'\left( {x_0} \right)\left( {x-x_0} \right)+f\left( {x_0} \right)$.}
\end{ex}

% nguồn: Ok_TracNghiemLop11_Lan1
\begin{ex}%[1D5B2-2]%[ma:4e6025e1]
	Phỏng vấn một số học sinh khối  11 vể thời gian (giờ) ngủ của một buổi tối, thu được bảng số liệu như sau.
	\begin{center}
		\begin{tabular}{|l|c|c|c|c|c|}
			\hline Thời gian  (giờ) &{$[4; 5)$}&{$[5; 6)$}&{$[6; 7)$}&{$[7; 8)$}&{$[8; 9)$}\\
			\hline Số học sinh & $6$  & $10$ & $13$ & $9$ & $7$ \\
			\hline 
		\end{tabular}
	\end{center}
	Hãy cho biết $75 \%$ học sinh  khối 11 ngủ ít nhất bao nhiêu giờ?
	 \choice
		 {$7{,}675$}
		 {\True $7{,}53$}
		{$8$} 
		{  $7{,}9$}
	 \loigiai{
		Cỡ mẫu là  $n=45$.\\
		Gọi $x_1, \ldots, x_{45}$ là mẫu số liệu được sắp xếp theo thứ tự không giảm. Khi đó, trung vị là $x_{23}$. Do đó, tứ phân vị thứ ba $Q_3$ là $x_{34}$. Do $x_{34}$ đều thuộc nhóm $[7;8)$ nên nhóm này chứa $Q_3$. Do đó, $p=4; \;a_4=7;\; m_4=9; \;m_1+m_2+m_3=29;  \;a_5-a_4=1$ và ta có
		$$ 
		Q_3=7+\dfrac{\dfrac{3  \cdot 45}{4}-29}{9}\cdot 1\approx7{,}53.
		 $$
		Vậy $75\%$ học sinh khối  11 ngủ ít nhất $7{,}53$ giờ.
	 }
\end{ex}

% nguồn: Ok_TracNghiemLop11_Lan1
\begin{ex}%[1D5K2-2]%[ma:89b34689]
	Ký hiệu $d$ là tiếp tuyến của đồ thị hàm số $y = x^4 - 4x^2 + 2m^2 + 1\quad (C)$ tại giao điểm của $(C)$ với trục hoành đồng thời $(C)$ đi  qua điểm $A(1; 0)$. Hỏi có bao nhiêu đường thẳng $d$ thỏa mãn bài toán?
	 \choice
		 {$3$}
		{$2$} 
		{$8$} 
		 {\True $4$}
	\loigiai{Vì $(C)$ đi qua điểm $A(1; 0)$ nên $1-4+2m^2+1=0\Leftrightarrow  m^2=1$. Do đó, $y=x^4-4x^2+3$.\\
	Phương trình hoành độ giao điểm giữa $(C)$ và trục hoành là $$x^4-4x^2+3=0\Leftrightarrow\hoac{& x^2=1 \\ & x^2=3}\Leftrightarrow\hoac{& x=\pm 1  \\ & x=\pm \sqrt{3}.}$$
	Như thế  đồ thị $(C)$ cắt trục hoành tại bốn điểm $M(x_0; 0)$ với $x_0\in\{\pm 1; \pm \sqrt{3}\}$.\\
	Phương trình đường tiếp tuyến của đồ thị $(C)$ tại $M(x_0;0)$ là $y=y'\left(x_0\right)(x-x_0)+0$.\\ 
	Ta có $y'\left(x_0\right)\in \{-4; 4; 4\sqrt{3};  -4\sqrt{3}\}$.\\
	Vậy có  tất cả $4$ đường thẳng $d$ thỏa mãn bài toán.
	 }
\end{ex}

% nguồn: Ok_TracNghiemLop11_Lan1
\begin{ex}%[1D5Y2-2]%[ma:f82e548f]
	Tiếp tuyến của đường cong $(C)\colon y=x\sqrt{x+1}$ tại điểm $M(3;6)$ có hệ  số góc bằng
	 \choice
		 {\True $\dfrac{11}{4}$}
		 {$\dfrac{1}{4}$}
		{$-\dfrac{11}{4}$} 
		 {$-\dfrac{1}{4}$}
	 \loigiai{
		Hệ số góc của tiếp tuyến tại  điểm $M(3;6)$ là $y'(3)=\dfrac{11}{4}$.
	 }
\end{ex}


%%%%%%%%%%%%%%%% ĐÚNG/SAI %%%%%%%%%%%%%%%%

% nguồn: Ok_dungsai_lop11_Lan1
\begin{ex}%[1D5H2-2]%[ma:92b382d2]
 Khi đo  mắt cho học sinh khối $10$ ở một trường THPT nhân viên y tế ghi nhận lại ở bảng sau
 \begin{tabular}{|c|c|c|c|c|c|}
  \hline 
  Thời gian  & $[0{,}25;0{,}75)$ & $[0{,}75;1{,}25)$ & $[1{,}25;1{,}75)$ & $[1{,}75;2{,}25)$ & $[2{,}25;2{,}75)$ \\
   \hline
  Số lần & $25$ & $32$ & $14$ & $12$  & $4$ \\
  \hline 
 \end{tabular}
 \choiceTF 
  {\True Số trung  bình của mẫu số liệu trên là $1{,}14$}
  {\True Nhóm chứa mốt của  số liệu là $[0{,}75;1{,}25)$}
  {\True Mốt của mẫu số liệu là $M_0=0{,}89$} 
  {\True Trung vị  của mẫu số liệu là $M_e=1{,}039$}
  \loigiai{
  \begin{tabular}{|l|l|l|l|l|l|}
    \hline
   Thời gian & {$[0{,}25; 0{,}75)$} & {$[0{,}75;  1{,}25)$} & {$[1{,}25; 1{,}75)$} & {$[1{,}75; 2{,}25)$} & {$[2{,}25; 2{,}75)$} \\
   \hline 
   Giá trị  đại diện & $0{,}50$ & $1{,}00$ & $1{,}50$ & $2{,}00$ & $2{,}50$ \\
   \hline 
   Số lần & $25$ &  $32$ & $14$ & $12$ & $4$ \\
    \hline
  \end{tabular}
  \begin{itemchoice}
   \itemch Đúng.  Số trung bình của mẫu số liệu trên là
   \[\dfrac{0{,}50\cdot 25+1{,}00\cdot 32+1{,}50\cdot 14+2{,}00\cdot 12+2{,}50\cdot  4}{87}=1{,}14.
   \] 
   \itemch Đúng. Nhóm chứa mốt của số liệu  là $[0{,}75;1{,}25)$.
   \itemch Đúng. Mốt của mẫu số liệu là $$M_0=0{,}75+\dfrac{32-25}{(32-25)+(32-14)} (1{,}25-0{,}75)=0{,}89.$$ 
   \itemch Đúng. Gọi $x_1, x_2, \ldots x_{87}$ lần lượt là chỉ số mắt cận của các học sinh sắp xếp theo thứ tự không giảm.\\ 
   Ta có $x_1, \ldots, x_{25} \in [0{,}25;0{,}75);x_{26}, \ldots, x_{57}  \in [0{,}75;1{,}25)$; nên trung vị của mẫu là $x_{44} \in [0{,}75;1{,}25)$.\\
   Ta xác định được $n=87$, $n_m=32$, $C=25$, $u_m=0{,}75$;  $u_{m+1}=1{,}25$.\\
    Nên $M_e=0{,}75+\dfrac{\dfrac{87}{2}-25}{32} (1{,}25-0{,}75)=1{,}039$.
  \end{itemchoice}
 } 
\end{ex}

% nguồn: Ok_dungsai_lop11_Lan2
\begin{ex}%[1D5H2-2]%[ma:0df09d5a]
	Khi điều tra về chiều cao của $40$ học sinh lớp 11A bộ phận y tế của nhà trường thu được mẫu số liệu  ghép nhóm sau
	\begin{center}
		\begin{tabular}{|c|c|c|c|c|c|}
			 \hline
			Chiều cao (cm)  & $[140;146)$ & $[146;152)$ & $[152;158)$ & $[158;164)$ & $[164;170]$ \\
			\hline 
			 Số học sinh & $3$ & $12$ & $11$ & $8$ & $6$ \\
			\hline 
		\end{tabular}
	\end{center}
	 \choiceTF
		{\True Số học sinh có chiều cao trên  $152$ cm là $25$}
		 {\True Tất cả học sinh của lớp đều có chiều cao không dưới $140$ cm}
		{\True  Nhóm chứa trung vị là nhóm thứ $3$}
		{\True Chiều  cao trung bình của $40$ học sinh là $155{,}3$ cm}
	\loigiai{ 
		 Bảng tần số ghép nhóm theo giá trị đại diện là
		\begin{center}
			\begin{tabular}{|c|c|c|c|c|c|}
				\hline 
				Chiều cao (cm) & $[140;146)$ & $[146;152)$ & $[152;158)$ & $[158;164)$ & $[164;170]$  \\
				 \hline
				Giá  trị đại diện & $143$ & $149$ & $155$ & $161$ & $167$ \\
				 \hline
				Số học sinh & $3$ & $12$ & $11$ & $8$ &  $6$ \\
				 \hline
			\end{tabular}
		\end{center}
		\begin{itemchoice}
			\itemch \textbf{Đúng.}\\ 
			Số học sinh có chiều cao  trên $152$ cm là $25$.
			\itemch \textbf{Đúng.}\\ 
			Tất cả học sinh của lớp đều có chiều cao không dưới $140$  cm.
			 \itemch \textbf{Đúng.}\\
			Ta có $n=40$, $\dfrac{n}{2}=20$. Nhóm  thứ $3$ là nhóm đầu tiên có tần số tích lũy lớn hơn $20$.\\
			Do đó nhóm chứa trung vị là nhóm thứ $3$. 
			\itemch \textbf{Đúng.}\\ 
			 Chiều cao trung bình của $40$ học sinh là
			\[\overline{x}=\dfrac{143\cdot 3+149\cdot 12+155\cdot 11+161\cdot 8+167\cdot 6}{40}=155{,}3.\] 
		\end{itemchoice}
	 }
\end{ex}

% nguồn: Ok_dungsai_lop11_Lan2
\begin{ex}%[1D5V2-2]%[ma:a944d2d2]
	Kết quả đo chiều cao của $25$  học sinh nam lớp 10A của một trường THPT được cho trong bảng sau
	\begin{center}
		\begin{tabular}{|c|c|c|c|c|c|}
			 \hline
			 Chiều cao (cm) & $[150; 155)$ & $[155; 160)$ & $[160; 165)$ & $[165; 170)$ & $[170; 175)$ \\ \hline
			Số  học sinh & $1$ & $7$ & $12$ & $3$ & $2$ \\ \hline
		\end{tabular}
	\end{center}
	 \choiceTF
		{\True Giá  trị đại diện của nhóm $[160;165)$ là $162{,}5$}
		{Độ dài của nhóm $[155;160)$  là $4$}
		{\True Chiều cao trung bình các bạn học sinh  nam trong lớp $\overline{x}=162{,}1$}
		{Trung vị của mẫu số liệu ghép nhóm trên là  $M_e=162$}
	\loigiai{ 
		\begin{center}
			Bảng tần số ghép  nhóm theo giá trị đại diện là
			\begin{tabular}{|c|c|c|c|c|c|}
				 \hline
				Nhóm & $[150;155)$  & $[155;160)$ & $[160;165)$ & $[165;170)$ & $[170;175)$ \\ \hline
				Giá trị đại diện & $152{,}5$ & $157{,}5$ & $162{,}5$ & $167{,}5$  & $172{,}5$ \\ \hline
				Tần số &  $1$ & $7$ & $12$ & $3$ & $2$ \\ \hline
			\end{tabular}
		\end{center}
		\begin{itemchoice}
			\itemch  \textbf{Đúng.}\\
			 Giá trị đại diện của nhóm $[160;165)$ bằng $\dfrac{1}{2} (160+165)=162{,}5$
			\itemch \textbf{Sai.}\\ 
			Độ dài của nhóm $[155;160)$ bằng $160-155=5$. 
			 \itemch \textbf{Đúng.}\\
			Chiều  cao trung bình học sinh nam trong lớp
			\[\overline{x}=\dfrac{152{,}5.1+157{,}5.7+162{,}5.12+167{,}5.3+172{,}5.2}{25}=162{,}1.\] 
			\itemch  \textbf{Sai.}\\
			Gọi $x_{1};x_{2}$; $\ldots$;$x_{25}$ là chiều cao của $25$ học sinh nam sắp xếp theo thứ tự không giảm. Trung vị của dãy số liệu là $x_{13} \in  [160;165)$.\\
			Xác định được $n=25$; $n_{m}=12$; $C=8$; $u_{m}=160$; $u_{m+1}=165$. Trung vị của mẫu  số liệu ghép nhóm trên là
			\[M_{e}=u_{m}+\dfrac{\dfrac{n}{2}-C}{n_{m}  }\cdot(u_{m+1}-u_{m})=160+\dfrac{\dfrac{25}{2}-8}{12}\cdot(165-160)=161{,}875.\]
		\end{itemchoice}
	} 
\end{ex}

% nguồn: Ok_dungsai_lop11_Lan2
\begin{ex}%[1D5V2-2]%[ma:2312b2e3]
	Thời gian truy cập Internet mỗi buổi tối  của một số học sinh được cho trong bảng sau
	\begin{center}
		\begin{tabular}{|c|c|c|c|c|c|}
			\hline 
			Thời gian (Phút)  & $[9{,}5; 12{,}5)$ & $[12{,}5; 15{,}5)$ & $[15{,}5; 18{,}5)$ & $[18{,}5; 21{,}5)$ & $[21{,}5; 24{,}5)$ \\
			\hline 
			 Số học sinh & $3$ & $12$ & $15$ & $24$ & $2$ \\
			\hline 
		\end{tabular}
	\end{center}
	 \choiceTF
		 {\True Nhóm chứa mốt là $[18{,}5;21{,}5)$}
		{\True Thời gian học sinh truy cập nhiều nhất của các học sinh đó  là $16$ phút}
		{Thời gian truy cập trung bình của các học sinh đó  là $16$ phút}
		{Trung vị của mẫu số liệu trên  là $20$}
	 \loigiai{
		Bảng tần số ghép nhóm theo  giá trị đại diện là
		\begin{center}
			\begin{tabular}{|c|c|c|c|c|c|}
				 \hline
				Thời gian (Phút) & $[9{,}5; 12{,}5)$ & $[12{,}5; 15{,}5)$  & $[15{,}5; 18{,}5)$ & $[18{,}5; 21{,}5)$ & $[21{,}5; 24{,}5)$ \\
				\hline 
				Giá trị đại  diện & $11$ & $14$ & $17$ & $20$ & $23$ \\
				 \hline
				Số học sinh & $3$ &  $12$ & $15$ & $24$ & $2$ \\
				\hline 
			\end{tabular}
		\end{center}
		\begin{itemchoice}
			\itemch  \textbf{Đúng.}\\
			Tần số  lớn nhất là $24$ nên nhóm chứa mốt là $[18,5;21,5)$.
			\itemch  \textbf{Đúng.}\\
			Ta có $j=4$,  $a_{4}=18{,}5$, $m_{4}=24$, $m_{3}=15$, $m_{5}=2$, $h=3$. Do đó
			 \[M_{o}=18{,}5+\dfrac{24-15}{(24-7)+(14-2)}\cdot3\approx 16.\]
			\itemch \textbf{Sai.}\\ 
			Thời gian truy cập trung bình của các  học sinh đó là
			\[\dfrac{11\cdot 3+14\cdot 12+17\cdot 15+20\cdot 24+23\cdot 2}{56}  \approx 9.\]
			\itemch  \textbf{Sai.}\\
			Ta  có $n=56$. Trung vị là $\dfrac{x_{28} +x_{29} }{2}$. Do $x_{28}$ và $x_{29}$ thuộc $[15{,}5;18{,}5)$ nên nhóm này chứa trung vị. Do đó $p=3$; $a_{3}=15{,}5$; $m_{3}=15$; $m_{1} +m_{2}=3+12=15$; $a_{4}-a_{3}=3$. Khi đó
			\[M_e= 15{,}5  + \dfrac{28-15}{15{,}5}\cdot3=\dfrac{1\,117}{62}.\]
		\end{itemchoice}
	} 
\end{ex}

% nguồn: Ok_dungsai_lop11_Lan2
\begin{ex}%[1D5V2-2]%[ma:aa9fdc8c]
	Thống kê điểm trung bình môn Toán của một số học sinh lớp $11$ được  cho ở bảng sau
	\begin{center}
		\begin{tabular}{|c|c|c|c|c|c|c|c|}
			\hline 
			Khoảng điểm & $[6{,}5; 7)$ & $[7; 7{,}5)$ &  $[7{,}5; 8)$ & $[8; 8{,}5)$ & $[8{,}5; 9)$ & $[9; 9{,}5)$ & $[9{,}5; 10)$ \\ \hline
			Tần số  & $8$ & $1$0 & $1$6 & $2$4 & $1$3 & $7$ & $4$ \\ \hline
		\end{tabular}
	\end{center}
	 \choiceTF
		{\True Điểm trung bình  môn Toán của một số học sinh lớp $11$ nằm trong khoảng $(8{,}1;8{,}3)$}
		{Trung vị  của mẫu số liệu ghép nhóm có giá trị lớn hơn $8{,}2$}
		{\True Nhóm  chứa mốt của mẫu số liệu trên là nhóm $[8;8{,}5)$}
		{Mốt  của mẫu số liệu trên thuộc khoảng $(8{,}3;8{,}5)$}
	\loigiai{ 
		Bảng  tần số ghép nhóm theo giá trị đại diện là
		\begin{center}
			\begin{tabular}{|c|c|c|c|c|c|c|c|}
				\hline 
				Khoảng điểm & $[6{,}5; 7)$ & $[7; 7{,}5)$ & $[7{,}5; 8)$ &$ [8; 8{,}5)$ & $[8{,}5;  9)$ & $[9; 9{,}5)$ & $[9{,}5; 10)$ \\ \hline
				Giá trị đại diện & $6{,}75$  & $7{,}25$ & $7{,}75$ & $8{,}25$ & $8{,}75$ & $9{,}25$ & $9{,}75$ \\ \hline
				Tần số & $8$ & $10$ & $16$ & $24$ & $13$ & $7$ & $4$ \\  \hline
			\end{tabular}
		\end{center}
		\begin{itemchoice}
			\itemch \textbf{Đúng.}\\ 
			Điểm trung bình  môn Toán của một số học sinh lớp $11$ là
			 \[\overline{x}=\dfrac{8\cdot6{,}75+10\cdot7{,}25+16\cdot7{,}75+24\cdot8{,}25+13\cdot8{,}75+7\cdot9{,}25+4\cdot9{,}75}{82} \approx 8{,}12.\]
			 \itemch \textbf{Sai.}\\
			Trung vị của mẫu số liệu ghép  nhóm là
			\[M_{e}=8+\dfrac{\dfrac{82}{2}-(8+10+16)}{24}\cdot(8{,}5-8)\approx  8{,}15.\]
			\itemch  \textbf{Đúng.}\\
			Vì  tần số lớn nhất là $24$ thuộc nhóm $[8;8,5)$ nên nhóm chứa mốt của mẫu số liệu trên là nhóm $[8;8,5)$.
			\itemch \textbf{Sai.}\\ 
			Mốt  $M_{o}$ chứa trong nhóm $[8;8{,}5)$. Do đó $u_{m}=8$; $u_{m+1}=8{,}5$, $u_{m+1}-u_{m}=8{,}5-8=0{,}5$; $n_{m-1}=16$; $n_{m}=24$; $n_{m+1}=13$.
			\[M_{o}=8+\dfrac{24-16}{(24-16)+(24-13)}\cdot(8{,}5-8)=8{,}21.\] 
		\end{itemchoice}
	} 
\end{ex}


%%%%%%%%%%%%%%%% TRẢ LỜI NGẮN %%%%%%%%%%%%%%%%

% nguồn: Ok_TraLoiNgan_Lop11_Lan1
\begin{ex}%[1D5N2-2]%[ma:3a008cbd]
	Người ta ghi chép lại trọng lượng (gam) một loại cá rô được nuôi trong ao theo một chế  độ đặc biệt sau $6$ tháng, họ có bảng tần số ghép nhóm sau
	\begin{center}
		\begin{tabular}{|c|c|c|c|c|c|c|}
			 \hline
			Trọng lượng & $[60;70)$ & $[70;80)$ & $[80;90)$ & $[90;100)$ & $[100;110)$ &  $[110;120)$ \\
			\hline 
			Số cá & $13$ & $24$ & $55$ & $61$ & $31$ &  $16$ \\
			\hline 
		\end{tabular}
	\end{center}
	Tìm trung vị của mẫu số liệu ghép nhóm  đã cho (Kết quả làm tròn đến hàng phần chục).
	 \shortans{91,3}
	 \loigiai{
		Cỡ mẫu của mẫu số liệu  là $n=13+24+55+61+31+16=200$.\\
		Gọi $x_1,  x_2, \ldots, x_{200}$ là mẫu số liệu được sắp xếp theo thứ tự không giảm.\\
		Trung vị của mẫu số liệu này là $\dfrac{x_{100}+x_{101}}{2} \in [90;100)$.\\ 
		Ta có  $n_m=61$; $C=13+24+55=92$; $u_m=90$; $u_{m+1}=100$.\\
		Trung vị của mẫu số liệu ghép nhóm  là
		\[M_e=90+\dfrac{\dfrac{200}{2}-92}{61} (100-90)=\dfrac{5570}{61} \approx 91{,}3.\] 
	} 
\end{ex}

% nguồn: Ok_TraLoiNgan_Lop11_Lan1
\begin{ex}%[1D5N2-2]%[ma:53c52f40]
	Kết  quả khảo sát hàm lượng vitamin $C$ của một số loại trái cây cho ở bảng sau
	\begin{center}
		\begin{tabular}{|c|c|c|c|c|c|}
			\hline 
			Hàm lượng vitamin C$(\%)$ & $[6;7)$ &$[7;8)$& $[8;9)$& $[9;10)$& $[10;11)$\\ 
			 \hline
			Số lượng & $3$ & $4$ & $5$ & $2$  & $1$ \\
			 \hline
		\end{tabular}
	\end{center}
	Hãy  tìm trung vị của mẫu số liệu ghép nhóm trên.
	\shortans{8,1} 
	\loigiai{ 
		Gọi  $x_1;x_2;\ldots;x_{15}$ là hàm lượng vitamin $C$ của một số loại trái cây xếp theo thứ tự không giảm.\\
		Do $x_1;x_2;x_3 \in [6;7);x_4;\ldots;x_7 \in [7;8);x_8;\ldots;x_{12} \in [8;9)$  nên trung vị của mẫu là $x_8 \in [8;9)$.\\
		Ta xác định được  $n=15$, $n_m=5$, $C=3+4=7$, $u_m=8$, $u_{m+1}=9$.
		\[M_e=8+\dfrac{\dfrac{15}{2}-7}{5}  (9-8)=8{,}1.\]
	} 
\end{ex}

% nguồn: Ok_TraLoiNgan_Lop11_Lan1
\begin{ex}%[1D5N2-2]%[ma:7f94bf6a]
	Điều tra về số lượng học sinh khối $11$ trong một lớp học, người ta thu được dữ liệu của $100$ lớp học và  có bảng phân phối tần số ghép nhóm sau
	\begin{center}
		\begin{tabular}{|c|c|c|c|c|c|}
			 \hline
			 Nhóm & $[36;38)$ & $[38;40)$ & $[40;42)$ & $[42;44)$ & $[44;46)$ \\
			 \hline
			Tần số & $9$ & $15$ & $25$ & $30$  & $21$ \\
			 \hline
		\end{tabular}
	\end{center}
	Tìm trung vị của mẫu số liệu ghép  nhóm trên (Kết quả làm tròn đến hàng đơn vị).
	 \shortans{42}
	\loigiai{ 
		Cỡ mẫu của mẫu số liệu  là $n=100$.\\
		Gọi $x_1, x_2, x_3, \ldots, x_{100}$ là số học sinh trong một lớp học khối 11 được điều tra được sắp  xếp theo thứ tự không giảm.\\
		Trung vị  của mẫu số liệu là $\dfrac{x_{50}+x_{51}}{2} \in [42;44)$.\\
		Ta có $n_m=30$; $C=9+15+25=49$; $u_m=42$; $u_{m+1}=44$.\\ 
		Trung  vị của mẫu số liệu ghép nhóm là
		\[M_e=42+\dfrac{\dfrac{100}{2}-49}{30} (44-42)=\dfrac{631}{15} \approx  42.\]
	} 
\end{ex}

% nguồn: Ok_TraLoiNgan_Lop11_Lan1
\begin{ex}%[1D5N2-2]%[ma:2eb19b29]
	Thời gian (phút) truy cập internet mỗi buổi tối của một số  học sinh được cho trong bảng sau
	\begin{center}
		\begin{tabular}{|c|c|c|c|c|c|}
			\hline 
			Nhóm & $[9{,}5;12{,}5)$ & $[12{,}5;15{,}5)$  & $[15{,}5;18{,}5)$ & $[18{,}5;21{,}5)$ & $[21{,}5;24{,}5)$ \\
			\hline 
			Số học $\sinh$ & $3$ & $12$ & $15$ &  $24$ & $2$ \\
			\hline 
		\end{tabular}
	\end{center}
	Tìm trung vị của mẫu số liệu ghép nhóm  trên.
	\shortans{18,1} 
	 \loigiai{
		Cỡ mẫu của mẫu số  liệu là $n=3+12+15+24+2=56$.\\
		Gọi $x_1{,}  x_2{,} x_3{,} \ldots{,} x_{56}$ là thời gian truy cập Internet của lần lượt 56 học sinh theo thứ tự không giảm.\\
		Trung vị của mẫu số liệu là $\dfrac{x_{28}+x_{29}}{2} \in  [15{,}5;18{,}5)$.\\
		Ta có $n_m=15$; $C=3+12=15$;  $u_m=15{,}5$; $u_{m+1}=18{,}5$.\\
		 Khi đó trung vị của mẫu số liệu phép nhóm là
		\[M_e=15{,}5+\dfrac{\dfrac{56}{2}-15}{15} (18{,}5-15{,}5)=\dfrac{181}{10}=18{,}1.\] 
	 }
\end{ex}

% nguồn: Ok_TraLoiNgan_Lop11_Lan1
\begin{ex}%[1D5N2-2]%[ma:05334f24]
	Kết quả thu thập điểm số môn Toán của $25$ học sinh khi  tham gia kì thi học sinh giỏi toán $11$ (thang điểm 20) cho ta bảng tần số ghép nhóm sau
	\begin{center}
		\begin{tabular}{|c|c|c|c|c|c|}
			 \hline
			Nhóm & $[0;4)$ & $[4;8)$ & $[8;12)$ & $[12;16)$ &  $[16;20)$ \\
			\hline 
			Số  học sinh & $1$ & $7$ & $12$ & $3$ & $2$ \\
			\hline 
		\end{tabular}
	\end{center}
	Tìm trung vị của mẫu số liệu ghép nhóm trên. 
	 \shortans{9,5}
	\loigiai{ 
		Cỡ mẫu của mẫu số liệu là $n=25$.\\ 
		Gọi $x_1, x_2, x_3, \ldots, x_{25}$ là điểm số của $25$ học sinh trong kì thi đó và dãy này được sắp  xếp theo thứ tự không giảm.\\
		Trung vị của mẫu  số liệu là $x_{13} \in [8;12)$.\\
		Ta có $n=25$,  $n_m=12$, $C=1+7=8$, $u_m=8$, $u_{m+1}=12$.\\
		Trung vị mẫu  số liệu ghép nhóm là
		\[M_e=u_m+\dfrac{\dfrac{n}{2}-C}{n_m} \left(u_{m+1}-u_m \right)=8+\dfrac{\dfrac{25}{2}-8}{12}  (12-8)=9{,}5.\]
	} 
\end{ex}

% nguồn: Ok_TraLoiNgan_Lop11_Lan1
\begin{ex}%[1D5N2-2]%[ma:6aaafe41]
	Nghiên cứu thời gian chạy một vòng sân trường $300$ m  của $41$ học sinh lớp 11A trường THPT được giáo viên bộ môn Thể dục ghi lại, có kết quả sau
	\begin{center}
		\begin{tabular}{|c|c|c|c|c|c|}
			 \hline
			Thời gian &$[40;45)$&$[45;50)$&$[50;55)$&$[55;60)$&$[60;65)$\\ 
			 \hline
			 Số học sinh & $5$ & $8$ & $13$ & $9$ & $6$ \\
			 \hline
		\end{tabular}
	\end{center}
	Hãy tìm trung vị của mẫu số liệu  ghép nhóm trên (Kết quả làm tròn hàng phần chục).
	\shortans{52,9} 
	 \loigiai{
		Gọi $x_1; x_2; \ldots; x_{41}$ là thời gian chạy của học sinh lớp 11A được xếp theo  thứ tự không giảm.\\
		Do $x_1;x_2;\ldots;x_5 \in [40;45);x_6;\ldots;x_{13} \in [45;50);x_{14};\ldots;x_{26} \in [50;55)$ nên trung vị của mẫu là $x_{21} \in  [50;55)$.\\
		Ta xác định được $n=41$,  $n_m=13$, $C=5+8=13$, $u_m=50$, $u_{m+1}=55$.
		\[M_e=50+\dfrac{\dfrac{41}{2}-13}{13}  (55-50)=52{,}9.\]
	} 
\end{ex}

% nguồn: Ok_TraLoiNgan_Lop11_Lan1
\begin{ex}%[1D5N2-2]%[ma:a013b68b]
	Số tiền mà  học sinh lớp $11$ chi cho ăn uống sinh hoạt trong một tuần được tổng hợp ở bảng sau
	\begin{center}
		\begin{tabular}{|c|c|c|c|c|c|}
			 \hline
			Số tiền (nghìn đồng) &$[300;340)$&$[340;380)$&$[380;420)$&$[420;460)$&$[460;500)$\\ 
			\hline 
			Số học sinh & $5$ & $9$ & $12$ & $7$  & $6$ \\
			\hline 
		\end{tabular}
	\end{center}
	Hãy tìm trung vị của mẫu số liệu ghép nhóm  trên (Kết quả làm tròn đến hàng đơn vị).
	 \shortans{398}
	\loigiai{ 
		Gọi $x_1;x_2;\ldots;x_3$ là số tiền chi  ăn uống sinh hoạt trong tuần của học sinh xếp theo thứ tự không giảm.\\
		Do $x_1;\ldots;x_5 \in [300;340);x_6;\ldots;x_9 \in [340;380);x_{15};\ldots;x_{26}  \in [380;420)$ nên trung vị của mẫu là$x_{19} \in [380;420)$.\\
		Ta  xác định được $n=39$, $n_m=12$, $C=5+9=14$, $u_m=380$, $u_{m+1}=420$.\\
		\[M_e=380+\dfrac{\dfrac{39}{2}-14}{12}  (420-380)\approx 398.\]
	 }
\end{ex}

% nguồn: Ok_TraLoiNgan_Lop11_Lan1
\begin{ex}%[1D5N2-2]%[ma:f65301e9]
	Chiều cao (đơn vị: m) của  $35$ cây bạch đàn được cho ở bảng sau
	\begin{center}
		\begin{tabular}{|c|c|c|c|c|c|}
			\hline 
			Số đo chiều  cao $(m)$&$[6{,}5;7)$&$[7;7{,}5)$&$[7{,}5;8)$&$[8;8{,}5)$&$[8{,}5;9)$\\
			 \hline
			Số cây & $6$  & $9$ & $15$ & $4$ & $1$ \\
			\hline 
		\end{tabular}
	\end{center}
	Hãy tìm  trung vị của mẫu số liệu ghép nhóm trên (Kết quả làm tròn hàng phần trăm).
	 \shortans{7,58}
	 \loigiai{
		Gọi $x_1;x_2;\ldots;x_{35}$ là chiều cao của  các cây bạch đàn xếp theo thứ tự không giảm.\\
		Do  $x_1;\ldots;x_6 \in [6{,}5;7);x_7;\ldots;x_{15} \in [7;7{,}5);x_{16};\ldots;x_{10} \in [7{,}5;8)$ nên trung vị của mẫu là $x_{16} \in [7{,}5;8)$.\\
		Ta xác định được  $n=35$, $n_m=15$, $C=6+9=15$, $u_m=7{,}5$; $u_{m+1}=8$.\\
		\[M_e=7{,}5+\dfrac{\dfrac{35}{2}-15}{15} (8-7{,}5)=7{,}58.\] 
	} 
\end{ex}

% nguồn: Ok_TraLoiNgan_Lop11_Lan1
\begin{ex}%[1D5N2-2]%[ma:eb46e45d]
	Một  mẫu số liệu có bảng tần số ghép nhóm như sau
	\begin{center}
		\begin{tabular}{|c|c|c|c|c|c|}
			 \hline
			Nhóm & $[1;5)$ & $[5;9)$ & $[9;13)$ & $[13;17)$ & $[17;21)$  \\
			\hline 
			Tần số & $4$ &  $8$ & $13$ & $6$ & $4$ \\
			\hline 
		\end{tabular}
	\end{center}
	Tìm trung vị của mẫu số liệu ghép nhóm trên  (Kết quả làm tròn đến hàng phần mười).
	 \shortans{10,7}
	\loigiai{ 
		Cỡ mẫu của mẫu số  liệu là $n=4+8+13+6+5=35$.\\
		Gọi $x_1, x_2, \ldots,  x_{35}$ là mẫu số liệu được sắp xếp theo thứ tự không giảm.\\
		Trung vị của mẫu số liệu  này là $x_{18} \in [9;13)$.\\
		Ta có $n_m=13$; $C=4+8=12$; $u_m=9$;  $u_{m+1}=13$.
		Trung  vị của mẫu số liệu ghép nhóm là $$M_e=9+\dfrac{\dfrac{35}{2}-12}{13} (13-9)=\dfrac{139}{13} \approx 10{,}7.$$
	 }
\end{ex}

% nguồn: Ok_TraLoiNgan_Lop11_Lan1
\begin{ex}%[1D5N2-2]%[ma:af6e4f3f]
	Cho bảng thống kê  cân nặng của $20$ quả cam như sau
	\begin{center}
		\begin{tabular}{|l|c|c|c|c|}
			 \hline
			 Trọng lượng (g)&$[115;130)$&$[130;145)$&$[145;160)$&$[160;175)$\\
			\hline 
			Số quả&$3$&$5$&$8$&$4$\\ 
			 \hline
		\end{tabular}
	\end{center}
	Tìm số trung vị của mẫu số liệu đã cho (làm tròn đến hàng đơn  vị).
	 \shortans{149}
	 \loigiai{
		\begin{itemize}
			\item Nhóm  chứa trung vị là nhóm thứ $3$: $[145;160)$.
			\item Trung vị là $M_e=145+\dfrac{10-8}{8}\cdot 15=\dfrac{595}{4} \approx 149$  (g).
	\end{itemize}}
\end{ex}

% nguồn: Ok_TraLoiNgan_Lop11_Lan1
\begin{ex}%[1D5N2-2]%[ma:40bb350b]
	Số bài  tập của các bộ môn được giáo viên cho học sinh về làm ở khối $11$ ở một nhóm học sinh trường THPT được giao về làm trong 01 tuần được cho như sau
	\begin{center}
		\begin{tabular}{|c|c|c|c|c|c|}
			\hline 
			Số bài  tập &$[5;12)$&$[12;19)$&$[19;26)$&$[26;33)$&$[33;40)$\\
			 \hline
			Số học sinh & $27$ & $58$ & $22$ & $23$ & $10$ \\ 
			\hline 
		\end{tabular}
	\end{center}
	Hãy tìm trung vị của mẫu số liệu ghép nhóm trên (Kết quả làm  tròn hàng phần chục).
	\shortans{23,7} 
	\loigiai{ 
		Gọi $x_1;x_2;\ldots;x_{140}$ là số lượng bài tâp xếp theo  thứ tự không giảm.\\
		Do $x_1;\ldots;x_{27} \in [5;12);x_{28};\ldots;x_{85} \in [12;19)$ nên trung vị của  mẫu là$\dfrac{1}{2} \left(x_{70}+x_{71} \right)\in [12;19)$.\\
		Ta xác  định được $n=140$, $n_m=58$, $C=27$, $u_m=12$, $u_{m+1}=19$.\\
		\[M_e=12+\dfrac{\dfrac{140}{2}-27}{58} (19-12)=23{,}7.\] 
	 }
\end{ex}

% nguồn: Ok_TraLoiNgan_Lop11_Lan1
\begin{ex}%[1D5N2-2]%[ma:f6a3bf51]
	Trong đợt kiểm tra học kỳ $II$ môn Thể dục ở  một trường THPT được giáo viên tổng hợp thời gian chạy của $41$ học sinh ở cự ly$1\,500$ m dưới bảng như sau
	\begin{center}
		\begin{tabular}{|c|c|c|c|c|c|}
			 \hline
			 Thời gian (đơn vị phút) &$[7;9)$&$[9;11)$&$[11;13)$&$[13;15)$&$[15;17)$\\
			 \hline
			Số học sinh & $5$ & $8$ & $13$  & $9$ & $6$ \\
			\hline 
		\end{tabular}
	\end{center}
	Hãy tìm trung vị của mẫu số liệu ghép nhóm trên (Kết quả làm tròn đến hàng phần  mười)
	 \shortans{12,2}
	 \loigiai{
		Gọi $x_1;x_2;\ldots;x_{41}$là  thời gian chạy của học sinh lớp được xếp theo thứ tự không giảm.\\
		Do $x_1;x_2;\ldots;x_5 \in [7;9);x_6 \ldots x_{13} \in [9;11);x_{14} \ldots x_{26} \in [11;13)$ nên trung vị của mẫu là  $x_{21} \in [11;13)$.\\
		 Ta xác định được $n=41$, $n_m=13$, $C=5+8=13$, $u_m=11$, $u_{m+1}=13$.\\
		\[M_e=11+\dfrac{\dfrac{41}{2}-13}{13} (13-11)\approx12{,}2.\] 
	} 
\end{ex}

% nguồn: Ok_TraLoiNgan_Lop11_Lan1
\begin{ex}%[1D5N2-2]%[ma:3d4a1d94]
	Một học viện bóng đá điều tra về lứa tuổi của $100$ học  viên trẻ đăng kí đầu tiên để tham gia khóa học mới và thu được bảng sau
	\begin{center}
		\begin{tabular}{|c|c|c|c|c|c|}
			\hline 
			Nhóm tuổi & $[8;9]$ & $[10;11]$ & $[12;13]$ &  $[14;15]$ & $[16;17]$ \\
			 \hline
			Số học viên & $14$ & $20$ &  $33$ & $18$ & $15$ \\
			 \hline
		\end{tabular}
	\end{center}
	Tìm trung vị của mẫu số liệu ghép nhóm trên (Kết quả làm tròn một  chữ số sau dấu phẩy).
	 \shortans{12,5}
	\loigiai{ 
		Vì độ tuổi được điều tra là số nguyên nên ta  hiệu chỉnh bảng số liệu trên về bảng tần số ghép nhóm sau
		\begin{center}
			\begin{tabular}{|c|c|c|c|c|c|}
				\hline 
				Nhóm tuổi & $[7{,}5;9{,}5)$ &  $[9{,}5;11{,}5)$ & $[11{,}5;13{,}5)$ & $[13{,}5;15{,}5)$ & $[15{,}5;17{,}5)$ \\
				 \hline
				Số học viên & $14$ &  $20$ & $33$ & $18$ & $15$ \\
				\hline 
			\end{tabular}
		\end{center}
		Cỡ mẫu của mẫu số liệu là $n=100$.\\ 
		Gọi $x_1, x_2, \ldots, x_{100}$ là mẫu  số liệu được sắp xếp theo thứ tự không giảm.\\
		Trung vị của mẫu số  liệu này là $\dfrac{x_{50}+x_{51}}{2} \in [11{,}5;13{,}5)$.\\
		Ta  có $n_m=33$; $C=14+20=34$; $u_m=11{,}5$; $u_{m+1}=13{,}5$.
		Trung vị của mẫu số liệu ghép  nhóm là
		\[M_e=11,5+\dfrac{\dfrac{100}{2}-34}{33} (13{,}5-11{,}5)=\dfrac{823}{66}  \approx 12{,}5.\]
	} 
\end{ex}

% nguồn: Ok_TraLoiNgan_Lop11_Lan1
\begin{ex}%[1D5H2-2]%[ma:2f51881b]
	Số câu hỏi mà một nhóm học sinh  làm được trong $1$ bài kiểm tra được cho bởi bảng sau:
	\begin{center}
		\begin{tabular}{|c|c|c|c|c|c|}
			 \hline
			Số bài tập & $[5;12)$ & $[12;19)$ & $[19;26)$ & $[26;33)$ & $[33;40]$ \\ 
			 \hline
			Số học sinh & $27$ & $58$ & $22$ &  $23$ & $10$ \\
			\hline 
		\end{tabular}
	\end{center}
	Trung vị của mẫu số liệu ghép nhóm trên bằng bao nhiêu? (\textit{Kết quả được  làm tròn đến hàng phần mười}).
	\shortans{17,2} 
	 \loigiai{
		Gọi $x_1;x_2;\ldots;x_{140}$ là số câu hỏi mà học sinh làm được  xếp theo thứ tự không giảm.\\
		Do $x_1,\ldots,x_{27}\in [5;12);  x_{28},\ldots,x_{85}\in [12;19)$.\\
		Nên trung vị của mẫu số  liệu $x_1;x_2;\ldots;x_{140}$ là $\dfrac{1}{2}(x_{70}+x_{71})\in[12;19)$.\\
		Ta xác định được $n=140$;  $n_m=58$; $C=27$; $u_m=12$; $u_{m+1}=19$.\\
		Vậy trung vị  của mẫu số liệu ghép nhóm là
		\[ 
		M_{\text{e}}=u_m+\dfrac{\dfrac{n}{2}-C}{n_m}\cdot(u_{m+1}-u_{m})=12+\dfrac{\dfrac{140}{2}-27}{58}\cdot(19-12)\approx17{,}2. 
		\] 
	} 
\end{ex}

% nguồn: Ok_TraLoiNgan_Lop11_Lan1
\begin{ex}%[1D5H2-2]%[ma:1529ff13]
	Thời gian (phút) di chuyển đến trường của  nhóm học sinh trường THPT A được tổng hợp dưới bảng sau:
	\begin{center}
		\begin{tabular}{|c|c|c|c|c|c|c|c|}
			\hline 
			Thời gian  & $[15;20)$ & $[20;25)$ & $[25;30)$ & $[30;35)$ & $[35;40)$ & $[40;44)$ & $[45;50)$ \\
			 \hline
			Số học sinh & $6$ & $14$ & $25$  & $37$ & $13$ & $9$ & $21$ \\
			\hline 
		\end{tabular}
	\end{center}
	Trung vị của mẫu số liệu ghép nhóm trên bằng  bao nhiêu? (\textit{Kết quả được làm tròn đến hàng đơn vị}).
	\shortans{32} 
	\loigiai{ 
		Gọi $x_1;x_2;\ldots;x_{125}$ là thời gian di chuyển của học sinh lớp được xếp  theo thứ tự không giảm.\\
		Do $x_1,\ldots,x_6\in [15;20);  x_7,\ldots,x_{20}\in [20;25); x_{21},\ldots,x_{45}\in[25;30); x_{46},\ldots,x_{82}\in[30;35)$.\\
		Nên trung  vị của mẫu số liệu $x_1;x_2;\ldots;x_{125}$ là $x_{63}\in[30;35)$.\\
		Ta xác định  được $n=125$; $n_m=37$; $C=6+14+25=45$; $u_m=30$; $u_{m+1}=35$.\\
		Vậy trung  vị của mẫu số liệu ghép nhóm là
		\[ 
		 M_{\text{e}}=u_m+\dfrac{\dfrac{n}{2}-C}{n_m}\cdot(u_{m+1}-u_{m})=30+\dfrac{\dfrac{125}{2}-45}{37}\cdot(35-30)\approx32.
		\] 
	} 
\end{ex}

% nguồn: Ok_TraLoiNgan_Lop11_Lan1
\begin{ex}%[1D5H2-2]%[ma:7900e005]
	Trong một bài kiểm tra nội dung chạy ở cự ly $1\,500$\,m, kết quả của các em học sinh  được giáo viên ghi lại như bảng sau:
	\begin{center}
		\begin{tabular}{|c|c|c|c|c|c|}
			 \hline
			Thời gian (phút) & $[7;9)$ & $[9;11)$ & $[11;13)$  & $[13;15)$ & $[15;17)$ \\
			 \hline
			Số học sinh & $5$ & $8$ & $13$ & $9$  & $6$ \\
			\hline 
		\end{tabular}
	\end{center}
	Trung vị của mẫu số liệu ghép nhóm trên bằng bao nhiêu? (\textit{Kết quả được làm tròn  đến hàng phần mười}).
	 \shortans{12,2}
	 \loigiai{
		Gọi $x_1;x_2;\ldots;x_{41}$ là  thời gian chạy của học sinh lớp được xếp theo thứ tự không giảm.\\
		Do $x_1,\ldots,x_5\in [7;9); x_6,\ldots,x_{13}\in [9;11); x_{14},\ldots,x_{26}\in[11;13)$.\\ 
		Nên trung vị của mẫu số liệu $x_1;x_2;\ldots;x_{41}$ là  $x_{21}\in[11;13)$.\\
		Ta xác định được $n=41$; $n_m=13$;  $C=5+8=13$; $u_m=11$; $u_{m+1}=13$.\\
		Vậy trung vị của mẫu  số liệu ghép nhóm là
		\[ 
		M_{\text{e}}=u_m+\dfrac{\dfrac{n}{2}-C}{n_m}\cdot(u_{m+1}-u_{m})=11+\dfrac{\dfrac{41}{2}-13}{13}\cdot(13-11)\approx12{,}2. 
		\] 
	} 
\end{ex}

% nguồn: Ok_TraLoiNgan_Lop11_Lan1
\begin{ex}%[1D5H2-2]%[ma:c2f4b316]
	Chiều cao (đơn vị: m) của $35$ cây bạch đàn  được cho ở bảng sau:
	\begin{center}
		\begin{tabular}{|c|c|c|c|c|c|}
			\hline 
			Số đo chiều cao ($m$) & $[6{,}5;7)$ & $[7;7{,}5)$ & $[7{,}5;8)$ &  $[8;8{,}5)$ & $[8{,}5;9)$ \\
			 \hline
			Số cây & $6$ & $9$  & $15$ & $4$ & $1$ \\
			 \hline
		\end{tabular}
	\end{center}
	Trung vị của mẫu số liệu ghép nhóm trên bằng bao nhiêu? (\textit{Kết quả được làm tròn đến hàng phần  mười}).
	 \shortans{7,6}
	 \loigiai{
		Gọi  $x_1;x_2;\ldots;x_{35}$ là chiều cao của các cây bạch đàn xếp theo thứ tự không giảm.\\
		Do $x_1,\ldots,x_6\in [6{,}5;7); x_7,\ldots,x_{15}\in  [7;7{,}5); x_{16},\ldots,x_{30}\in[7{,}5;8)$.\\
		Nên trung vị của mẫu số liệu  $x_1;x_2;\ldots;x_{35}$ là $x_{18}\in[7{,}5;8)$.\\
		Ta xác định được $n=35$; $n_m=15$; $C=6+9=15$; $u_m=7{,}5$;  $u_{m+1}=8$.\\
		Vậy trung vị của mẫu số  liệu ghép nhóm là
		\[ 
		 M_{\text{e}}=u_m+\dfrac{\dfrac{n}{2}-C}{n_m}\cdot(u_{m+1}-u_{m})=7{,}5+\dfrac{\dfrac{35}{2}-15}{15}\cdot(8-7{,}5)\approx7{,}6.
		\] 
	 }
\end{ex}

% nguồn: Ok_TraLoiNgan_Lop11_Lan1
\begin{ex}%[1D5H2-2]%[ma:f982bfd6]
	Khi đo mắt cho học sinh khối $10$ ở một trường THPT nhân viên y tế ghi nhận  lại ở bảng sau:
	\begin{center}
		\begin{tabular}{|c|c|c|c|c|c|}
			\hline Thời gian  & $[0{,}25;0{,}75)$ & $[0{,}75;1{,}25)$ & $[1{,}25;1{,}75)$ & $[1{,}75;2{,}25)$ & $[2{,}25;2{,}75)$ \\
			\hline  Số lần & $25$ & $32$ & $14$ & $12$ & $4$ \\
			 \hline
		\end{tabular}
	\end{center}
	Mốt của mẫu số liệu là  bao nhiêu?
	\shortans{0,89} 
	 \loigiai{
		Ta có bảng 
		\begin{center}
			\begin{tabular}{|c|c|c|c|c|c|}
				\hline Thời gian & $[0{,}25;0{,}75)$ & $[0{,}75;1{,}25)$  & $[1{,}25;1{,}75)$ & $[1{,}75;2{,}25)$ & $[2{,}25;2{,}75)$ \\
				 \hline
				Giá  trị đại diện & $0{,}50$ & $1{,}00$ & $1{,}50$ & $2{,}00$ & $2{,}50$ \\
				\hline Số lần  & $25$ & $32$ & $14$ & $12$ & $4$ \\
				 \hline
			\end{tabular}
		\end{center}
		Mốt của mẫu số liệu  là
		 $$
		 M_{0}=0{,}75+\dfrac{32-25}{(32-25)+(32-14)}\cdot(1{,}25-0{,}75)=0{,}89.
		$$ 
	} 
\end{ex}

% nguồn: Ok_TraLoiNgan_Lop11_Lan1
\begin{ex}%[1D5H2-2]%[ma:9de590ae]
	Doanh thu của  50 cửa hàng của một công ty trong một tháng (đơn vị: triệu đồng) được thống kê trong bảng tần số ghép nhóm sau:
	\begin{center}
		\begin{tabular}{|c|c|c|c|c|c|}
			\hline Doanh thu & {$[26; 48)$} & {$[48; 70)$} & {$[70; 92)$}  & {$[92; 114)$} & {$[114; 136)$} \\
			\hline 
			Số cửa hàng & $6$ &  $14$ & $17$ & $8$ & $5$ \\
			 \hline
		\end{tabular}
	\end{center}
	Tính số trung vị của bảng  số liệu ghép nhóm trên. (Kết quả lấy một chữ số sau dẩu phẩy).
	 \shortans{76,5}
	 \loigiai{
		\begin{enumerate}[a)]
			\item \begin{center}
			\begin{tabular}{|c|c|c|c|c|c|}
				\hline Doanh thu & {$[26;  48)$} & {$[48; 70)$} & {$[70; 92)$} & {$[92; 114)$} & {$[114; 136)$} \\
				\hline 
				Số cửa hàng & $6$ & $14$ & $17$ & $8$ & $5$ \\ 
				 \hline
				GTĐD & $37$ & $59$ & $81$ & $103$ &  $125$ \\
				 \hline
			\end{tabular}
		\end{center}
		Doanh thu trung bình là $\overline{x}=\dfrac{6\cdot 37+14\cdot 59+17\cdot 81+8\cdot 103+5\cdot 125}{50}\approx 77{,}5$ (triệu đồng). 
		 \item
		\begin{itemize}
			\item Nhóm chứa trung vị là nhóm thứ $3$: $[70;92)$.\\ Trung vị của mẫu  số liệu ghép nhóm $\dfrac{x_{25}+x_{26}}{2}\in [70;92)$.
			 \item Do đó, $M_e=a_3+\dfrac{\frac{n}{2}-(m_1+m_2)}{m_3}\cdot (a_4-a_3)=70+\dfrac{\frac{50}{2}-(6+14)}{17}\cdot (92-70)\approx 76{,}5$.
		\end{itemize}
	\end{enumerate}
} 
\end{ex}

% nguồn: Ok_TraLoiNgan_Lop11_Lan1
\begin{ex}%[1D5H2-2]%[ma:b4506ccf]
	Người ta tiến hành phỏng vấn $40$ người về một mẫu áo khoác. Người điều tra yêu cầu cho điểm 
	mẫu áo đó theo thang điểm là $100$. Kết quả được trình bày trong  bảng ghép nhóm sau
	\begin{center}
		\begin{tabular}{|c|c|c|c|c|c|}
			 \hline
			Nhóm&$[50;60)$&$[60;70)$&$[70;80)$&$[80;90)$&$[90;100)$\\ 
			 \hline
			Tần số  &$4$&$5$&$23$&$6$&$2$\\
			\hline 
		\end{tabular}
	\end{center}
	Hãy tìm trung vị của mẫu số liệu ghép nhóm trên (kết  quả làm tròn đến hàng đơn vị).
	 \shortans{75}
	 \loigiai{
		Ta có $\dfrac{n}{2}=\dfrac{40}{2}=20$ và $4+5<20<4+5+23$ nên nhóm  $[70;80)$ là nhóm đầu tiên có tần số tích luỹ lớn hơn hoặc bằng $20$.\\
		Nhóm $[70;80)$ có $r=70$, $d=10$,  $n_3=23$, nhóm $[60;70)$ có $cf_2=9$.\\
		Vậy trung vị của mẫu ghép  nhóm là
		$$M_e=70+\dfrac{20-9}{23}\cdot 10=\dfrac{1720}{23}\approx 75.$$ 
	} 
\end{ex}

% nguồn: Ok_TraLoiNgan_Lop11_Lan1
\begin{ex}%[1D5H2-2]%[ma:97de7f16]
	Thời gian (phút) truy cập Internet mỗi buổi tối của một số  học sinh được cho trong bảng sau
	\begin{center}
		\begin{tabular}{|c|*{8}{w{c}{1.7cm}|}}
			 \hline
			Thời gian & $[9{,}5; 13{,}5)$ & $[13{,}5; 17{,}5)$ & $[17{,}5; 21{,}5)$ & $[21{,}5; 25{,}5)$ & $[25{,}5; 29{,}5)$ & $[29{,}5; 33{,}5)$ & $[33{,}5;  37{,}5)$\\
			 \hline
			Số học  sinh & $21$ & $24$ & $2$ & $21$ & $5$ & $23$ & $4$ \\
			\hline 
		\end{tabular}
	\end{center}
	Tứ phân  vị thứ nhất của mẫu số liệu ghép nhóm này (làm tròn đến hàng phần chục).
	 \shortans{14,2}
	 \loigiai{
		Bước 1: Tổng số học sinh là $100$.\\ 
		Bước 2: Xác định nhóm chứa tứ  phân vị thứ nhất $Q_1$. Tứ phân vị thứ nhất nằm trong nhóm thứ 2.
		 Tính $Q_1$ theo công thức:
		$$Q_1 = L_1 + \frac{\frac{n}{4}  - CF_1}{F_1} \cdot c_1$$
		 Trong đó
		$L_1 =  13{,}5$ là cận dưới của nhóm chứa $Q_1$,
		$CF_1 = 21$ là tần số tích  lũy trước nhóm chứa $Q_1$,
		 $F_1 = 24$ là tần số của nhóm chứa $Q_1$,
		$c_1  = 4$ là độ rộng của nhóm.
		Vậy 
		$$Q_1 = 13{,}5 + \frac{\frac{100}{4} - 21}{24}  \cdot 4{,}0 = 14{,}2.$$
	} 
\end{ex}

% nguồn: Ok_TraLoiNgan_Lop11_Lan1
\begin{ex}%[1D5H2-2]%[ma:f716bc9b]
	Số tiền mà một nhóm học sinh lớp  $11$ chi cho ăn uống sinh hoạt trong $01$ tuần được tổng hợp ở bảng sau:
	\begin{center}
		\begin{tabular}{|c|c|c|c|c|c|}
			\hline 
			Số tiền (nghìn đồng) & $[300;340)$ & $[340;380)$ & $[380;420)$ &  $[420;460)$ & $[460;500)$ \\
			 \hline
			Số học sinh & $5$ &  $9$ & $12$ & $7$ & $6$ \\
			\hline 
		\end{tabular}
	\end{center}
	Trung vị của mẫu số liệu ghép nhóm trên bằng bao nhiêu? (\textit{Kết quả được làm tròn đến hàng đơn  vị}).
	\shortans{398} 
	\loigiai{ 
		Gọi $x_1;x_2;\ldots;x_{39}$ là số tiền chi ăn uống sinh hoạt trong tuần của học sinh được xếp theo thứ tự không giảm.\\ 
		Do $x_1,\ldots,x_5\in [300;340);  x_6,\ldots,x_{14}\in [340;380); x_{15},\ldots,x_{26}\in[380;420)$.\\
		Nên trung vị của mẫu  số liệu $x_1;x_2;\ldots;x_{39}$ là $x_{20}\in[380;420)$.\\
		Ta xác định được $n=39$;  $n_m=12$; $C=5+9=14$; $u_m=380$; $u_{m+1}=420$.\\
		Vậy trung vị của mẫu số liệu  ghép nhóm là
		 \[
		 M_{\text{e}}=u_m+\dfrac{\dfrac{n}{2}-C}{n_m}\cdot(u_{m+1}-u_{m})=380+\dfrac{\dfrac{39}{2}-14}{12}\cdot(420-380) \approx 398.
		\] 
	} 
\end{ex}

% nguồn: Ok_TraLoiNgan_Lop11_Lan1
\begin{ex}%[1D5H2-2]%[ma:75f7125b]
	Ở lớp $11$A của trường  THPT, kết quả phần thi lý thuyết trong kỳ thi Nghề môn Tin học được giáo viên tổng hợp như sau:
	\begin{center}
		\begin{tabular}{|c|c|c|c|c|c|}
			 \hline
			Điểm  thi & $[6;7)$ & $[7;8)$ & $[8;9)$ & $[9;10)$ & $[10;11)$ \\
			 \hline
			Số học  sinh & $3$ & $4$ & $5$ & $2$ & $1$ \\
			\hline 
		\end{tabular}
	\end{center}
	Trung vị của mẫu số  liệu ghép nhóm trên bằng bao nhiêu? (\textit{Kết quả làm tròn đến hàng phần mười}).
	 \shortans{8,1}
	 \loigiai{
		Gọi $x_1;x_2;\ldots;x_{15}$ là điểm thi của học sinh lớp được xếp theo thứ tự không  giảm.\\
		Do $x_1,\ldots,x_3\in [6;7); x_4,\ldots,x_7\in [7;8); x_8,\ldots,x_{12}\in[8;9)$.\\ 
		Nên số trung vị của mẫu  số liệu $x_1;x_2;\ldots;x_{41}$ là $x_8\in[8;9)$.\\
		Ta xác định được $n=15$; $n_m=5$; $C=3+4=7$; $u_m=8$; $u_{m+1}=9$.\\ 
		Vậy trung vị của mẫu số  liệu ghép nhóm là
		 \[
		 M_{\text{e}}=u_m+\dfrac{\dfrac{n}{2}-C}{n_m}\cdot(u_{m+1}-u_{m})=8+\dfrac{\dfrac{15}{2}-7}{5}\cdot(9-8)=8{,}1.
		\] 
	 }
\end{ex}

% nguồn: Ok_TraLoiNgan_Lop11_Lan1
\begin{ex}%[1D5H2-2]%[ma:01064008]
	Cân nặng (đơn vị: kg) của nhóm học  sinh trường THPT được tổng hợp dưới bảng sau:
	\begin{center}
		\begin{tabular}{|c|c|c|c|c|c|}
			 \hline
			Cân  nặng & $[40;45)$ & $[45;50)$ & $[50;55)$ & $[55;60)$ & $[60;65)$ \\
			\hline 
			Số học sinh & $7$ & $5$ & $11$ & $5$ & $7$  \\
			 \hline
		\end{tabular}
	\end{center}
	Trung vị của mẫu số liệu ghép nhóm trên bằng bao nhiêu? (\textit{Kết quả  được làm tròn đến hàng phần chục}).
	\shortans{52,5} 
	 \loigiai{
		Gọi  $x_1;x_2;\ldots;x_{35}$ là cân nặng của học sinh lớp được xếp theo thứ tự không giảm.\\
		Do $x_1,\ldots,x_7\in [40;45); x_8,\ldots,x_{12}\in [45;50); x_{13},\ldots,x_{23}\in[50;55)$.\\ 
		Nên trung vị của mẫu số liệu  $x_1;x_2;\ldots;x_{35}$ là $x_{18}\in[50;55)$.\\
		Ta xác  định được $n=35$; $n_m=11$; $C=7+5=12$; $u_m=50$; $u_{m+1}=55$.\\
		Vậy trung  vị của mẫu số liệu ghép nhóm là
		\[ 
		M_{\text{e}}=u_m+\dfrac{\dfrac{n}{2}-C}{n_m}\cdot(u_{m+1}-u_{m})=50+\dfrac{\dfrac{35}{2}-12}{11}\cdot(55-50)\approx52{,}5. 
		 \]
	} 
\end{ex}

% nguồn: Ok_TraLoiNgan_Lop11_Lan2
\begin{ex}%[1D5H2-2]%[ma:3ef60742]
	Người ta ghi chép lại trọng lượng (gam) một loại cá rô được nuôi trong ao theo một chế độ đặc biệt sau $6$ tháng, họ có bảng tần số  ghép nhóm sau
	\begin{center}
		\begin{tabular}{|c|c|c|c|c|c|c|}
			 \hline
			Trọng lượng & $[60; 70)$ & $[70; 80)$ & $[80; 90)$ & $[90; 100)$ & $[100; 110)$ & $[110;  120)$\\
			\hline 
			Số cá & $13$ &  $24$ & $55$ & $6$1 & $31$ & $16$\\
			\hline 
		\end{tabular}
	\end{center}
	Tìm trung vị  của mẫu số liệu ghép nhóm đã cho (làm tròn đến hàng phần mười). \shortans{91,3}
	\loigiai{ 
		Trả lời:  \\
		Cỡ mẫu của mẫu số liệu là $n=13+24+55+61+31+16=200$.\\ 
		Gọi $x_1, x_2,\ldots, x_{200}$ là mẫu số liệu được sắp xếp theo thứ tự không giảm.\\ 
		 Trung vị của mẫu số liệu này là $\dfrac{x_{100}+x_{101}}{2}\in[90; 100)$.\\
		Ta có $n_m=61;  C=13+24+55=92; u_m=90; u_{m+1}=100$.\\
		Trung vị của mẫu số liệu ghép  nhóm là
	\[M_e=90+\dfrac{\dfrac{200}{2}-92}{61}\cdot  (100-90)=\dfrac{5570}{61}\approx 91{,}3\,\, \text{(g)}.\]}
\end{ex}

% nguồn: Ok_TraLoiNgan_Lop11_Lan2
\begin{ex}%[1D5H2-2]%[ma:6dc4189e]
	Chiều cao (đơn vị:  mét) của $35$ cây bạch đàn được cho ở bảng sau
	\begin{center}
		\begin{tabular}{|c|c|c|c|c|c|}
			 \hline
			Số đo chiều cao (m) & $[6{,}5; 7)$ &  $[7; 7{,}5)$ & $[7{,}5; 8)$ & $[8; 8{,}5)$ & $[8{,}5; 9)$\\
			\hline 
			Số cây & $6$ & $9$ & $15$ & $4$ &  $1$\\
			 \hline
		\end{tabular}
	\end{center}
	Hãy tìm trung vị  của mẫu số liệu ghép nhóm trên (làm tròn đến hàng phần mười).\shortans{7,6}
	 \loigiai{
		Trả lời:  \\
		 Gọi $x_1; x_2;\ldots; x_{35}$ là chiều cao của các cây bạch đàn xếp theo thứ tự không giảm.\\
		Do $x_1;\ldots; x_6\in[6{,}5; 7); x_7;\ldots; x_{15}\in[7; 7{,}5); x_{16};\ldots; x_{30}\in[7{,}5; 8); x_{31}; \ldots; x_{34}\in [8; 8{,}5); x_{35}\in [8{,}5; 9)$ nên trung vị của mẫu là $x_{18}\in[7,5;  8)$.\\
		Ta xác định  được $n=35; n_m=15; C=6+9=15; u_m=7{,}5; u_{m+1}=8$.\\
		Trung vị của mẫu  số liệu ghép nhóm là
	\[M_e=7{,}5+\dfrac{\dfrac{35}{2}-15}{15}\cdot (8-7{,}5)\approx  7{,}6.\]}
\end{ex}

% nguồn: Ok_TraLoiNgan_Lop11_Lan2
\begin{ex}%[1D5H2-2]%[ma:f0780365]
	Kết quả bài kiểm tra toán Giữa HKI của  lớp 11A được ghi lại ở bảng sau
	\begin{center}
		\begin{tabular}{|c|c|c|c|c|c|}
			\hline Điểm số &{$[0;  2)$}&{$[2; 4)$}&{$[4; 6)$}&{$[6; 8)$}&{$[8; 10,5)$}\\
			\hline Số học sinh & 1 & 5 & 22 & 10 & 7 \\ 
			\hline 
		\end{tabular}
	\end{center}
	Dựa vào bảng số liệu trên, giáo viên toán có thể nhận định $50 \%$ học sinh trong lớp có điểm từ bao nhiêu trở lên (làm tròn  đến hàng phần mười)? \shortans{5,5}
	 \loigiai{
		Số học  sinh trong lớp 10A1 là: $n=1+5+22+10+7=45$.\\
		Gọi $x_1; x_2; x_3; \ldots; x_{45}$ lần lượt là điểm số của các học sinh theo thứ  tự không giảm.
		 Do
		$$ 
		x_1 \in[0;  2); x_2, \ldots, x_6 \in[2; 4); x_7, \ldots, x_{28}\in[4; 6); x_{29}, \ldots, x_{38}\in[6; 8); x_{39}; \ldots; x_{45}\in[8; 10{,}5).
		$$ 
		Trung vị của dãy số liệu là  $x_{23}\in[4; 6)$.\\
		Ta có $n=45; p=3;  m_3=22; m_1+m_2=1+5=6; a_3=4; a_4=6$.\\
		Trung  vị của mẫu số liệu ghép nhóm là
		$$ 
		M_e=4+\dfrac{\dfrac{45}{2}-6}{22}\cdot(6-4)=5{,}5 . 
		$$ 
	Vậy có  $50 \%$ học sinh trong lớp $10$A1 từ $5{,}5$ điểm trở lên.}
\end{ex}

% nguồn: Ok_TraLoiNgan_Lop11_Lan2
\begin{ex}%[1D5H2-2]%[ma:7f5df69f]
	Một học viện bóng đá điều tra về lứa tuổi  của $100$ học viên trẻ đăng kí đầu tiên để tham gia khóa học mới và thu được bảng sau
	\begin{center}
		\begin{tabular}{|c|c|c|c|c|c|}
			\hline 
			Nhóm  tuổi & $[8; 9]$ & $[10; 11]$ & $[12; 13]$ & $[14; 15]$ & $[16; 17]$\\
			 \hline
			Số học viên & $14$ & $20$ & $33$ & $18$  & $15$\\
			 \hline
		\end{tabular}
	\end{center}
	Tìm trung vị của mẫu số  liệu ghép nhóm trên (làm tròn đến hàng phần mười). \shortans{12,5}
	\loigiai{ 
		Vì độ tuổi được điều tra là số nguyên nên ta hiệu chỉnh bảng số liệu trên về bảng tần số ghép  nhóm sau
		\begin{center}
			\begin{tabular}{|c|c|c|c|c|c|}
				 \hline
				Nhóm tuổi & $[7{,}5; 9{,}5)$ & $[9{,}5; 11{,}5)$ & $[11{,}5; 13{,}5)$ & $[13{,}5; 15{,}5)$  & $[15{,}5; 17{,}5)$\\
				\hline 
				Số học viên & $14$ & $20$ & $33$ & $18$ & $15$\\ 
				 \hline
			\end{tabular}
		\end{center}
		Cỡ mẫu của mẫu số  liệu là $n=100$.\\
		Gọi $x_1, x_2,\ldots, x_{100}$  là mẫu số liệu được sắp xếp theo thứ tự không giảm.\\
		Trung vị của  mẫu số liệu này là $\dfrac{x_{50}+x_{51}}{2}\in[11{,}5; 13{,}5)$.\\
		Ta có  $n_m=33; C=14+20=34; u_m=11{,}5; u_{m+1}=13{,}5$.\\
		 Trung vị của mẫu số liệu ghép nhóm là
	\[M_e=11{,}5+\dfrac{\dfrac{100}{2}-34}{33}\cdot (13{,}5-11{,}5)=\dfrac{823}{66}\approx 12{,}5.\]} 
\end{ex}

% nguồn: Ok_TraLoiNgan_Lop11_Lan2
\begin{ex}%[1D5H2-2]%[ma:9cee459c]
	Kết quả điều tra chiều cao của  học sinh trong trường THPT được biểu diễn ở biểu đồ dưới đây.
	\begin{center}
		\begin{tikzpicture}[>=stealth,line join=round,line cap=round,font=\footnotesize,scale=0.85,line width=1pt,yscale=.7]
				 \draw[->] (0,0)--(0,13) node[above]{(\text{Số học sinh})};
				 \draw[shift={(0,3)}] (0,0)--(-2pt,0) node[left]{\scriptsize $3$};
				\draw[shift={(0,5)}] (0,0)--(-2pt,0) node[left]{\scriptsize  $5$};
				\draw[shift={(0,6)}] (0,0)--(-2pt,0) node[left]{\scriptsize  $6$};
				 \draw[shift={(0,9)}] (0,0)--(-2pt,0) node[left]{\scriptsize $9$};
				\draw[shift={(0,12)}]  (0,0)--(-2pt,0) node[left]{\scriptsize $12$};
				\draw[dashed,thin,line  width=0.01pt] (0,3)--(11,3);
				 \draw[dashed,thin,line width=0.01pt] (0,5)--(11,5);
				\draw[dashed,thin,line width=0.01pt]  (0,6)--(11,6);
				 \draw[dashed,thin,line width=0.01pt] (0,9)--(11,9);
				\draw[dashed,thin,line  width=0.01pt] (0,12)--(11,12);
				\draw[line cap=round,pattern=dots,pattern color =black]  (1,0)--(1,6)--(3,6) node[midway,above]{$ $}--(3,0)--cycle;
				 \draw[line cap=round,pattern=dots,pattern color =black] (3,0)--(3,12)--(5,12) node[midway,above]{$ $}--(5,0)--cycle;
				\draw[line cap=round,pattern=dots,pattern color =black] (5,0)--(5,9)--(7,9)  node[midway,above]{$ $}--(7,0)--cycle;
				\draw[line cap=round,pattern=dots,pattern color =black] (7,0)--(7,5)--(9,5) node[midway,above]{$ $}--(9,0)--cycle; 
				\draw[line cap=round,pattern=dots,pattern color =black] (9,0)--(9,3)--(11,3)  node[midway,above]{$ $}--(11,0)--cycle;
				\node [below]  at (2,0){$ [160; 163)$};
				\node [below] at  (4,0){$ [163; 166)$};
				\node  [below] at (6,0){$ [166; 169)$};
				 \node [below] at (8,0){$ [169; 172)$};
				\node [below] at (10,0){$ [172;  173)$};
				\draw[->] (0,0) node [below  left=-2pt]{$O$}--(13,0) node[below]{(\text{cm})};
			\end{tikzpicture}
	\end{center}
	Trung vị của mẫu số liệu ghép nhóm trên là  bao nhiêu (làm tròn đến hàng đơn vị)? \shortans{166}
	\loigiai{ 
		Gọi $x_1$,  $x_2$, …, $x_{35}$ là cân nặng của các em học sinh trong trường THPT xếp theo thứ tự không giảm.\\
		Do $x_1,\ldots ,x_6\in\left[160;163\right)$; $x_7,\ldots ,x_{18}\in\left[163;166\right)$; $x_{19},\ldots ,x_{27}\in\left[166;169\right)$; $x_{28},\ldots ,x_{32}\in\left[169;172\right)$; $x_{33},x_{34},x_{35}\in\left[172;175\right)$ nên trung vị của mẫu số liệu  $x_1,x_2,\ldots ,x_{35}$ là $x_{18}\in\left[163;166\right)$.\\
		Ta xác định được $ n=35$, $n_m=12$, $ C=6$,  $u_m=163$, $u_{m+1}=166$.\\ Vậy trung vị của mẫu số liệu ghép nhóm là
	 \[M_e=163+\dfrac{\dfrac{35}{2}-6}{12}\cdot\left(166-163\right)\approx 166.\]}
\end{ex}

% nguồn: Ok_TraLoiNgan_Lop11_Lan2
\begin{ex}%[1D5H2-2]%[ma:44f6a677]
	Kết quả khảo  sát cân nặng của $25$ quả bơ (đơn vị: gam) ở một lô hàng cho trong bảng sau
	\begin{center}
		\begin{tabular}{|c|c|c|c|c|c|}
			 \hline
			Cân nặng (gam) & $\left[150;\,155\right)$ & $\left[155;\,160\right)$  & $\left[160;\,165\right)$ & $\left[165;\,170\right)$ & $\left[170;\,175\right)$\\
			\hline 
			Số quả bơ &  $1$ & $7$ & $12$ & $3$ & $2$\\
			 \hline
		\end{tabular}
	\end{center}
	Hãy tìm trung vị của mẫu số liệu  ghép nhóm trên (làm tròn đến hàng đơn vị). \shortans{162}
	\loigiai{ 
		 Cỡ mẫu là: $ n=1+7+12+3+2=25$.\\
		Gọi $x_1; x_2,\ldots ,x_{25}$  là cân nặng của $25$ quả bơ xếp theo thứ tự không giảm.\\
		Do $x_1\in\left[150;155\right);{x_2},\ldots ,x_8\in\left[155;160\right);{x_9},\ldots ,x_{20}\in\left[160;165\right); x_{21}; \ldots; x_{23} \in [165; 170); x_{24}; x_{25}\in [170; 175)$ nên trung vị  của mẫu số liệu $x_1;{x_2};\ldots;{x_{25}}$ là $x_{13}\in\left[160;165\right)$.\\
		Ta xác định được $ n=25; n_m=12; C=1+7=8; u_m=160;  u_{m+1}=165$.\\
		Suy ra, $ p=3; a_2 =160; m_1 +m_2=1+7=8; a_4  - a_3=5$ và ta có\\
		Vậy  trung vị của mẫu số liệu ghép nhóm là
	\[M_e=160+\dfrac{\dfrac{25}{2}-8}{12}\cdot\left(165-160\right)\approx 162.\]} 
\end{ex}

% nguồn: Ok_TraLoiNgan_Lop11_Lan2
\begin{ex}%[1D5H2-2]%[ma:f2d63a67]
	Biểu đồ bên dưới thống kê chiều cao của  các vận động viên của hai đội bóng rổ Sao La và Kim Ngưu.
	\begin{center}
		\begin{tikzpicture}[>=stealth,line join=round,line cap=round,font=\footnotesize,scale=0.8,line width=1pt]
				\draw[->]  (0,0)--(0,12) node[above]{(\text{Số vận động viên})};
				\draw[shift={(0,1)}] (0,0)--(-2pt,0) node[left]{\scriptsize  $1$};
				 \draw[shift={(0,2)}] (0,0)--(-2pt,0) node[left]{\scriptsize $2$};
				\draw[shift={(0,3)}] (0,0)--(-2pt,0)  node[left]{\scriptsize $3$};
				\draw[shift={(0,4)}] (0,0)--(-2pt,0)  node[left]{\scriptsize $4$};
				 \draw[shift={(0,5)}] (0,0)--(-2pt,0) node[left]{\scriptsize $5$};
				\draw[shift={(0,6)}] (0,0)--(-2pt,0) node[left]{\scriptsize  $6$};
				\draw[shift={(0,7)}] (0,0)--(-2pt,0) node[left]{\scriptsize  $7$};
				\draw[shift={(0,8)}]  (0,0)--(-2pt,0) node[left]{\scriptsize $8$};
				\draw[shift={(0,9)}] (0,0)--(-2pt,0) node[left]{\scriptsize  $9$};
				\draw[shift={(0,10)}]  (0,0)--(-2pt,0) node[left]{\scriptsize $10$};
				\draw[shift={(0,11)}] (0,0)--(-2pt,0) node[left]{\scriptsize  $11$};
				\path (5,13) node  {
				\normalsize{\textbf{Chiều cao của các  vận động viên}}
					 };
				 \draw[dashed,thin,line width=0.01pt] (0,1)--(13,1);
				\draw[dashed,thin,line  width=0.01pt] (0,2)--(13,2);
				\draw[dashed,thin,line width=0.01pt] (0,3)--(13,3); 
				\draw[dashed,thin,line  width=0.01pt] (0,4)--(13,4);
				\draw[dashed,thin,line  width=0.01pt] (0,5)--(13,5);
				\draw[dashed,thin,line width=0.01pt] (0,6)--(13,6); 
				 \draw[dashed,thin,line width=0.01pt] (0,7)--(13,7);
				\draw[dashed,thin,line width=0.01pt]  (0,8)--(13,8);
				\draw[dashed,thin,line width=0.01pt]  (0,9)--(13,9);
				\draw[dashed,thin,line width=0.01pt] (0,10)--(13,10); 
				\draw[dashed,thin,line width=0.01pt]  (0,11)--(13,11);
				\draw[line cap=round,pattern=north east lines,pattern color =yellow] (13,10)--(13,11)--(12,11)--(12,10)--(13,10)  node[above right]{\text{Đội Sao La}};
				\draw[line cap=round,pattern=dots,pattern color =black] (13,12)--(13,13)--(12,13)--(12,12)--(13,12)  node[above right]{\text{Đội Kim Ngưu}};
				\draw[line cap=round,pattern=dots,pattern color =black] (1,0)--(1,1)--(2,1) node[midway,above]{$  $}--(2,0)--cycle;
				\draw[line cap=round,pattern=north east lines,pattern color  =yellow] (2,0)--(2,1)--(3,1) node[midway,above]{$ $}--(3,0)--cycle;
				\draw[line cap=round,pattern=dots,pattern color =black] (3,0)--(3,4)--(4,4)  node[midway,above]{$ $}--(4,0)--cycle;
				\draw[line cap=round,pattern=north east lines,pattern color =yellow] (4,0)--(4,3)--(5,3) node[midway,above]{$ $}--(5,0)--cycle; 
				\draw[line  cap=round,pattern=dots,pattern color =black] (5,0)--(5,5)--(6,5) node[midway,above]{$ $}--(6,0)--cycle;
				\draw[line cap=round,pattern=north east lines,pattern color  =yellow] (6,0)--(6,4)--(7,4) node[midway,above]{$ $}--(7,0)--cycle;
				\draw[line cap=round,pattern=dots,pattern color =black] (7,0)--(7,5)--(8,5) node[midway,above]{$ $}--(8,0)--cycle; 
				\draw[line cap=round,pattern=north east lines,pattern  color =yellow] (8,0)--(8,10)--(9,10) node[midway,above]{$ $}--(9,0)--cycle;
				\draw[line  cap=round,pattern=dots,pattern color =black] (9,0)--(9,4)--(10,4) node[midway,above]{$ $}--(10,0)--cycle;
				\draw[line cap=round,pattern=north east lines,pattern color =yellow] (10,0)--(10,1)--(11,1) node[midway,above]{$ $}--(11,0)--cycle; 
				\node [below] at (2,0){$  [170; 750)$};
				\node [below] at (4,0){$  [175; 180)$};
				\node [below] at (6,0){$ [180; 185)$}; 
				\node [below] at (8,0){$ [185;  190)$};
				 \node [below] at (10,0){$ [190; 195)$};
				\draw[->] (0,0) node  [below left=-2pt]{$O$}--(13,0) node[below]{(\text{cm})};
			\end{tikzpicture}
	\end{center}
	 Tính tổng số trung vị của các vận động viên của hai đội bóng (làm tròn đến hàng đơn vị). \shortans{370}
	 \loigiai{
		Gọi $x_1,x_2,...,x_{19}$ là chiều cao của 19 thành viên đội Sao  La xếp theo thứ tự không giảm.\\
		Số trung vị của mẫu số liệu trên là  $x_{10}$ thuộc nhóm $\left[185;190\right)$.\\
		Số trung vị của mẫu  số liệu đội Sao La là \[M_o=185+\dfrac{\dfrac{19}{2}-\left(1+3+4\right)}{10}\left(190-185\right)=185{,}75.\]
		Gọi $y_1,y_2,...,y_{19}$ là chiều cao của 19 thành viên đội Kim Ngưu  xếp theo thứ tự không giảm.\\
		Số trung vị của mẫu số liệu trên là $y_{10}$  thuộc nhóm $\left[180;185\right)$.\\
		Số  trung vị của mẫu số liệu đội Kim Ngưu là \[M_o^{\prime}=180+\dfrac{\dfrac{19}{2}-\left(1+4\right)}{5}\left(185-180\right)= 184{,}5.\]
	Vậy tổng số trung vị của các vận động viên của hai đội bóng là $370$.} 
\end{ex}

% nguồn: Ok_TraLoiNgan_Lop11_Lan2
\begin{ex}%[1D5H2-2]%[ma:c8b9b44b]
	Thời gian  (phút) truy cập internet mỗi buổi tối của một số học sinh được cho trong bảng sau
	\begin{center}
		\begin{tabular}{|c|c|c|c|c|c|}
			 \hline
			Nhóm & $[9{,}5; 12{,}5)$ & $[12{,}5; 15{,}5)$ & $[15{,}5; 18{,}5)$ & $[18{,}5; 21{,}5)$ & $[21{,}5;  24{,}5)$\\
			 \hline
			Số học $\sinh$ & $3$ & $12$ & $15$ & $24$ &  $2$\\
			\hline 
		\end{tabular}
	\end{center}
	Tìm trung vị của mẫu số  liệu ghép nhóm trên (làm tròn đến hàng phần mười). \shortans{18,1}
	\loigiai{ 
		Cỡ mẫu của mẫu số liệu là $n=3+12+15+24+2=56$.\\ 
		Gọi $x_1, x_2, x_3,\ldots, x_{56}$ là thời gian truy cập Internet của lần lượt $56$ học sinh theo thứ tự không giảm.\\ 
		 Trung vị của mẫu số liệu là $\dfrac{x_{28}+x_{29}}{2}\in[15{,}5; 18{,}5)$.\\
		 Ta có: $n_m=15; C=3+12=15; u_m=15{,}5; u_{m+1}=18{,}5$.\\
		Khi đó trung vị của mẫu số  liệu phép nhóm là
	\[M_e=15{,}5+\dfrac{\dfrac{56}{2}-15}{15}\cdot (18{,}5-15{,}5)=\dfrac{181}{10}=18{,}1.\]} 
\end{ex}

% nguồn: Ok_TraLoiNgan_Lop11_Lan2
\begin{ex}%[1D5H2-2]%[ma:b619ecc1]
	Ghi lại tốc độ bóng trong $200$ lần giao bóng của một vận động viên môn quần vợt cho  kết quả như bảng bên.
	\begin{center}
		\begin{tabular}{|c|c|c|}
			 \hline
			 Tốc độ (km/h) & Số lần \\
			\hline 
			$150\le v < 155$ & $18$  \\
			 \hline
			$155\le v < 160$ &  $28$ \\
			\hline 
			$160\le v < 165$ & $35$ \\ 
			 \hline
			$165\le v <  170$ & $43$ \\
			\hline 
			$170\le  v < 175$ & $41$ \\
			\hline 
			$175\le v <  180$ & $35$ \\
			\hline 
		\end{tabular}
	\end{center}
	Tính trung vị của mẫu số liệu ghép nhóm này (làm tròn đến hàng đơn vị).  \shortans{167}
	 \loigiai{
		Cỡ mẫu là $  n=18+28+35+43+41+35=200$.\\
		Gọi $ x_1, x_2\ldots x_{200}$ là tốc độ giao bóng của $200$ lần và giả  sử dãy này được sắp xếp theo thứ tự tăng dần.\\ Khi đó trung vị là $\dfrac{x_{100}+x_{101}}{2}$.\\
		Do hai giá trị $ x_{100}, x_{101}$ thuộc  nhóm $[165;170)$ nên nhóm này chứa trung vị.\\
		Suy ra  $ p=4; a_4=165; m_4=43; m_1+m_2+m_3=18+28+35=81; a_5-a_4=5$ và ta có\\
	\[M_e=165+\dfrac{\dfrac{200}{2}-81}{43}\cdot 5\approx  167.\]}
\end{ex}

% nguồn: Ok_TraLoiNgan_Lop11_Lan2
\begin{ex}%[1D5H2-2]%[ma:d453d031]
	Điểm thi giữa kỳ $1$ môn toán của một lớp học sinh khối $11$ thu được  mẫu số liệu ghép nhóm sau
	\centerline{\begin{tabular}{|c|c|c|c|}
	 \hline
	Điểm thi & $\left[1{,}5;4{,}5\right)$ & $\left[4{,}5;7{,}5\right)$ & $\left[7{,}5;10{,}5\right)$\\ 
	\hline 
	Số  học sinh & $7$ & $18$ & $10$\\
	 \hline
\end{tabular}}
Trung vị của mẫu số liệu ghép nhóm trên bằng bao nhiêu (làm tròn đến hàng phần mười)? \shortans{6,3} 
\loigiai{ 
	Cỡ mẫu là $n=7+18+10=35.$\\ 
	Gọi $x_1,x_2,...,x_{35}$ là số điểm của 35 học  sinh và giả sử dãy này được sắp xếp theo thứ tự không giảm. Khi đó, trung vị là $x_{18}$ thuộc nhóm $\left[4{,}5;7{,}5\right)$.\\
	Ta xác định được $n=35; n_m=18; C=7;  u_m=4{,} 5; u_{m+1}=7{,}5$.\\
	Trung vị của mẫu số liệu ghép nhóm là 
 \[M_e=4{,}5+\dfrac{\dfrac{35}{2}-7}{18}\left(7{,}5-4{,}5\right)\approx 6{,}3.\]}
\end{ex}

% nguồn: Ok_TraLoiNgan_Lop11_Lan2
\begin{ex}%[1D5H2-2]%[ma:92742562]
	Biểu đồ bên dưới thống  kê chiều cao của các vận động viên của hai đội bóng rổ Sao La và Kim Ngưu.
	\begin{center}
		\begin{tikzpicture}[>=stealth,line join=round,line cap=round,font=\footnotesize,scale=0.8,line width=1pt]
				\draw[->]  (0,0)--(0,12) node[above]{(\text{Số vận động viên})};
				\draw[shift={(0,1)}] (0,0)--(-2pt,0)  node[left]{\scriptsize $1$};
				\draw[shift={(0,2)}]  (0,0)--(-2pt,0) node[left]{\scriptsize $2$};
				 \draw[shift={(0,3)}] (0,0)--(-2pt,0) node[left]{\scriptsize $3$};
				\draw[shift={(0,4)}]  (0,0)--(-2pt,0) node[left]{\scriptsize $4$};
				\draw[shift={(0,5)}] (0,0)--(-2pt,0) node[left]{\scriptsize $5$}; 
				\draw[shift={(0,6)}] (0,0)--(-2pt,0) node[left]{\scriptsize  $6$};
				\draw[shift={(0,7)}] (0,0)--(-2pt,0) node[left]{\scriptsize $7$}; 
				\draw[shift={(0,8)}] (0,0)--(-2pt,0) node[left]{\scriptsize $8$}; 
				 \draw[shift={(0,9)}] (0,0)--(-2pt,0) node[left]{\scriptsize $9$};
				\draw[shift={(0,10)}] (0,0)--(-2pt,0) node[left]{\scriptsize  $10$};
				\draw[shift={(0,11)}] (0,0)--(-2pt,0)  node[left]{\scriptsize $11$};
				\path  (5,13) node {
				\normalsize{\textbf{Chiều cao của các  vận động viên}}
					 };
				\draw[dashed,thin,line width=0.01pt]  (0,1)--(13,1);
				\draw[dashed,thin,line width=0.01pt]  (0,2)--(13,2);
				 \draw[dashed,thin,line width=0.01pt] (0,3)--(13,3);
				\draw[dashed,thin,line  width=0.01pt] (0,4)--(13,4);
				\draw[dashed,thin,line  width=0.01pt] (0,5)--(13,5);
				 \draw[dashed,thin,line width=0.01pt] (0,6)--(13,6);
				\draw[dashed,thin,line width=0.01pt]  (0,7)--(13,7);
				\draw[dashed,thin,line width=0.01pt] (0,8)--(13,8); 
				 \draw[dashed,thin,line width=0.01pt] (0,9)--(13,9);
				\draw[dashed,thin,line width=0.01pt] (0,10)--(13,10); 
				\draw[dashed,thin,line width=0.01pt] (0,11)--(13,11); 
				\draw[line cap=round,pattern=north  east lines,pattern color =yellow] (13,10)--(13,11)--(12,11)--(12,10)--(13,10) node[above right]{\text{Đội Sao La}};
				 \draw[line cap=round,pattern=dots,pattern color =black] (13,12)--(13,13)--(12,13)--(12,12)--(13,12) node[above right]{\text{Đội Kim Ngưu}};
				 \draw[line cap=round,pattern=dots,pattern color =black] (1,0)--(1,1)--(2,1) node[midway,above]{$ $}--(2,0)--cycle;
				\draw[line cap=round,pattern=north east lines,pattern  color =yellow] (2,0)--(2,1)--(3,1) node[midway,above]{$ $}--(3,0)--cycle;
				\draw[line cap=round,pattern=dots,pattern color  =black] (3,0)--(3,4)--(4,4) node[midway,above]{$ $}--(4,0)--cycle;
				\draw[line cap=round,pattern=north east lines,pattern color =yellow] (4,0)--(4,3)--(5,3) node[midway,above]{$ $}--(5,0)--cycle; 
				 \draw[line cap=round,pattern=dots,pattern color =black] (5,0)--(5,5)--(6,5) node[midway,above]{$ $}--(6,0)--cycle;
				\draw[line cap=round,pattern=north east lines,pattern color =yellow] (6,0)--(6,4)--(7,4)  node[midway,above]{$ $}--(7,0)--cycle;
				\draw[line cap=round,pattern=dots,pattern  color =black] (7,0)--(7,5)--(8,5) node[midway,above]{$ $}--(8,0)--cycle;
				\draw[line cap=round,pattern=north east lines,pattern color  =yellow] (8,0)--(8,10)--(9,10) node[midway,above]{$ $}--(9,0)--cycle;
				\draw[line cap=round,pattern=dots,pattern color =black] (9,0)--(9,4)--(10,4)  node[midway,above]{$ $}--(10,0)--cycle;
				\draw[line cap=round,pattern=north east lines,pattern  color =yellow] (10,0)--(10,1)--(11,1) node[midway,above]{$ $}--(11,0)--cycle;
				\node [below] at (2,0){$ [170; 175)$}; 
				\node [below] at  (4,0){$ [175; 180)$};
				\node [below] at (6,0){$ [180; 185)$}; 
				\node [below] at (8,0){$ [185;  90)$};
				 \node [below] at (10,0){$ [190; 195)$};
				\draw[->] (0,0) node [below left=-2pt]{$O$}--(13,0) node[below]{(\text{cm})}; 
			\end{tikzpicture}
	\end{center}
	Tính tổng chiều cao trung bình của các vận  động viên của hai đội bóng (làm tròn đến hàng đơn vị). \shortans{369}
	\loigiai{ 
		Ta có bảng tần số ghép nhóm  như sau\\
		\centerline{\begin{tabular}{|c|c|c|c|c|c|}
		 \hline
		Chiều cao  & $\left[170;175\right)$ & $\left[175;180\right)$ & $\left[180;185\right)$ & $\left[185;190\right)$ & $\left[190;195\right)$\\
		 \hline
		Chiều cao đại diện & $172{,}5$ & $177{,}5$ & $182{,}5$ & $187{,}5$ &  $192{,}5$\\
		\hline 
		Đội  Kim Ngưu & $1$ & $4$ & $5$ & $5$ & $4$\\
		 \hline
		Đội Sao La & $1$ & $3$ & $4$ & $10$ & $1$\\ 
		 \hline
	\end{tabular}}
	Theo số trung bình 
	\begin{itemize}
		 \item
		Chiều cao trung bình của các vận động viên đội Sao La xấp xỉ  bằng
		\[\overline{x_1}=\dfrac{1\cdot 172{,}5+4\cdot 177{,}5+5\cdot 182{,}5+5\cdot 187{,}5+4\cdot 192{,}5}{1+4+5+5+4}\approx  184{,}3\,\, \text{(cm)}.\]
		\item Chiều cao trung bình của các  vận động viên đội Kim Ngưu xấp xỉ bằng
		\[\overline{x_2}=\dfrac{1\cdot 172{,}5+3\cdot 177{,}5+4\cdot  182{,}5+10\cdot 187{,}5+1\cdot 192{,}5}{1+3+4+10+1}\approx 184{,}3\,\, \text{(cm)}.\]
	\end{itemize}
Vậy tổng chiều cao trung bình của các vận  động viên của hai đội bóng là $184{,}3+184{,}3\approx 387$.}
\end{ex}

% nguồn: Ok_TraLoiNgan_Lop11_Lan2
\begin{ex}%[1D5H2-2]%[ma:47dba55e]
	Kết quả thu thập điểm số môn Toán của $25$ học sinh khi tham gia kì thi học sinh giỏi  Toán khối $11$ (thang điểm $20$) cho ta bảng tần số ghép nhóm sau
	\begin{center}
		\begin{tabular}{|c|c|c|c|c|c|}
			\hline 
			Nhóm & $[0; 4)$ & $[4; 8)$ & $[8; 12)$ & $[12; 16)$ &  $[16; 20)$\\
			 \hline
			Số học $\sinh$  & $1$ & $7$ & $12$ & $3$ & $2$\\
			\hline 
		\end{tabular}
	\end{center}
	Tìm trung vị của  mẫu số liệu ghép nhóm trên (làm tròn đến hàng phần mười). \shortans{9,5}
	 \loigiai{
		Cỡ mẫu của mẫu số liệu là $n=25$. Gọi $x_1, x_2, x_3,\ldots, x_{25}$ là điểm số của $25$ học sinh trong kì thi đó và dãy này được sắp  xếp theo thứ tự không giảm.\\
		Trung vị của  mẫu số liệu là $x_{13}\in[8; 12)$.\\
		Ta có  $n=25, n_m=12, C=1+7=8, u_m=8, u_{m+1}=12$.\\
		Trung vị mẫu số liệu ghép  nhóm là \\
	\[M_e=u_m+\dfrac{\dfrac{n}{2}-C}{n_m}\left(u_{m+1}-u_m\right)=8+\dfrac{\dfrac{25}{2}-8}{12}(12-8)=9{,}5.\]} 
\end{ex}

% nguồn: Ok_TraLoiNgan_Lop11_Lan2
\begin{ex}%[1D5H2-2]%[ma:120906da]
	Lợi nhuận bán hàng trong $20$ ngày được lựa chọn ngẫu nhiên trong tháng $8$ tại một cửa hàng được ghi lại ở  bảng sau
	\begin{center}
		\begin{tabular}{|c|c|c|c|c|c|}
			\hline 
			Doanh thu & $\left[7;\,10\right)$ & $\left[10;\,13\right)$ & $\left[13;\,16\right)$ & $\left[16;\,19\right)$ & $\left[19;\,22\right)$\\ 
			 \hline
			Số ngày  & $3$ & $4$ & $6$ & $5$ & $2$\\
			\hline 
		\end{tabular}
	\end{center}
	Tìm trung  vị của mẫu số liệu trên (làm tròn đến hàng phần mười)? \shortans{14,5}
	\loigiai{ 
		Gọi $x_1,\,x_2,\ldots,x_{20}$ là doanh thu trong $20$ ngày  được xếp theo thứ tự không giảm.\\
		Ta  có\\
		 $x_1,\,x_2,\,x_3\in\left[7;\,10\right);\,\,x_4,\,....,\,x_7\in\left[10;\,13\right);\,\,x_8,\,....,\,x_{13}\in\left[13;\,16\right);\,;\,\,x_{14},\,....,\,x_{18}\in\left[16;\,19\right);\,x_{19},\,x_{20}\in\left[19;\,22\right)$\\
		nên trung vị của mẫu  số liệu $x_1,\,x_2,\,...,\,x_{20}$ là $\dfrac{1}{2}\left(x_{10}+x_{11}\right)\in\left[13;\,16\right)$.\\
		Suy ra $\left[13;\,16\right)$ là nhóm  chứa trung vị. Khi đó ta xác định được\\
		$n=20,\,\,n_m=6,\,C=3+4=7,\,u_m=13,\,u_{m+1}=16$.\\ 
		Vậy trung vị của mẫu số liệu ghép  nhóm là\\
	\[M_e=u_m+\dfrac{\dfrac{n}{2}-C}{n_m}\cdot \left(u_{m+1}-u_m\right)=13+\dfrac{\dfrac{20}{2}-7}{6}\cdot  \left(16-13\right)=14,5.\]}
\end{ex}

% nguồn: Ok_TraLoiNgan_Lop11_Lan2
\begin{ex}%[1D5H2-2]%[ma:3b2a97ea]
	Kết quả khảo sát hàm lượng vitamin C của một  số loại trái cây cho ở bảng sau
	\begin{center}
		\begin{tabular}{|c|c|c|c|c|c|}
			\hline 
			Hàm lượng vitamin C $(\%)$ & $[6; 7)$ & $[7; 8)$ & $[8; 9)$ & $[9; 10)$  & $[10; 11)$\\
			\hline 
			Số lượng & $3$ & $4$ & $5$  & $2$ & $1$\\
			\hline 
		\end{tabular}
	\end{center}
	 Hãy tìm trung vị của mẫu số liệu ghép nhóm trên (làm tròn đến hàng phần mười). \shortans{8,1}
	\loigiai{ 
		 Gọi $x_1; x_2;\ldots; x_{15}$ là hàm lượng vitamin $C$ của một số loại trái cây xếp theo thứ tự không giảm.\\
		Do $x_1; x_2; x_3\in[6; 7); x_4;\ldots; x_7\in[7; 8); x_8;\ldots; x_{12}\in[8; 9)$ nên trung vị của mẫu  là $x_8$ $\in[8; 9)$.\\
		Ta xác định được $n=15,  n_m=5, C=3+4=7, u_m=8, u_{m+1}=9$.\\
		Trung vị của mẫu số liệu ghép nhóm là 
	 \[M_e=8+\dfrac{\dfrac{15}{2}-7}{5}\cdot (9-8)\approx 8{,}1\%.\]}
\end{ex}

% nguồn: Ok_TraLoiNgan_Lop11_Lan2
\begin{ex}%[1D5H2-2]%[ma:b27f737d]
	Bảng số liệu dưới đây thể hiện cân nặng  của các thành viên trong một câu lạc bộ thể thao.
	\begin{center}
		\begin{tabular}{|c|c|c|c|c|c|c|}
			 \hline
			Cân nặng &  $\left[40;50\right)$ & $\left[50;60\right)$ & $\left[60;70\right)$ & $\left[70;80\right)$ & $\left[80;90\right)$ & $\left[90;100\right)$\\
			 \hline
			Tần số  & $6$ & $8$ & $12$ & $14$ & $7$ & $3$\\
			 \hline
		\end{tabular}
	\end{center}
	Tìm trung vị của mẫu  số liệu ghép nhóm này (làm tròn đến hàng phần mười).
	\shortans{69,2} 
	\loigiai{ 
		Cỡ mẫu là  $ n=6+8+12+14+7+3=50$.\\
		Gọi $x_1,\ldots ,x_{50}$ là cân nặng của $50 $ thành viên và giả sử dãy này đã được xếp theo thứ tự tăng dần. Khi đó trung vị sẽ  nằm ở vị trí $\dfrac{x_{25}+x_{26}}{2}$.\\
		 Do 2 giá trị $x_{25},\,x_{26}$ thuộc nhóm $\left[60;70\right)$ nên nhóm này chứa trung vị.\\
		Ta  có $ p=3;{a_3}=60;{m_3}=12;{m_1}+m_2=6+8=14;{a_4}-a_3=70-60=10$.\\
	Vậy trung vị là $M_e=60+\dfrac{\dfrac{50}{2}-14}{12}\cdot 10=69{,}2$.} 
\end{ex}

% nguồn: Ok_TraLoiNgan_Lop11_Lan2
\begin{ex}%[1D5H2-2]%[ma:879e08af]
	Nghiên cứu thời gian chạy 1 vòng  sân trường $300$ m của $41$ học sinh lớp $11$A trường THPT được giáo viên bộ môn Thể dục ghi lại, có kết quả sau
	\begin{center}
		\begin{tabular}{|c|c|c|c|c|c|}
			 \hline
			Thời gian (giây) & $[40; 45)$ & $[45; 50)$ & $[50; 55)$ & $[55; 60)$ & $[60; 65)$\\ 
			\hline 
			Số học sinh &  $5$ & $8$ & $13$ & $9$ & $6$\\
			\hline 
		\end{tabular}
	\end{center}
	Hãy tìm trung  vị của mẫu số liệu ghép nhóm trên (làm tròn đến hàng phần mười). \shortans{52,9}
	\loigiai{ 
		Gọi $x_1; x_2;\ldots; x_{41}$ là thời gian chạy của học sinh lớp $11~A$ được xếp theo thứ tự không giảm. Do $x_1; x_2;\ldots; x_5\in[40; 45); x_6;\ldots; x_{13}\in[45; 50); x_{14};\ldots; x_{26}\in[50; 55)$ nên trung vị của mẫu  là $x_{21}\in[50; 55)$.\\
		Ta xác định được $n=41, n_m=13, C=5+8=13, u_m=50, u_{m+1}=55$.\\ 
		Trung vị của mẫu số liệu ghép nhóm  là
	\[M_e=50+\dfrac{\dfrac{41}{2}-13}{13}\cdot (55-50)\approx 52{,}9\,\, \text{(giây)}.\]} 
\end{ex}

% nguồn: Ok_TraLoiNgan_Lop11_Lan2
\begin{ex}%[1D5H2-2]%[ma:8026b718]
	Điều tra về số lượng học sinh khối  $11$ trong một lớp học, người ta thu được dữ liệu của $100$ học sinh và có bảng phân phối tần số ghép nhóm sau
	\begin{center}
		\begin{tabular}{|c|c|c|c|c|c|}
			\hline 
			Nhóm & $[36; 38)$ & $[38; 40)$ & $[40; 42)$ &  $[42; 44)$ & $[44; 46)$\\
			 \hline
			Tần số &  $9$ & $15$ & $25$ & $30$ & $21$\\
			\hline 
		\end{tabular}
	\end{center}
	Tìm trung vị của mẫu số liệu ghép nhóm trên (làm tròn  đến hàng phần mười). \shortans{42,1}
	\loigiai{ 
		Cỡ mẫu của mẫu  số liệu là $n=100$.\\
		Gọi $x_1, x_2, x_3,\ldots,  x_{100}$ là số học sinh trong một lớp học khối $11$ được điều tra được sắp xếp theo thứ tự không giảm.\\
		Trung vị của mẫu số liệu là $\dfrac{x_{50}+x_{51}}{2}\in[42;  44)$.\\
		Ta  có $n_m=30; C=9+15+25=49; u_m=42; u_{m+1}=44$.\\
		Trung vị của mẫu số liệu ghép nhóm  là
	\[M_e=42+\dfrac{\dfrac{100}{2}-49}{30}(44-42)=\dfrac{631}{15}\approx  42{,}07.\]}
\end{ex}

% nguồn: Ok_TraLoiNgan_Lop11_Lan2
\begin{ex}%[1D5H2-2]%[ma:864afb9c]
	Thời gian truy cập Internet mỗi buổi  tối của một số học sinh được cho trong bảng sau
	\begin{center}
		\begin{tabular}{|c|c|c|c|c|c|}
			 \hline
			Thời gian (phút) & $\left[9{,}5;\,12{,}5\right)$ & $\left[12{,}5;\,15{,}5\right)$ & $\left[15{,}5;\,18{,}5\right)$  & $\left[18{,}5;\,21{,}5\right)$ & $\left[21{,}5;\,24{,}5\right)$\\
			\hline 
			Số  học sinh & $3$ & $12$ & $15$ & $24$ & $2$\\
			 \hline
		\end{tabular}
	\end{center}
	 Tính trung vị của mẫu số liệu ghép nhóm này (làm tròn đến hàng phần mười).\shortans{18,1}
	 \loigiai{
		Cỡ mẫu là $ n=3+12+15+24+2=56$.\\ 
		Gọi $ x_1,\ldots, x_{56}$ là thời gian vào Internet của 56 học sinh và giả sử dãy này đã được sắp xếp theo thứ tự tăng dần.\\ 
		Khi đó, trung vị là $\dfrac{x_{28}+x_{29}}{2}$. Do 2 giá trị $  x_{28}, x_{29}$ thuộc nhóm $[15,5; 18,5)$ nên nhóm này chứa trung vị.\\
		Do đó, $ p=3; a_3=15,5; m_3=15;  m_1+m_2=3+12=15; a_4-a_3=3$ và ta có\\
	 $ M_e=15,5+\dfrac{\dfrac{56}{2}-15}{15}\cdot 3=18{,}1.$}
\end{ex}

% nguồn: Ok_TraLoiNgan_Lop11_Lan2
\begin{ex}%[1D5H2-2]%[ma:8762bb98]
	Khảo sát thời gian truy cập Internet  mỗi buổi tối của một số học sinh khối $11$ được cho trong bảng sau
	\begin{center}
		\begin{tabular}{|c|c|c|c|c|c|}
			 \hline
			Thời gian & $\left[9{,}5;19{,}5\right)$ & $\left[19{,}5;29{,}5\right)$ &  $\left[29{,}5;39{,}5\right)$ & $\left[39{,}5;49{,}5\right)$ & $\left[49{,}5;59{,}5\right)$\\
			 \hline
			Số học sinh &  $3$ & $10$ & $21$ & $26$ & $14$\\
			 \hline
		\end{tabular}
	\end{center}
	Trung vị của mẫu số  liệu ghép nhóm trên bằng bao nhiêu (làm tròn đến hàng phần mười)? \shortans{40,7}
	\loigiai{ 
		Cỡ mẫu là $n=3+10+21+26+14=74$.\\ 
		Gọi $x_1,x_2,...,x_{74}$ là thời gian của $74$ học sinh truy cập Internet mỗi buổi tối và giả sử dãy này  được sắp xếp theo thứ tự không giảm. Khi đó, trung vị là $\dfrac{x_{37}+x_{38}}{2}$ thuộc nhóm $\left[39{,}5;49{,}5\right)$.\\
		 Ta xác định được $n=74;\,\, n_m=26;\,\, C=34;\,\, u_m=39{,}5;\,\, u_{m+1}=49{,}5$.\\
		Trung vị của mẫu số  liệu ghép nhóm là
	\[M_e=39{,}5+\dfrac{\dfrac{74}{2}-34}{26}\left(49{,}5-39{,}5\right)=\dfrac{1057}{26}\approx 40{,}7.\]} 
\end{ex}

% nguồn: Ok_TraLoiNgan_Lop11_Lan2
\begin{ex}%[1D5H2-2]%[ma:917b1878]
	Số tiền mà học  sinh lớp $11$ chi cho ăn uống sinh hoạt trong 01 tuần được tổng hợp ở bảng sau
	\begin{center}
		\begin{tabular}{|c|c|c|c|c|c|}
			 \hline
			Số tiền (nghìn đồng) & $[300;  340)$ & $[340; 380)$ & $[380; 420)$ & $[420; 460)$ & $[460; 500)$\\
			\hline 
			Số học sinh & $5$  & $9$ & $12$ & $7$ & $6$\\
			 \hline
		\end{tabular}
	\end{center}
	Hãy tìm trung vị của mẫu số liệu ghép nhóm  trên (làm tròn đến hàng đơn vị). \shortans{398}
	 \loigiai{
		Gọi $x_1; x_2;\ldots; x_3$ là số tiền chi ăn uống sinh hoạt trong  tuần của học sinh xếp theo thứ tự không giảm.\\
		Do  $x_1;\ldots; x_5\in[300; 340); x_6;\ldots; x_9\in[340; 380); x_{15};\ldots; x_{26}\in[380; 420); x_{27}; \ldots; x_{33}\in [420; 460)$; $x_{34}; \ldots; x_{39}\in [460; 500)$ nên trung vị của mẫu là $x_{19}\in[380; 420)$.\\
		Ta xác định  được $n=39, n_m=12; C=5+9=14; u_m=380; u_{m+1}=420$.\\
		Trung vị của mẫu  số liệu ghép nhóm là
	\[M_e=380+\dfrac{\dfrac{39}{2}-14}{12}\cdot  (420-380)\approx 398.\]}
\end{ex}

% nguồn: Ok_TraLoiNgan_Lop11_Lan2
\begin{ex}%[1D5H2-2]%[ma:be6dce66]
	Trong một hội thao, thời gian chạy $200$ m của một nhóm vận  động viên được ghi lại trong bảng sau
	\begin{center}
		\begin{tabular}{|c|c|c|c|c|c|}
			\hline Thời gian (giây)  &{$[21; 21,5)$}&{$[21,5; 22)$}&{$[22; 22,5)$}&{$[22,5; 23)$}&{$[23; 23,5)$}\\
			\hline Số vận động viên & $5$ & $12$ & $32$ &  $45$ & $30$ \\
			\hline 
		\end{tabular}
	\end{center}
	Dựa vào bảng dữ liệu trên,  ban tổ chức muốn chọn ra khoảng $50 \%$ số vận động viên chạy nhanh nhất để tiếp tục thi vòng 2. Ban tổ chức nên chọn các vận động viên có thời gian chạy không quá bao nhiêu giây (làm tròn đến hàng phần mười)? \shortans{22,6}
	\loigiai{ 
		Cỡ mẫu là  $n=5+12+32+45+30=124$.
		Gọi $x_1,\ldots, x_{124}$ là thời gian chạy của $124$ vận động viên tham gia hội thao và giả sử dãy này được  sắp xếp theo thứ tự tăng dần.
		Khi đó, trung vị là $\dfrac{x_{62}+x_{63}}{2}$. Do 2 giá trị $x_{62}, x_{63}$ thuộc nhóm $[22{,}5; 23)$ nên nhóm  này chứa trung vị.
		 \\
		Suy  ra
		 \\
		 $$
		p=4;  a_4=22{,}5; m_4=45; m_1+m_2+m_3=5+12+32=49; a_5-a_4=23-22{,}5=0{,}5.
		 $$
	Ta có trung vị $M_e=22{,}5+\dfrac{\dfrac{124}{2}-49}{45}\cdot 0{,}5 \approx 22{,}6$.} 
\end{ex}

% nguồn: Ok_TraLoiNgan_Lop11_Lan1
\begin{ex}%[1D5V2-2]%[ma:60b2d4b2]
	Thời gian (phút) truy  cập internet mỗi buổi tối của một số học sinh được cho trong bảng sau\\
	\centerline{\begin{tabular}{|c|c|c|c|c|c|}
	 \hline
	Nhóm & $[9{,}5; 12{,}5)$ & $[12{,}5; 15{,}5)$ &  $[15{,}5; 18{,}5)$ & $[18{,}5; 21{,}5)$ & $[21{,}5; 24{,}5)$\\
	 \hline
	 Số học sinh & $3$ & $12$ & $15$ & $24$ & $2$\\
	\hline 
\end{tabular}}\\
Tìm trung vị của mẫu số liệu ghép nhóm  trên.
\shortans{18,1} 
\loigiai{ 
	 Cỡ mẫu của mẫu số liệu là $n=3+12+15+24+2=56$.\\
	Ta có  $\dfrac{n}{2}=28$, $cf_2=15<\dfrac{n}{2}<cf_3=30$ suy ra trung vị nằm trong nhóm $[15{,}5; 18{,}5)$.\\
	Khi  đó trung vị của mẫu số liệu phép nhóm là
 $$M_e=15{,}5+\dfrac{28-15}{15}(18{,}5-15{,}5)=\dfrac{181}{10}=18{,}1.$$}
\end{ex}

% nguồn: Ok_TraLoiNgan_Lop11_Lan1
\begin{ex}%[1D5V2-2]%[ma:f416c335]
	Cho  mẫu số liệu ghép nhóm về chiều cao (đơn vị: cm) của $25$ cây dừa giống như sau:
	\begin{center}
		\begin{tabular}{|l|c|c|c|c|c|}
			\hline 
			Chiều cao  & {$[0; 10)$} & {$[10; 20)$} & {$[20; 30)$} & {$[30; 40)$} & {$[40; 50)$} \\
			\hline 
			Số cây & $4$& $6$& $7$& $5$&  $3$\\
			\hline 
		\end{tabular}
	\end{center}
	 Tính trung vị của mẫu số liệu ghép nhóm này (Kết quả làm tròn đến chữ số thập phân thứ $1$)
	 \shortans{23,6}
	 \loigiai{
		Cở mẫu là  $n=4+6+7+5+3=25$.\\
		Gọi $x_1,\ldots,x_{25}$ là chiều cao của $25$ cây dừa và giả sử dãy này được sắp xếp  theo thư tự tăng dần. Khi đó, trung vị là $\dfrac{x_{12}+x_{13}}{2}$. Do đó hai giá trị $x_{12}$ và $x_{13}$ thuộc nhóm $[20; 30)$ nên nhóm này chứa trung vị. Do đó, $p=3$, $a_3=20$, $m_3=7$, $m_1+m_2=4+6=10$, $a_4-a_3=30-20=10$.
		\[M_e=20+\dfrac{\dfrac{25}{2}-10}{7}\cdot 10\approx 23{,}6\] 
	} 
\end{ex}


%%%%%%%%%%%%%%%% TỰ LUẬN %%%%%%%%%%%%%%%%

% nguồn: Ok_TuanLuan_Toan11_Lan1
\begin{bt}%[1D5H2-2]%[ma:4f168dee]
	Ghi lại tốc độ bóng trong $200$ lần giao bóng của một vận động viên môn quần vợt cho kết quả  như bảng sau
	\begin{center}
		\begin{tabular}{|c|c|}
			\hline  Tốc độ $v$ (km/h)&Số lần\\
			\hline  $150\le v<155$&$18$\\
			 \hline $155\le v<160$&$28$\\
			\hline $160\le v<165$&$35$\\ 
			\hline $165\le v<170$&$43$\\ 
			\hline $170\le v<175$&$41$\\ 
			\hline $175\le v<180$&$35$\\ 
			\hline 
		\end{tabular}
	\end{center}
	Tính trung vị của mẫu số liệu ghép nhóm này. 
	\loigiai{ 
		 Cỡ mẫu là $n=200$.\\
		Gọi ${x_1},\ldots,{x_{56}}$ là tốc độ giao bóng của $200$ lần giao bóng và giả  sử dãy này được sắp xếp theo thứ tự tăng dần.\\
		Khi đó, trung vị là $\dfrac{{x_{100}}+{x_{101}}}{2}$.\\ 
		Do $2$ giá trị ${x_{100}}$,  ${x_{101}}$ thuộc nhóm $[165;170)$ nên nhóm này chứa trung vị.\\
		Do đó, $p=4$, ${a_4}=165$, ${m_3}=43$, ${m_1}+{m_2}+{m_3}=18+28+35=81$, ${a_5}-{a_4}=5$ và  ta có ${M_e}=165+\dfrac{\dfrac{200}{2}-81}{43}\cdot 5=167{,}21$.
	} 
\end{bt}

% nguồn: Ok_TuanLuan_Toan11_Lan1
\begin{bt}%[1D5H2-2]%[ma:16b3d922]
	Điều tra điểm môn  Văn của một nhóm học sinh, ta được bảng số liệu như sau
	\begin{center}
		\begin{tabular}{|l|l|l|l|l|l|l|l|l|}
			\hline  $5$&$6$&$6$&$7$&$7$&$8$&$8$&$8{,}5$&$9$\\
			\hline 
		\end{tabular}
	\end{center}
	 Tính số trung vị của bảng nói trên.
	 \loigiai{
		Ta có  $N=9$
		là số lẻ.\\ 
		 Số liệu thứ $\dfrac{N+1}{2}=5$ là số trung vị.\\
	Do đó, số  trung vị là $M_e=7$ (điểm).}
\end{bt}

% nguồn: Ok_TuanLuan_Toan11_Lan1
\begin{bt}%[1D5H2-2]%[ma:86fcef9a]
	 Tốc độ phát triển của một loại Virus trong $10$ ngày với các điều kiện khác nhau (đơn vị nghìn con) được thống kê như sau
	\begin{center}
		\begin{tabular}{|l|l|l|l|l|l|l|l|l|l|}
			 \hline $20$&$100$&$30$&$980$&$440$&$20$&$20$&$150$&$60$&$270$\\
			 \hline
		\end{tabular}
	\end{center}
	Trong  trường hợp này ta chọn số nào dưới đây là giá trị đại diện là tốt nhất? Tính giá trị đại diện đó.
	\loigiai{ 
		Ta chọn số trung vị làm đại diện là tốt nhất vì có sự chênh lệch lớn giữa các số  liệu trong mẫu.\\
		Sắp xếp  lại số liệu mẫu
		\begin{center}
			\begin{tabular}{|c|c|c|c|c|c|c|c|c|c|}
				\hline  $20$&$20$&$20$&$30$&$60$&$100$&$150$&$270$&$440$&$980$\\
				 \hline
			\end{tabular}
		\end{center}
		Kích  thước mẫu là $N=10$. Số liệu thứ $5$ và $6$ lần lượt là $60$, $100$.\\
		Vậy  giá trị đại diện cho bảng số liệu là ${M_e}=\dfrac{60+100}{2}=80$ (nghìn con).
	} 
\end{bt}

% nguồn: Ok_TuanLuan_Toan11_Lan1
\begin{bt}%[1D5H2-2]%[ma:b124d5e6]
	 Điều tra số học sinh của $30$ lớp học, ta được bảng số liệu như sau
	\begin{center}
		\begin{tabular}{|l|l|l|l|l|l|l|l|l|l|l|l|l|l|l|}
			 \hline $35$&$39$&$39$&$40$&$40$&$41$&$41$&$41$&$41$&$44$&$44$&$45$&$45$&$45$&$46$\\
			\hline  $48$&$48$&$48$&$48$&$49$&$49$&$49$&$49$&$49$&$49$&$50$&$50$&$50$&$50$&$51$\\
			 \hline
		\end{tabular}
	\end{center}
	 Tính số trung vị của bảng nói trên.
	\loigiai{Ta có $N=30$ là số chẵn. Số liệu thứ $15$ và $16$ lần lượt là  $46$, $48$.\\
	Vậy số trung vị là ${M_e}=\dfrac{46+48}{2}=47$  (học sinh).}
\end{bt}

% nguồn: Ok_TuanLuan_Toan11_Lan1
\begin{bt}%[1D5H2-2]%[ma:75ef4765]
	Điều tra số học sinh giỏi khối $10$ của $15$  trường cấp ba trên địa bàn tỉnh $A$, ta được bảng số liệu như sau
	\begin{center}
		\begin{tabular}{|l|l|l|l|l|l|l|l|l|l|l|l|l|l|l|}
			 \hline $22$&$29$&$29$&$29$&$30$&$31$&$32$&$32$&$33$&$34$&$34$&$35$&$35$&$35$&$36$\\
			 \hline
		\end{tabular}
	\end{center}
	Tính số trung vị của bảng nói  trên.
	 \loigiai{
		Ta có $N=15$  là số lẻ. Số liệu thứ $\dfrac{15+1}{2}=8$.\\
		Vậy số trung vị là ${M_e}=32$ (học sinh). 
	} 
\end{bt}

% nguồn: Ok_TuanLuan_Toan11_Lan1
\begin{bt}%[1D5H2-2]%[ma:f328d388]
	 Thời gian truy cập Internet mỗi buổi tối của một số học sinh được cho trong bảng sau
	\begin{center}
		\begin{tabular}{|l|c|c|c|c|c|}
			\hline Thời  gian (phút)&{$[9{,}5;12{,}5)$}&{$[12{,}5;15{,}5)$}&{$[15{,}5;18{,}5)$}&{$[18{,}5;21{,}5)$}&{$[21{,}5;24{,}5)$}\\
			\hline Số  học sinh&$3$&$12$&$15$&$24$&$2$\\
			 \hline
		\end{tabular}
	\end{center}
	Tính trung vị của mẫu số liệu  ghép nhóm này.
	\loigiai{ 
		Cỡ mẫu là $n=3+12+15+24+2=56$.\\ 
		Gọi ${x_1}$,$\ldots$,${x_{56}}$ là thời  gian vào Internet của $56$ học sinh và giả sử dãy này được sắp xếp theo thứ tự tăng dần.\\
		Khi đó, trung vị là  $\dfrac{{x_{28}}+{x_{29}}}{2}$.\\
		Do $2$ giá trị ${x_{28}}$, ${x_{29}}$ thuộc nhóm $[15{,}5;18{,}5)$ nên nhóm  này chứa trung vị.\\
	Do đó, $p=3$, ${a_3}=15{,}5$, ${m_3}=15$, ${m_1}+{m_2}=3+12=15$, ${a_4}-{a_3}=3$ và ta có ${M_e}=15{,}5+\dfrac{\dfrac{56}{2}-15}{15}\cdot  3=18{,}1$.}
\end{bt}

% nguồn: Ok_TuanLuan_Toan11_Lan2
\begin{ex}%[1D5H2-2]%[ma:1a901acd]
	Ghi lại tốc độ bóng trong  $200$ lần giao bóng của một vận động viên môn quần vợt cho kết quả như bảng sau
	\begin{center}
		\begin{tabular}{|c|c|}
			\hline  Tốc độ $v$ (km/h)&Số lần\\
			\hline $150\le  v<155$&$18$\\
			\hline  $155\le v<160$&$28$\\
			 \hline $160\le v<165$&$35$\\
			 \hline $165\le v<170$&$43$\\
			\hline $170\le  v<175$&$41$\\
			\hline $175\le v<180$&$35$\\ 
			 \hline
		\end{tabular}
	\end{center}
	 Tính trung vị của mẫu số liệu ghép nhóm này.
	\loigiai{ 
		 Cỡ mẫu là $n=200$.\\
		Gọi ${x_1},\ldots,{x_{56}}$ là tốc độ giao bóng của $200$ lần giao bóng và giả sử dãy này  được sắp xếp theo thứ tự tăng dần.\\
		Khi  đó, trung vị là $\dfrac{{x_{100}}+{x_{101}}}{2}$.\\
		Do  $2$ giá trị ${x_{100}}$, ${x_{101}}$ thuộc nhóm $[165;170)$ nên nhóm này chứa trung vị.\\
		Do đó, $p=4$, ${a_4}=165$, ${m_3}=43$, ${m_1}+{m_2}+{m_3}=18+28+35=81$,  ${a_5}-{a_4}=5$ và ta có ${M_e}=165+\dfrac{\dfrac{200}{2}-81}{43}\cdot 5=167{,}21$.
	 }
\end{ex}

% nguồn: Ok_TuanLuan_Toan11_Lan2
\begin{ex}%[1D5H2-2]%[ma:cef837e4]
	Điều tra điểm môn Văn của một nhóm học sinh, ta được bảng số liệu như  sau
	\begin{center}
		\begin{tabular}{|l|l|l|l|l|l|l|l|l|}
			\hline  $5$&$6$&$6$&$7$&$7$&$8$&$8$&$8{,}5$&$9$\\
			\hline 
		\end{tabular}
	\end{center}
	Tính số trung vị của bảng nói trên. 
	 \loigiai{
		Ta có $N=9$ 
		 là số lẻ.\\
		Số liệu thứ $\dfrac{N+1}{2}=5$  là số trung vị.\\
	Do đó, số trung  vị là $M_e=7$ (điểm).}
\end{ex}

% nguồn: Ok_TuanLuan_Toan11_Lan2
\begin{ex}%[1D5H2-2]%[ma:bfff1a32]
	Tốc độ phát triển của một loại Virus trong $10$ ngày với các điều kiện khác nhau (đơn vị nghìn  con) được thống kê như sau
	\begin{center}
		\begin{tabular}{|l|l|l|l|l|l|l|l|l|l|}
			\hline $20$&$100$&$30$&$980$&$440$&$20$&$20$&$150$&$60$&$270$\\ 
			\hline 
		\end{tabular}
	\end{center}
	Trong trường hợp này ta chọn số nào dưới đây  là giá trị đại diện là tốt nhất? Tính giá trị đại diện đó.
	 \loigiai{
		Ta chọn số trung vị làm đại diện là tốt nhất vì có sự chênh lệch lớn giữa các số liệu  trong mẫu.\\
		Sắp xếp  lại số liệu mẫu
		\begin{center}
			\begin{tabular}{|c|c|c|c|c|c|c|c|c|c|}
				\hline $20$&$20$&$20$&$30$&$60$&$100$&$150$&$270$&$440$&$980$\\ 
				\hline 
			\end{tabular}
		\end{center}
		Kích thước mẫu là $N=10$. Số liệu thứ $5$ và  $6$ lần lượt là $60$, $100$.\\
		Vậy giá  trị đại diện cho bảng số liệu là ${M_e}=\dfrac{60+100}{2}=80$ (nghìn con).
	} 
\end{ex}

% nguồn: Ok_TuanLuan_Toan11_Lan2
\begin{ex}%[1D5H2-2]%[ma:f387fc5f]
	 Điều tra số học sinh của $30$ lớp học, ta được bảng số liệu như sau
	\begin{center}
		\begin{tabular}{|l|l|l|l|l|l|l|l|l|l|l|l|l|l|l|}
			\hline $35$&$39$&$39$&$40$&$40$&$41$&$41$&$41$&$41$&$44$&$44$&$45$&$45$&$45$&$46$\\ 
			\hline  $48$&$48$&$48$&$48$&$49$&$49$&$49$&$49$&$49$&$49$&$50$&$50$&$50$&$50$&$51$\\
			 \hline
		\end{tabular}
	\end{center}
	Tính số  trung vị của bảng nói trên.
	\loigiai{Ta có  $N=30$ là số chẵn. Số liệu thứ $15$ và $16$ lần lượt là $46$, $48$.\\
	Vậy  số trung vị là ${M_e}=\dfrac{46+48}{2}=47$ (học sinh).}
\end{ex}

% nguồn: Ok_TuanLuan_Toan11_Lan2
\begin{ex}%[1D5H2-2]%[ma:6c6fe0be]
	Điều tra số học sinh giỏi khối $10$ của $15$ trường cấp ba trên địa bàn tỉnh $A$, ta được bảng  số liệu như sau
	\begin{center}
		\begin{tabular}{|l|l|l|l|l|l|l|l|l|l|l|l|l|l|l|}
			 \hline $22$&$29$&$29$&$29$&$30$&$31$&$32$&$32$&$33$&$34$&$34$&$35$&$35$&$35$&$36$\\
			\hline 
		\end{tabular}
	\end{center}
	 Tính số trung vị của bảng nói trên.
	 \loigiai{
		Ta có $N=15$ là số lẻ.  Số liệu thứ $\dfrac{15+1}{2}=8$.\\
		Vậy số trung vị là ${M_e}=32$  (học sinh).
	 }
\end{ex}

% nguồn: Ok_TuanLuan_Toan11_Lan2
\begin{ex}%[1D5H2-2]%[ma:acdeeeb7]
	Thời gian truy cập Internet mỗi buổi tối của  một số học sinh được cho trong bảng sau
	\begin{center}
		\begin{tabular}{|l|c|c|c|c|c|}
			\hline Thời  gian (phút)&{$[9{,}5;12{,}5)$}&{$[12{,}5;15{,}5)$}&{$[15{,}5;18{,}5)$}&{$[18{,}5;21{,}5)$}&{$[21{,}5;24{,}5)$}\\
			\hline Số học  sinh&$3$&$12$&$15$&$24$&$2$\\
			 \hline
		\end{tabular}
	\end{center}
	 Tính trung vị của mẫu số liệu ghép nhóm này.
	\loigiai{ 
		Cỡ mẫu  là $n=3+12+15+24+2=56$.\\
		Gọi ${x_1}$,$\ldots$,${x_{56}}$  là thời gian vào Internet của $56$ học sinh và giả sử dãy này được sắp xếp theo thứ tự tăng dần.\\
		Khi đó, trung  vị là $\dfrac{{x_{28}}+{x_{29}}}{2}$.\\
		Do $2$ giá trị ${x_{28}}$, ${x_{29}}$ thuộc nhóm  $[15{,}5;18{,}5)$ nên nhóm này chứa trung vị.\\
	Do đó, $p=3$, ${a_3}=15{,}5$, ${m_3}=15$, ${m_1}+{m_2}=3+12=15$, ${a_4}-{a_3}=3$ và ta  có ${M_e}=15{,}5+\dfrac{\dfrac{56}{2}-15}{15}\cdot 3=18{,}1$.}
\end{ex}

% nguồn: Ok_TuanLuan_Toan11_Lan2
\begin{ex}%[1D5H2-2]%[ma:2886ea68]
	Thời gian sử dụng  điện thoại trong một ngày của 30 sinh viên được ghi lại ở bảng sau (đơn vị: phút).\[
	\begin{array}{|c|c|c|c|c|c|c|c|c|c|}
		\hline 85 & 195  & 187 & 198 & 43 & 223 & 280 & 71 & 205 & 277 \\
		\hline 298 & 142 &  162 & 89 & 167 & 122 & 175 & 168 & 148 & 253 \\
		\hline 234 & 187 & 85 & 193 & 224 & 233  & 117 & 81 & 39 & 85 \\
		 \hline
	\end{array}
	\] 
	a) Tìm các tứ phân vị của dãy  số liệu trên.\\
	b) Tổng hợp lại dãy số liệu trên vào bảng tần số ghép nhóm với nhóm đầu tiên là ${[0; 60)}$. Hãy ước lượng các tứ phân vị  của mẫu số liệu ghép nhóm
	 \loigiai{
		a)  Sắp xếp mẫu số liệu theo thứ tự không giảm ta được:\\
		$39; 43; 71; 81; 85; 85; 85; 89; 117; 122; 142; 148; 162; 167; 168; 175;$\\ 
		$187; 187; 193; 195; 198; 205; 223; 224; 233; 234; 253; 277; 280; 298$.\\ 
		Cỡ mẫu ${n=30}$ là số chẵn nên trung vị của mẫu số  liệu là\\
		${Q_2=\dfrac{1}{2}\left(x_{15}+x_{16}\right)=\dfrac{1}{2}(168+175)=171,5}$.\\ 
		 Tứ phân vị thứ nhất là ${Q_1=x_8=89}$.\\
		 Tứ phân vị thứ ba là ${Q_3=x_{23}=223}$.\\
		b) Bảng  tần số ghép nhóm của mẫu số liệu như sau:\[
		\begin{array}{|c|c|c|c|c|c|}
			\hline  \text { Thời gian (phút) } & {[0; 60)} & {[60; 120)} & {[120; 180)} & {[180; 240)} & {[240; 300)} \\
			\hline \text { Số sinh viên } & 2 & 7  & 7 & 10 & 4 \\
			\hline 
		\end{array}
		 \]
		Gọi ${x_1; x_2; \ldots; x_{30}}$ là mẫu số liệu được xếp theo thứ tự không  giảm.\\
		Ta có: $x_1, x_2 \in[0; 60)$; $x_3, \ldots, x_9 \in[60; 120)$; $x_{10}, \ldots, x_{16} \in[120; 180)$; $x_{17}, \ldots, x_{26}  \in[180; 240)$; $x_{27}, \ldots, x_{30} \in[240; 300)$.\\
		Tứ phân vị thứ hai của mẫu số liệu ghép nhóm là ${\dfrac{1}{2}\left(x_{15}+x_{16}\right)}$. Do ${x_{15}}$ và ${x_{16}}$ thuộc nhóm ${[120; 180)}$ nên tứ phân vị thứ hai của mẫu số liệu ghép nhóm là\\ 
		${Q_2=120+\dfrac{\dfrac{30}{2}-(2+7)}{7} \cdot(180-120)=\dfrac{1200}{7} \approx  171,43 .
		}$\\ 
		Tứ phân vị thứ nhất của mẫu là  ${x_8 \in[60; 120)}$. Do đó, tứ phân vị thứ nhất của mẫu số liệu ghép nhóm là ${Q_1=60+\dfrac{\dfrac{30}{4}-2}{7} \cdot(120-60)=\dfrac{750}{7} \approx 107,14}$.\\
	Tứ phân vị thứ ba của mẫu số liệu là ${x_{23} \in[180; 240)}$. Do đó, tứ  phân vị thứ ba của mẫu số liệu ghép nhóm là ${Q_3=180+\dfrac{\dfrac{3.30}{4}-(2+7+7)}{10} \cdot(240-180)=219}$.}
\end{ex}

% nguồn: Ok_TuanLuan_Toan11_Lan2
\begin{ex}%[1D5H2-2]%[ma:aa9d97cf]
	Cho mẫu số liệu ghép nhóm về chiều cao của $21$  cây na giống.
	\begin{center}
		\begin{tabular}{|c|c|c|c|c|}
			\hline 
			Chiều  cao (cm) & $[0;5)$ & $[5;10)$ & $[10;15)$ & $[15;20)$ \\
			 \hline
			Số cây & $3$ & $8$ & $7$ & $3$  \\
			 \hline
		\end{tabular}
	\end{center}
	Gọi $x_1$, $x_2$, $\ldots$, $x_{21}$ là chiều cao của các cây giống, đã  được sắp xếp theo thứ tự tăng dần.\\ Khi đó, $x_1$, $\ldots$, $x_3$ thuộc $[0;5)$, $x_4$, $\ldots$, $x_{11}$ thuộc $[5;10)$, $\ldots$. Hỏi trung vị thuộc nhóm nào?
	\loigiai{ 
		Ta có cỡ mẫu $n=21$, là số lẻ nên trung vị là giá trị chính giữa của mẫu số  liệu và là giá trị ở vị trí thứ $11$ của mẫu số liệu.\\
		Mà $x_{11}$ thuộc $[5;10)$ nên  trung vị của mẫu số liệu thuộc nhóm $[5;10)$.
	 }
\end{ex}

% nguồn: Ok_TuanLuan_Toan11_Lan2
\begin{ex}%[1D5H2-2]%[ma:2914644e]
	Một công ty cung cấp nước sạch thống kê lượng nước các hộ gia đình trong một khu vực tiêu thụ  trong một tháng ở bảng sau:\[
	\begin{array}{|c|c|c|c|c|c|}
		\hline \text { Lượng nước tiêu thụ } (\text {m}^3) & {[3; 6)} & {[6; 9)} & {[9; 12)} & {[12; 15)} &  {[15; 18)} \\
		\hline \text { Số hộ gia đình } & 24 & 57 & 42 & 29 &  8 \\
		 \hline
	\end{array}
	 \]
	a) Hãy ước lượng số trung bình, mốt và trung vị của mẫu số liệu  ghép nhóm trên.\\
	b) Công ty muốn gửi  một thông báo khuyến nghị tiết kiệm nước đến ${25 \%}$ các hộ gia đình có lượng nước tiêu thụ cao nhất. Hỏi công ty nên gửi đến các hộ tiêu thụ từ bao nhiêu mét khối nước trở lên?
	\loigiai{ 
		a) Cỡ  mẫu ${n=160}$.\\
		Bảng tần số ghép nhóm của mẫu số liệu trên  như sau:\[
		\begin{array}{|c|c|c|c|c|c|}
			 \hline \text { Lượng nước tiêu thụ } (\text {m}^3) & {[3; 6)} & {[6; 9)} & {[9; 12)} & {[12; 15)} & {[15; 18)} \\
			\hline \text { Giá trị đại diện } & 4,5 & 7,5 & 10,5 &  13,5 & 16,5 \\
			\hline \text { Số hộ gia đình } & 24  & 57 & 42 & 29 & 8 \\
			\hline 
		\end{array}
		 \]
		Số trung bình của mẫu số liệu ghép nhóm đã cho  là\\
		${\bar{x}=\dfrac{1}{160}(24 \cdot 4,5+57 \cdot 7,5+42 \cdot  10,5+29 \cdot 13,5+8 \cdot 16,5)=9,375}$.\\
		Nhóm chứa mốt của mẫu số liệu ghép nhóm  trên là nhóm $\left[6;9 \right)$.\\
		Do đó:  $u_m=6$; $n_m=57$; $n_{m-1}=24$; $n_{m+1}=42$; $u_{m+1}=9$.\\
		 Mốt của mẫu số liệu là ${M_o=6+\dfrac{(57-24)}{(57-24)+(57-42)} \cdot(9-6)=8,0625}$.\\
		Gọi $x_1$; $x_2$; $\ldots$; $x_{160}$  là mẫu số liệu được xếp theo thứ tự không giảm.\\
		Ta có: $x_1$, $\ldots$, $x_{24} \in[3; 6)$; $x_{25}$, $\ldots$, $x_{81} \in[6;  9)$; $x_{82}$, $\ldots$, $x_{123} \in[9; 12)$; $x_{124}$, $\ldots$, $x_{152} \in[12; 15)$; $x_{153}$, $\ldots$, $x_{160} \in[15; 18)$.\\
		Cỡ mẫu  ${n=160}$ là số chẵn nên trung vị là ${M_e=\dfrac{1}{2}\left(x_{80}+x_{81}\right)}$.\\
		Do ${x_{80}}$ và ${x_{81}}$ thuộc nhóm ${[6;  9)}$ nên trung vị của mẫu số liệu là ${M_e=6+\dfrac{\dfrac{160}{2}-24}{57} \cdot(9-6) \approx 8,95}$.\\
		b) ${25 \%}$ các hộ gia đình có lượng nước tiêu thụ cao nhất có lượng nước tiêu thụ không nhỏ hơn $Q_3$, với $Q_3$ là tứ phân vị thứ  ba của mẫu số liệu.\\
		Tứ phân vị thứ ba của dãy số  liệu $x_1$; $x_2$; $\ldots$; $x_{160}$ là ${\dfrac{1}{2}\left(x_{120}+x_{121}\right)}$. Do ${x_{120}}$ và ${x_{121}}$ thuộc nhóm ${[9; 12)}$ nên tứ phân vị thứ ba của mẫu số liệu là ${Q_3=9+\dfrac{\dfrac{160.3}{4}-(24+57)}{42} \cdot(12-9) \approx 11,79}$.\\
	Vậy công ty nên gửi  thông báo tiết kiệm nước đến các hộ gia đình có lượng nước tiêu thụ từ $11,79$ $\text{m}^3$ nước trở lên.}
\end{ex}

% nguồn: Ok_TuanLuan_Toan11_Lan2
\begin{ex}%[1D5H2-2]%[ma:acd1f3b6]
	Quãng đường (km) các cầu thủ (không tính thủ môn) chạy trong một trận bóng đá tại giải ngoại  hạng Anh được cho trong bảng thống kê sau:\[
	\begin{array}{|c|c|c|c|c|c|}
		\hline \text { Quãng đường } & {[2; 4)} & {[4; 6)} & {[6; 8)}  & {[8; 10)} & {[10; 12)} \\
		\hline \text { Số cầu thủ } & 2  & 5 & 6 & 9 & 3 \\
		\hline 
	\end{array}
	\] 
	Tìm trung vị  của mẫu số liệu và giải thích ý nghĩa của giá trị thu được.
	 \loigiai{
		Cỡ mẫu ${n=2+5+6+9+3=25}$. Nhóm chứa trung vị là  ${[6; 8)}$.\\
		Trung vị là $M_e=6+\dfrac{\dfrac{25}{2}-(2+5)}{6}\cdot  (8-6)\approx 7,83$.\\
	Có ${50 \%}$  số cầu thủ chạy nhiều hơn $7,83$ (km) và có ${50 \%}$ số cầu thủ chạy ít hơn $7,83$ (km).}
\end{ex}

% nguồn: Ok_TuanLuan_Toan11_Lan2
\begin{ex}%[1D5H2-2]%[ma:d073352c]
	Bảng sau  thống kê khối lượng một số quả măng cụt được lựa chọn ngẫu nhiên trong một thùng hàng.\[
	\begin{array}{|c|c|c|c|c|c|}
		\hline \text { Khối lượng (gam) } & {[80; 82)} & {[82;  84)} & {[84; 86)} & {[86; 88)} & {[88; 90)} \\
		\hline \text { Số quả } & 18  & 20 & 24 & 15 & 13 \\
		\hline 
	\end{array}
	 \]
	 a) Hãy ước lượng số trung bình, mốt và trung vị của mẫu số liệu ghép nhóm trên.\\
	b) Người ta muốn chia măng cụt trong thùng ra làm ba loại theo cân nặng, bao gồm: loại nhỏ, loại vừa và loại to. Các loại này lần lượt  chiếm khoảng $25 \%$, $50 \%$ và $25 \%$ số măng cụt trong thùng. Hãy xác định ngưỡng cân nặng để phân loại quả.
	\loigiai{ 
		a) Cỡ mẫu ${n=90}$.\\ 
		Bảng tần số  ghép nhóm của mẫu số liệu trên như sau:\[
		\begin{array}{|c|c|c|c|c|c|}
			\hline \text { Khối lượng (gam) } & {[80; 82)} &  {[82; 84)} & {[84; 86)} & {[86; 88)} & {[88; 90)} \\
			\hline \text {  Khối lượng đại diện } & 81 & 83 & 85 & 87 & 89 \\
			\hline \text { Số quả }  & 18 & 20 & 24 & 15 & 13 \\
			\hline 
		\end{array}
		\] 
		Số  trung bình của mẫu số liệu ghép nhóm trên là\\
		${\bar{x}=\dfrac{1}{90}(18 \cdot 81+20 \cdot 83+24 \cdot 85+15 \cdot  87+13 \cdot 89)=\dfrac{254}{3} \approx 84,67}$.\\
		Nhóm chứa mốt của mẫu số liệu là nhóm $\left[84;86  \right).$\\
		 Do đó: $u_m=84$; $n_m=24$; $n_{m-1}=20$; $n_{m+1}=15$; $u_{m+1}=86$.\\
		Mốt của mẫu số liệu là  ${M_o=84+\dfrac{(24-20)}{(24-20)+(24-15)} \cdot (86-84) \approx 84,62}$.\\
		Gọi $x_1$; $x_2$; $\ldots$;  $x_{90}$ là mẫu số liệu được xếp theo thứ tự không giảm.\\
		Ta có: $x_1$, $\ldots$, $x_{18} \in[80; 82)$; $x_{19}$, $\ldots$, $x_{38} \in[82; 84)$; $x_{39}$, $\ldots$, $x_{62} \in[84; 86)$; $x_{63}$, $\ldots$, $x_{77} \in[86; 88)$; $x_{78}$, $\ldots$, $x_{90}  \in[88; 90)$.\\
		Cỡ mẫu ${n=90}$ là số chẵn nên trung  vị là ${M_e=\dfrac{1}{2}\left(x_{45}+x_{46}\right)}$.\\
		Do  ${x_{45}}$ và ${x_{46}}$ thuộc nhóm ${[84; 86)}$ nên trung vị của mẫu số liệu là\\
		${M_e=84+\dfrac{\dfrac{90}{2}-(18+20)}{24} \cdot  (86-84) \approx 84,58
		 }$.\\
		b) Gọi $Q_1$, $Q_3$ lần lượt là tứ phân vị thứ nhất và thứ ba  của mẫu số liệu. Theo đề bài, ta có:\\
		Măng cụt loại nhỏ có cân nặng nhỏ hơn  $Q_1$.\\
		Măng cụt loại vừa  có cân nặng trong ${\left[Q_1; Q_3\right)}$.\\
		Măng cụt loại to có cân nặng không nhỏ  hơn $Q_3$.\\
		Tứ phân vị thứ nhất của  dãy số liệu $x_1$; $x_2$; $\ldots$; $x_{90}$ là ${x_{23} \in[82; 84)}$. Tứ phân vị thứ ba của dãy số liệu $x_1$; $x_2$; $\ldots$; $x_{90}$ là ${x_{68} \in[86; 88)}$.\\
		Do đó, tứ phân vị thứ nhất là ${Q_1=82+\dfrac{\dfrac{90}{4}-18}{20}  \cdot(84-82)=82,45}$;\\
		 tứ phân vị thứ ba là ${Q_3=86+\dfrac{\dfrac{90 \cdot 3}{4}-(18+20+24)}{15} \cdot(88-86) \approx 86,73}$.\\
		Vậy măng cụt loại nhỏ có khối lượng (tính theo gam) thuộc  $\left[80;82,45 \right).$\\
		Măng cụt loại vừa  có khối lượng (tính theo gam) thuộc $[82,45;86,73)$.\\
	Măng cụt loại to có  khối lượng (tính theo gam) thuộc $\left[86,73;90 \right).$}
\end{ex}

% nguồn: Ok_TuanLuan_Toan11_Lan2
\begin{ex}%[1D5H2-2]%[ma:7b66f83a]
	Ghi lại tốc độ bóng trong $200$  lần giao bóng của một vận động viên môn quần vợt cho kết quả như bảng bên.
	\begin{center}
		\begin{tabular}{|c|c|}
			 \hline
			\multicolumn{1}{|c|}{\textbf{Tốc độ $v$ (km/h)}} & \multicolumn{1}{c|}{\textbf{Số lần}}  \\
			 \hline
			\multicolumn{1}{|c|}{$150\leq v<155$}  & \multicolumn{1}{c|}{$18$} \\
			 \hline
			\multicolumn{1}{|c|}{$155\leq  v<160$} & \multicolumn{1}{c|}{$28$} \\
			 \hline
			\multicolumn{1}{|c|}{$160\leq  v<165$} & \multicolumn{1}{c|}{$35$} \\
			\hline 
			\multicolumn{1}{|c|}{$165\leq  v<170$} & \multicolumn{1}{c|}{$43$} \\
			 \hline
			 \multicolumn{1}{|c|}{$170\leq v<175$} & \multicolumn{1}{c|}{$41$} \\
			\hline 
			 \multicolumn{1}{|c|}{$175\leq v<180$} & \multicolumn{1}{c|}{$35$} \\
			\hline 
		\end{tabular}
	\end{center}
	Tính trung vị của mẫu  số liệu ghép nhóm này.
	 \loigiai{
		 Cỡ mẫu là $n=200$.\\
		Gọi $x_1$, $x_2$, $\ldots$, $x_{200}$ là tốc độ giao bóng của vận động viên trong $20$ lần giao bóng và giả sử dãy này đã được sắp xếp theo thứ tự tăng  dần.\\
		Khi đó, trung  vị là $\dfrac{x_{100}+x_{101}}{2}$.\\
		Do $2$  giá trị $x_{100}$, $x_{101}$ thuộc nhóm $[165;170)$ (vì $18+28+35+43=124)$ nên nhóm này chứa trung vị.\\
		Do đó, $p=4$;  $a_4=165$; $m_4=43$; $m_1+m_2+m_3=18+28+35=81$; $a_5-a_4=170-165=5$ và ta có
		\[M_{\mathrm{e}}=165+\dfrac{\dfrac{200}{2}-81}{43}\cdot  5\approx 167{,}21.\]
	} 
\end{ex}

% nguồn: Ok_TuanLuan_Toan11_Lan2
\begin{ex}%[1D5V2-2]%[ma:e5772afc]
	Tìm tứ phân vị thứ nhất và tứ phân vị thứ ba cho mẫu số liệu ghép nhóm ở bảng sau: Ghi lại tốc độ bóng trong 200 lần giao bóng của một vận động viên môn quần vợt cho kết quả như bảng bên. Tính trung vị của mẫu số liệu ghép nhóm này. 
	\begin{center}
		\begin{tabular}{|c|c|}
			\hline Tốc  độ $v$ (km/h) & Số lần \\
			\hline $150 \leq v<155$ & 18 \\ 
			\hline $155 \leq v<160$ &  28 \\
			\hline $160 \leq v<165$ & 35 \\ 
			\hline $165 \leq  v<170$ & 43 \\
			\hline $170 \leq v<175$ & 41  \\
			 \hline $175 \leq v<180$ & 35 \\
			 \hline
		\end{tabular}
	\end{center}
	 \loigiai{
		Cỡ  mẫu: $n=200$.\\
		Tứ phân vị thứ nhất $Q_{1}$ là $\dfrac{x_{50}+x_{51}}{2}$. Do $x_{50}$; $x_{51}$ đều thuộc nhóm [160;165) nên tứ phân vị thứ nhất thuộc nhóm $[160;165)$. Do đó, $p=3$; $a_{3}=160$; $m_{3}=35$; $m_{1}+m_{2}=18+28=46$; $a_{4}-a_{3}=5$  và ta có:
		$Q_{1}=160+\dfrac{\dfrac{200}{4}-46}{35} \times 5=160{,}57$\\ 
		Tứ phân vị thứ ba $Q_{3}$ là $\frac{X_{150}+X_{151}}{2}$. Do $X_{150}, X_{151}$ đều thuộc nhóm [170;175) nên tứ phân vị thứ  ba thuộc nhóm $[170;175)$. \\
		 Do đó,$p=5$; $a_{5}=170$; $m_{5}=41$; $m_{1}+m_{2}+m_{3}+m_{4}=18+28+35+43=124$; $a_{6}-a_{5}=5$
		 và ta có
		\begin{eqnarray*}
			Q_{1}=170+\dfrac{\dfrac{600}{4}-124}{41} \times 5=173.17. 
		\end{eqnarray*}
	 }
\end{ex}

% nguồn: Ok_TuanLuan_Toan11_Lan2
\begin{ex}%[1D5V2-2]%[ma:ec2011ae]
	Tìm tứ phân vị thứ nhất $Q_{1}$ và tứ phân vị thứ ba $Q_{3}$ của mẫu số liệu ghép nhóm cho trong bảng sau: Thời gian (phút) truy cập  Internet mỗi buổi tối của một số học sinh được cho trong bảng sau
	\begin{center}
		\begin{tabular}{|l|c|c|c|c|c|}
			\hline Thời gian (phút) &  $[9{,}5;12{,}5)$ & $[12{,}5;15{,}5)$ & $[15{,}5;18{,}5)$ & $[18{,}5;21{,}5)$ & $[21{,}5;24{,}5)$ \\
			\hline Số học sinh & $3$ & $12$ & $15$ & $24$ & $2$  \\
			\hline 
		\end{tabular}
	\end{center}
	\loigiai{ 
		Cỡ  mẫu là $n=56$.\\
		Tứ phân vị thứ nhất $Q_{1}$ là $\dfrac{x_{14}+x_{15}}{2}$. Do $x_{14}, x_{15}$ đều thuộc nhóm  $[12,5;15,5)$ nên nhóm này chứa $Q_{1}$.\\
		Do đó, $p=2$; $a_{2}=12{,}5$; $m_{2}=12$; $m_{1}=3, a_{3}-a_{2}=3$ và ta  có
		\begin{eqnarray*}
			Q_{1}=12{,}5+\dfrac{\dfrac{56}{4}-3}{12} \cdot  3=15{,}25
		\end{eqnarray*}
		Với tứ phân vị thứ ba $Q_{3}$ là $\dfrac{X_{42}+X_{43}}{2}$. Do $x_{42}, X_{43}$ đều thuộc nhóm $[18,5;21,5)$  nên nhóm này chứa $Q_{3}$.\\
		Do  đó, $p=4$; $a_{4}=18{,}5$; $m_{4}=24$; $m_{1}+m_{2}+m_{3}=3+12+15=30$; $a_{5}-a_{4}=3$ và ta có
		\begin{eqnarray*}
			 Q_{3}=18,5+\dfrac{\dfrac{3.56}{4}-30}{24} \cdot 3=20.
		\end{eqnarray*}
	} 
\end{ex}

% nguồn: Ok_TuanLuan_Toan11_Lan2
\begin{ex}%[1D5V2-2]%[ma:c17ca0f1]
	Một cửa hàng đã ghi lại số tiền bán xăng cho $35$ khách hàng đi xe máy. Mẫu số liệu gốc có dạng: $x_{1}, x_{2}, \ldots, x_{35}$ trong đó $x_{i}$ là số tiền bán xăng  cho khách hàng thứ $i$. Vì một lí do nào đó, cửa hàng chỉ có mẫu số liệu ghép nhóm dạng sau
	\begin{center}
		\begin{tabular}{|l|c|c|c|c|}
			 \hline Số tiền (nghin đồng) & $[0;30)$ & $[30;60)$ & $[60;90)$ & $[90;120)$ \\
			\hline Số khách hàng & 3  & 15 & 10 & 7 \\
			 \hline
		\end{tabular}
	\end{center}
	Hãy tính các số đặc trung cho mẫu số liệu trong bảng trên và giải thích ý nghĩa của các giá trị  thu được.
	\loigiai{ 
		\begin{center}
			\begin{tabular}{|l|r|l|l|r|r|}
				\hline 
				 Số tiền (nghìn đồng) & $[0;30)$ & $[30;60)$ & $[60;90)$ & $[90;120)$ \\
				\hline Số khách hàng  & 3 & 15 & 10 & 7 \\
				\hline Giá trị đại diện & 15 & 45 &  75 & 105 \\
				\hline 
			\end{tabular}
		\end{center}
		Số trung bình là: $\overline{x}=\frac{3 \times  15+15 \times 45+10 \times 75+7 \times 105}{3+15+10+7}=63$.\\
		Cỡ mẫu là: $n=3+15+10+7=35$.\\ 
		\textbf{Ý nghĩa:} Xấp xỉ bằng số trung bình của mẫu số liệu gốc, cho biết vị trí  trung tâm của mẫu số liệu và đại diện cho mẫu số liệu.\\
		 Trung vị là $x_{18}$ thuộc nhóm [30;60), do đó $p=2$; $a_{2}=30$; $m_{2}=15$; $m_{1}=3, a_{3}-a_{2}=30$ và ta có: $M_{e}=30+\dfrac{\dfrac{35}{2}-3}{15} \times 30=59$.\\
		\textbf{Ý nghĩa:} Xấp xỉ bằng trung vị của mẫu số  liệu gốc, chia mẫu số liệu thành 2 phần, mỗi phần chứa $50 \%$ giá trị.\\
		Tứ  phân vị thứ nhất $Q_{1}$ là $x_{9}$ thuộc nhóm [30;60), do đó $p=2$; $a_{2}=30$; $m_{2}=15$; $m_{1}=3$; $a_{3}-a_{2}=30$ và ta có: $Q_{1}=30+\dfrac{\dfrac{35}{4}-3}{15} \times 30=41{,}5$.\\
		\textbf{Ý nghĩa:} Có 25\% số giá trị nhỏ hơn  41{,}5.\\
		Tứ phân vị thứ  ba $Q_{3}$ là $X_{27}$ thuộc nhóm [60;90),do đó $p=3$; $a_{3}=60$; $m_{3}=10$; $m_{1}+m_{2}=3+15=18$; $a_{4}-a_{3}=30$ và ta có:
		$Q_{3}=30+\dfrac{\dfrac{3 \times  35}{4}-18}{10} \times 30=54.75$.\\
		 \textbf{Ý nghĩa:} Có $75 \%$ số giá trị nhỏ hơn 54{,}75\% \\
		Tần số lớn nhất là 15 nên mốt thuộc nhóm [30;60). Ta có, $j=2$; $a_{2}=30$; $m_{2}=15$; $m_{1}=3$;  $m_{3}=10$; $h=30$,do dó: $M_{o}=30+\dfrac{15-3}{(15-3)+(15-10)} \times 30=51{,}18$\\
		\textbf{ Ý nghĩa:} Xấp xỉ bằng mốt của mẫu số liệu gốc, được dùng để đo xu thế trung tâm của mẫu số liệu. 
	 }
\end{ex}

% nguồn: Ok_TuanLuan_Toan11_Lan2
\begin{ex}%[1D5V2-2]%[ma:9c8065c3]
	Thời gian (phút) đề học sinh hoàn thành một câu hỏi thi được  cho như sau
	\begin{center}
		\begin{tabular}{|l|c|c|c|c|c|}
			\hline Thời gian (phút) & $[0{,}5;10{,}5)$ & $[10{,}5;20{,}5)$ &  $[20{,}5;30{,}5)$ & $[30{,}5;40{,}5)$ & $[40{,}5;50{,}5)$ \\
			\hline Số học sinh & 2 & 10 & 6 & 4  & 3 \\
			 \hline
		\end{tabular}
	\end{center}
	Tìm mốt của mẫu số liệu ghép nhóm  này.
	\loigiai{ 
		Tần  số lớn nhất là $10$ nên nhóm chứa mốt là nhóm $[10{,}5; 20{,}5]$. Ta có
		\begin{eqnarray*}
			j=2, a_{2}=10{,}5, m_{2}=10, m_{1}=2; m_{3}=6,  h=10
		\end{eqnarray*}
		Do đó  $M_{o}=10{,}5+\dfrac{10-2}{(10-2)+(10-6)} \times 10=17{,}16$.
	} 
\end{ex}

% nguồn: Ok_TuanLuan_Toan11_Lan2
\begin{ex}%[1D5V2-2]%[ma:7269c3b9]
	Các bạn học sinh lớp 11A1 trả lời 40 câu  hỏi trong một bải kiểm tra. Kết quả được thống kê ở bảng sau
	\begin{center}
		\begin{tabular}{|l|c|c|c|c|c|}
			\hline 
			Số câu trả lời đúng & $[16; 21)$ & $[21; 26)$ & $[26; 31)$ & $[31; 36)$ & $[36; 41)$ \\ 
			\hline 
			Số học sinh & $4$ & $6$ & $8$ & $18$ & $4$ \\ 
			 \hline
		\end{tabular}
	\end{center}
	\begin{listEX}
		\item Tính giá trị đại diện $x_{i}$, $1 \leq i \leq 5$, của từng nhóm  số liệu.
		\item Tính giá trị đại diện $x_{i}$, $1 \leq  i \leq 5$, của từng nhóm số liệu.
		\item Tính $\overline{x}=\dfrac{m_{1} x_{1}+m_{2} x_{2}+m_{3} x_{3}+m_{4}  x_{4}+m_{5} x_{5}}{40}$.
	\end{listEX}
	\loigiai{ 
		\begin{listEX}
			\item $x_{1}=18{,}5$; $x_{2}=23{,}5$;  $x_{3}=28{,}5$; $x_{4}=33{,}5$; $x_{5}=38{,}5$.
			\item $m_{1} x_{1}+m_{2} x_{2}+m_{3} x_{3}+m_{4} x_{4}+m_{5}  x_{5}=1200$.
			\item $\overline{x}=\dfrac{m_{1} x_{1}+m_{2} x_{2}+m_{3} x_{3}+m_{4} x_{4}+m_{5}  x_{5}}{40}=\dfrac{1200}{40}=30$.
		\end{listEX}
	} 
\end{ex}

% nguồn: Ok_TuanLuan_Toan11_Lan2
\begin{ex}%[1D5V2-2]%[ma:4e095a4a]
	Thời gian luyện tập trong  một ngày (tính theo giờ) của một số vận động viên được ghi lại ở bảng sau
	\begin{center}
		\begin{tabular}{|c|c|c|c|c|c|}
			\hline Thời gian luyện tập (giờ) & $[0;2)$ & $[2;4)$ & $[4;6)$ & $[6;8)$ & $[8;10)$ \\ 
			\hline Số vận động viên & $3$ & $8$ & $12$ & $12$ &  $4$ \\
			 \hline
		\end{tabular}
	\end{center}
	Huấn luyện viên muốn xác  định nhóm gồm $25 \%$ các vận động viên có số giờ luyện tập cao nhất. Hỏi huấn luyện viên nên chọn các vận động viên có thời gian luyện tập từ bao nhiêu giờ trở lên vào nhóm này?
	\loigiai{ 
		 Số vận động viên được khảo sát là $n=3+8+12+12+4=39$.\\
		Gọi $x_{1}$; $x_{2}$; $\ldots$; $x_{39}$ là thời gian luyện tập  của $39$ vận động viên được xếp theo thứ tự không giảm.\\
		Ta có $x_{1}$, $x_{2}$, $x_{3} \in[0;2)$; $x_{4}$,  $\ldots$, $x_{11} \in[2;4)$; $x_{12}$, $\ldots$, $x_{23} \in[4;6)$; $x_{24}$, $\ldots$, $x_{35} \in[6;8)$; $x_{36}$, $\ldots$, $x_{39} \in[8;10)$. \\
		Do đó  đối với dãy số liệu $x_{1}$; $x_{2}$; $\ldots$; $x_{39}$ thì
		\begin{itemize}
			\item Tứ phân vị thứ nhất là  $x_{10}$ thuộc nhóm $[2;4)$;
			\item Tứ phân vị thứ  hai là $x_{20}$ thuộc nhóm $[4;6)$;
			\item Tứ phân vị thứ ba là $x_{30}$  thuộc nhóm $[6;8)$.
		\end{itemize}
	Ta nói nhóm $[2;4)$ là nhóm chứa tử phân vị thứ nhất; nhóm $[4;6)$ là nhóm chứa tứ phân  vị thứ hai; nhóm $[6;8)$ là nhóm chứa tứ phân vị thứ ba.}
\end{ex}

% nguồn: Ok_TuanLuan_Toan11_Lan2
\begin{ex}%[1D5V2-2]%[ma:8b0f1811]
	 Ghi lại tốc độ bóng trong $200$ lần giao bóng của một vận động viên môn quần vợt cho kết quả như bảng bên. Tính trung vị của mẫu số liệu ghép nhóm này.
	\begin{center}
		\begin{tabular}{|c|c|}
			\hline Tốc độ  $v$ (km/h) & Số lần \\
			\hline $150 \leq  v<155$ & $18$ \\
			\hline $155 \leq v<160$ & $28$  \\
			\hline $160 \leq v<165$ &  $35$ \\
			\hline $165 \leq v<170$  & $43$ \\
			\hline  $170 \leq v<175$ & $41$ \\
			\hline $175 \leq v<180$  & $35$ \\
			 \hline
		\end{tabular}
	\end{center}
	 \loigiai{
		 Cỡ mẫu là: $n=18+28+35+43+43+41+35=200$.\\
		Gọi $x_{1}, x_{2} \ldots X_{200}$ là tốc độ giao bóng của $200$ lần và  giả sử dãy này được sắp xếp theo thứ tự tăng dần. Khi đó trung vị là $\dfrac{x_{100}+x_{101}}{2}$.\\
		Do hai  giá trị $x_{100}$, $x_{101}$ thuộc nhóm $[165;170)$ nên nhóm này chứa trung vị.\\
		Suy ra, $p=4$; $a_{4}=165$; $m_{4}=43$; $m_{1}+m_{2}+m_{3}=18+28+35=81$;  $a_{5}-a_{4}=5$ và ta có
		\begin{eqnarray*}
			M_{e}=165+\dfrac{\dfrac{200}{2}-81}{43} \times  5=167{,}21.
		\end{eqnarray*}
	 }
\end{ex}

% nguồn: Ok_TuanLuan_Toan11_Lan2
\begin{ex}%[1D5V2-2]%[ma:f908d215]
	Kết quả khảo sát cân nặng của $25$ quả cam ở mỗi lô hàng $A$  và $B$ được cho ở bảng sau
	\begin{center}
		\begin{tabular}{|c|c|c|c|c|c|}
			\hline Cân nặng (g) & $[150;155)$ & $[155;160)$ & $[160;165)$  & $[165;170)$ & $[170;175)$ \\
			\hline Số quả cam ở lô hàng $A$ & $2$ & $6$ & $12$ & $4$ &  $1$ \\
			\hline  Số quả cam ở lô hàng $B$ & $1$ & $3$ & $7$ & $10$ & $4$ \\
			\hline 
		\end{tabular}
	\end{center}
	\begin{listEX}
		\item Hãy ước lượng cân nặng trung  bình của mỗi quả cam ở lô hàng $A$ và lô hàng $B$.
		\item Nếu so sánh theo số trung bình thì cam ở lô hàng nào nặng hơn? 
	\end{listEX}
	\loigiai{ 
		Ta có bảng thống kê số lượng cam theo giá trị đại  diện
		\begin{center}
			\begin{tabular}{|l|c|c|c|c|c|}
				\hline Cân nặng đại diện (g) & $152{,}5$ & $157{,}5$ & $162{,}5$ & $167{,}5$  & $172{,}5$ \\
				\hline Số quả cam ở lô hàng $A$ & $2$ &  $6$ & $12$ & $4$ & $1$ \\
				\hline Số quả cam ở lô hàng $B$ & $1$ & $3$ & $7$ & $10$  & $4$ \\
				 \hline
			\end{tabular}
		\end{center}
		\begin{listEX}
			\item Cân nặng trung bình  của mỗi quả cam ở lô hàng $A$ xấp xỉ bẳng
			\begin{eqnarray*}
				(2\cdot 152{,}5+6 \cdot 157{,}5+12 \cdot 162{,}5+4{,}167{,}5+1 \cdot 172{,}5) :  25=161{,}7\,(text{g}).
			\end{eqnarray*}
			 Cân nặng trung bình của mỗi quả cam ở lô hàng $B$ xấp xỉ bằng
			\begin{eqnarray*}
				(1\cdot152{,}5+3\cdot 157{,}5+7 \cdot 162{,}5+10\cdot 167{,}5+4\cdot  172{,}5) : 25=165{,}1 \,(text{g}).
			\end{eqnarray*}
			\item Nếu so sánh theo số trung bình thì cam ở lô hàng $B$ nặng hơn cam  ở lô hàng $A$.
		\end{listEX}
	} 
\end{ex}

% nguồn: Ok_TuanLuan_Toan11_Lan2
\begin{ex}%[1D5V2-2]%[ma:a9332f64]
	Tổng số lượng mưa trong tháng 8 đo được tại một trạm quan trắc đặt tại Vũng Tàu từ năm 2002 đến năm 2020 được ghi  lại như dưới đây (đơn vị: mm)
	\begin{center}
		\begin{tabular}{cccccccccc}
			$121{,}8$ &  $158{,}3$ & $334{,}9$ & $200{,}9$ & $165{,}6$ & $161{,}5$ & $194{,}3$ & $220{,}7$ & $189{,}8$ & $243{,}2$ \\
			$165{,}9$ & $165{,}9$ & $134$ &  $173$ & $169$ & $189$ & $254$ & $168$ & $255$ &
		\end{tabular}
	\end{center}
	\begin{listEX}
		\item  Xác định số trung bình và tứ phân vị của mẫu số liệu trên.
		\item Hoàn thành bảng tần số  ghép nhóm theo mẫu sau
		\begin{center}
			\begin{tabular}{|c|c|c|c|c|}
				\hline Tőng lượng mưa trong tháng $8$ (mm) & $[120;175)$  & $[175;230)$ & $[230;285)$ & $[285;340)$ \\
				\hline So  năm & $?$ & $?$ & $?$ & $?$ \\
				 \hline
			\end{tabular}
		\end{center}
		\item Hãy ước  lượng số trung bình, tứ phân vị và mốt của mẫu số liệu ở bảng tần số ghép nhóm trên.
	\end{listEX}
	 \loigiai{
		\begin{listEX}
			\item  \textbf{Số trung bình}
			\begin{eqnarray*}
				\overline{x}=\dfrac{121{,}8+158{,}3+334{,}9+200{,}9+165{,}6+161{,}5+194{,}3+220{,}7+189{,}8+\cdots+255}{19} \approx 19288 
			\end{eqnarray*}
			\textbf{Tứ  phân vị}
			Xếp mẫu số liệu  không giảm ta được
			\begin{center}
				\begin{tabular}{cccccccccc}
					$121{,}8$ & $134$  & $158{,}3$ & $161{,}5$ & $165{,}6$ & $165{,}9$ & $165{,}9$ & $168$ & $169$ & $173$ \\
					$189$ & $189{,}8$ & $194{,}3$ & $200{,}9$ & $202{,}7$ &  $243{,}2$ & $254$ & $255$ & $334{,}9$ &
				\end{tabular}
			\end{center}
			Từ đó ta có 
			\begin{itemize}
				\item Tứ phân vị thứ nhất của mẫu  số liệu là: $165{,}6$.
				\item Tứ phân vị thứ hai  của mẫu số liệu là: $173$.
				\item Tứ phân vị thứ ba của mẫu  số liệu là: $202{,}7$.
			\end{itemize}
			\item Hoàn thành bảng  tần số ghép nhóm theo mẫu sau
			\begin{center}
				\begin{tabular}{|c|c|c|c|c|}
					\hline Tổng lượng mưa  trong tháng $8\, (mm)$ & $[120;175)$ & $[175;230)$ & $[230;285)$ & $[285;340)$ \\
					\hline Số năm  & $10$ & $5$ & $3$ & $1$ \\
					\hline 
				\end{tabular}
			\end{center}
			Giá trị đại diện của các  lớp
			$\begin{array}{ll}
			 c_{1}=\dfrac{120+175}{2}=147{,}5;& c_{2}=\dfrac{175+230}{2}=202{,}5 \\
			c_{3}=\dfrac{230+285}{2}=257{,}5;&  c_{4}=\dfrac{285+340}{2}=312{,}5
		\end{array}$ \\
		Tần  số các lớp: $n_{1}=10$; $n_{2}=5$; $n_{3}=3$; $n_{4}=1$.\\
		Số trung bình: $\overline{x}=\dfrac{n_{1} c_{1}+n_{2}  c_{2}+n_{3} c_{3}+n_{4} c_{4}}{n_{1}+n_{2}+n_{3}+n_{4}}=\frac{7145}{38} \approx 188{,}02$.\\
		Tứ phân vị thứ nhất.  Nhóm $[120;175)$
		\begin{eqnarray*}
			Q_{1}=120+\dfrac{\dfrac{1\cdot 19}{4}-0}{10}(175-120)=\frac{1169}{8} \approx 146{,}125. 
		\end{eqnarray*}
		Tứ phân  vị thứ hai. Nhóm $[175; 230)$
		\begin{eqnarray*}
			Q_{2}=175+\dfrac{\dfrac{2\cdot  19}{4}-(0+10)}{5}(230-175)=\frac{339}{2} \approx 169{,}5.
		\end{eqnarray*}
		Tứ  phân vị thứ ba. Nhóm $[230; 285)$
		\begin{eqnarray*}
			 Q_{3}=230+\dfrac{\dfrac{3\cdot 19}{4}-(0+10+5)}{3}(285-230)=\frac{865}{4} \approx 216{,}25.
		\end{eqnarray*}
		 Mốt $M_{0}$ chứa trong nhóm $[120;175)$\\
		 Do đó: $u_{m}=120$; $u_{m+1}=175 \Rightarrow u_{m+1}-u_{m}=175-120=55$\\
		$n_{m-1}=0$; $n_{m}=10$;  $n_{m+1}=5$
		\begin{eqnarray*}
			 M_{0}=120+\dfrac{10-0}{(10-0)+(10-5)}(175-120)=\dfrac{470}{3}=156{,}67.
		\end{eqnarray*}
	\end{listEX}
} 
\end{ex}

% nguồn: Ok_TuanLuan_Toan11_Lan2
\begin{ex}%[1D5V2-2]%[ma:978d2098]
	Số cuộc gọi điện thoại một người thực hiện mỗi ngày trong 30 ngày được lựa chọn ngẫu nhiên  được thống kê trong bảng sau
	\begin{center}
		\begin{tabular}{|l|l|l|l|l|l|}
			\hline Số cuộc gọi & $[3;5]$ &  $[6;8]$ & $[9;11]$ & $[12;14]$ & $[15;17]$ \\
			\hline Số ngày & 5 & 13 & 7 & 3  & 2 \\
			\hline 
		\end{tabular}
	\end{center}
	\begin{listEX}
		\item Tìm mốt của mẫu số liệu  ghép nhóm trên.
		\item Hãy dự đoán xem khả năng người đó thực hiện bao  nhiêu cuộc gọi mỗi ngày là cao nhất.
	\end{listEX}
	 \loigiai{
		Do số cuộc gọi là số  nguyên nên ta hiệu chỉnh lại như sau
		\begin{center}
			\begin{tabular}{|l|l|l|l|l|l|}
				\hline Số cuộc  gọi & $[2{,}5;5{,}5)$ & $[5{,}5;8{,}5)$ & $[8{,}5;11{,}5)$ & $[11{,}5;14{,}5)$ & $[14{,}5;17{,}5)$ \\
				\hline Số ngày & 5 & 13 & 7  & 3 & 2 \\
				 \hline
			\end{tabular}
		\end{center}
		\begin{listEX}
			\item  Nhóm chứa mốt của mẫu số liệu trên là nhóm $[5{,}5;8{,}5)$.\\
			Do đó $u_{m}=5{,}5$; $n_{m-1}=5$;  $n_{m}=13$; $n_{m+1}=7$; $u_{m+1}-u_{m}=8{,}5-5{,}5=3$.\\
			Mốt của mẫu số liệu ghép nhóm là $M_{o}=5,5+\dfrac{13-5}{(13-5)+(13-7)} \cdot 3=\dfrac{101}{14} \approx  7{,}2$.
			\item  Dựa vào kết quả trên ta có thể dự đoán rằng khả năng người đó thực hiện 7 cuộc gọi mỗi ngày là cao nhất.
		\end{listEX}
	} 
\end{ex}

% nguồn: Ok_TuanLuan_Toan11_Lan2
\begin{ex}%[1D5V2-2]%[ma:79d70b65]
	 Thống kê điểm trung bình môn Toán của một số hoc sinh lớp 11 được cho ở bảng sau:
	\begin{center}
		\begin{tabular}{|c|c|c|c|c|c|c|c|}
			\hline Khoảng điểm & $[6{,}5;7)$ & $[7;7{,}5)$ &  $[7{,}5;8)$ & $[8;8{,}5)$ & $[8{,}5;9)$ & $[9;9{,}5)$ & $[9{,}5;10)$ \\
			\hline Tần số & 8 & 10 & 16 & 24 & 13  & 7 & 4 \\
			\hline 
		\end{tabular}
	\end{center}
	Hãy ước lượng số trung  bình, tứ phân vị và mốt của mẫu số liệu ở bảng tần số ghép nhóm trên.
	 \loigiai{
		Bảng tần số ghép nhóm theo giá trị đại  diện là
		\begin{center}
			\begin{tabular}{|c|c|c|c|c|c|c|c|}
				\hline Khoảng điểm  & $[6{,}5;7)$ & $[7;7,{,}5)$ & $[7{,}5;8)$ & $[8;8{,}5)$ & $[8{,}5;9)$ & $[9;9{,}5)$ & $[9{,}5;10)$ \\
				\hline Giá trị đại diện & 6{,}75 & 7{,}25 & 7{,}75 & 8{,}25 & 8{,}75 & 9{,}25 & 9{,}75 \\ 
				\hline Tần số  & 8 & 10 & 16 & 24 & 13 & 7 & 4 \\
				\hline 
			\end{tabular}
		\end{center}
		Điểm trung bình môn Toán của một số hoc sinh lớp  11 là
		\begin{eqnarray*}
			\overline{x}=\dfrac{8\cdot 6{,}75+10\cdot 7{,}25+16\cdot 7{,}75+24\cdot 8{,}25+13\cdot 8{,}75+7\cdot 9{,}25+4\cdot 9{,}75}{82} \approx  8{,}12
		\end{eqnarray*}
		Tứ phân vị  thứ hai. Nhóm $[8;8{,}5)$
		$Q_{2}=8+\dfrac{\dfrac{2\cdot 82}{4}-(8+10+16)}{24}(8;5-8) \approx 8,15$.\\ 
		Tứ phân vị thứ  nhất. Nhóm $[7{,}5;8)$
		$Q_{1}=7,5+\dfrac{\dfrac{2\cdot 82}{4}-(8+10)}{16}(8-7;5) \approx 7,58$.\\ 
		Tứ phân vị thứ ba. Nhóm $[8{,}5;9)$ 
		$Q_{3}=8{,}5+\dfrac{\dfrac{3\cdot 82}{4}-(8+10+16+24)}{16}(9-8;5) \approx  8{,}63$.\\
		Mốt $M_{0}$ chứa trong nhóm $[8;8{,}5)$\\ 
		Do đó: $u_{m}=8$; $u_{m+1}=8{,}5  \Rightarrow u_{m+1}-u_{m}=8{,}5-8=0{,}5$
		$n_{m-1}=16$; $n_{m}=24$; $n_{m+1}=13$ 
		\begin{eqnarray*}
			M_{0}=8+\dfrac{24-16}{(24-16)+(24-13)}\cdot  (8{,}5-8)=8{,}21
		\end{eqnarray*}
	 }
\end{ex}

% nguồn: Ok_TuanLuan_Toan11_Lan2
\begin{ex}%[1D5V2-2]%[ma:162d8949]
	Kết quả khảo sát cân nặng của 25 quả  bơ ở một lô hàng cho trong bảng sau
	\begin{center}
		\begin{tabular}{|c|c|c|c|c|c|}
			\hline Cân nặng (g) & $[150;155)$  & $[155;160)$ & $[160;165)$ & $[165;170)$ & $[170;175)$ \\
			 \hline Số quả bơ & 1 & 7 & 12 & 3 & 2 \\
			\hline 
		\end{tabular}
	\end{center}
	Hãy tìm trung vị của  mẫu số liệu ghép nhóm trên.
	 \loigiai{
		Cỡ mẫu là:  $n=1+7+12+3+2=25$.\\
		 Gọi $x_{1}$; $x_{2}, \ldots, x_{25}$ là cân nặng của 25 quả bơ xếp theo thứ tự không giảm.\\
		Do $x_{1} \in[150;155)$; $x_{2}, \ldots, x_{8} \in[155;160)$; $x_{9}, \ldots, x_{20} \in[160;165)$ nên trung vị của mẫu số liệu $x_{1}$; $x_{2}$; $\ldots$; $x_{25}$ là  $x_{13} \in[160;165)$.\\
		Ta xác định được $n=25,  n_{m}=12, C=1+7=8, u_{m}=160, u_{m+1}=165$.\\
		Suy ra, $p=3$;  $a_{2}=160$; $m_{1}+m_{2}=1+7=8$; $a_{4}-a_{3}=5$ và ta có\\
		Vậy trung vị của mẫu  số liệu ghép nhóm là
		$M_{e}=160+\dfrac{\dfrac{25}{2}-8}{12} \cdot(165-160)=161,875$. 
	 }
\end{ex}

% nguồn: Ok_TuanLuan_Toan11_Lan2
\begin{ex}%[1D5V2-2]%[ma:7a303e88]
	Một công ty xây dựng khảo sát khách hàng xem họ có nhu  cầu mua nhà ở mức giá nào. Kết quả khảo sát được ghi lại ở bảng sau
	\begin{center}
		\begin{tabular}{|l|l|l|l|l|l|}
			\hline Mức  giá (triệu đồng/m²) & $[10;14)$ & $[14;18)$ & $[18;22)$ & $[22;26)$ & $[26;30)$ \\
			\hline Số khách hàng &  $54$ & $78$ & $120$ & $45$ & $12$ \\
			 \hline
		\end{tabular}
	\end{center}
	\begin{listEX}
		\item Tìm mốt của mẫu số  liệu ghép nhóm trên.
		\item Công ty nên xây nhà ở mức giá nào để nhiều  người có nhu cầu mua nhất?
	\end{listEX}
	 \loigiai{
		\begin{listEX}
			\item Nhóm chứa mốt của mẫu  số liệu trên là nhóm [18; 22).\\
			Do  đó $u_{m}=18$, $n_{m-1}=78$, $n_{m}=120$, $n_{m+1}=45$, $u_{m+1}-u_{m}=22-18=4$.\\
			Mốt của mẫu số  liệu ghép nhóm là $M_{o}=18+\dfrac{120-78}{(120-78)+(120-45)} \cdot 4=\dfrac{758}{39} \approx 19{,}4$.
			\item  Dựa vào kết quả trên ta có thể dự đoán rằng nếu công ty xây nhà ở mức giá 19,4 triệu đồng/ $m^{2}$ thì sẽ có nhiều người có nhu cầu mua nhất.
		\end{listEX}
	 }
\end{ex}

% nguồn: Ok_TuanLuan_Toan11_Lan2
\begin{ex}%[1D5V2-2]%[ma:f25d50a6]
	Bảng số liệu ghép nhóm sau cho biết  chiều cao (cm) của 50 học sinh lớp 11A.
	\begin{center}
		\begin{tabular}{|l|c|c|c|c|c|}
			\hline  Khoảng chiều cao $(\mathrm{cm})$ & $[145;150)$ & $[150;155)$ & $[155;160)$ & $[160;165)$ & $[165;170)$ \\
			\hline Số học sinh & 7 & 14 & 10  & 10 & 9 \\
			\hline 
		\end{tabular}
	\end{center}
	Tính mốt của mẫu số liệu  ghép nhóm này. Có thể kết luận gì từ giá trị tính được?
	 \loigiai{
		Tần số lớn nhất là 14 nên nhóm  chứa mốt là nhóm [150; 155). Ta có, $j=2, a_{2}=150, m_{2}=14{,}m_{1}=7, m_{3}=10, h=5$.\\
		Do đó $M_{0}=150+\dfrac{14-7}{(14-7)+(14-10)} \cdot 5  \approx 153{,}18$.
		Số học sinh có chiều cao khoảng $153{,}18$  \, (cm) là nhiều nhất.
	 }
\end{ex}

% nguồn: Ok_TuanLuan_Toan11_Lan2
\begin{ex}%[1D5V2-2]%[ma:1ea94684]
	Để kiểm tra thời gian sả dụng pin của một chiếc điện thoại mới, chị An thống kê  thời gian sử dụng điện thoại của mình từ lúc sạc đầy pin cho đến khi hết pin ở bảng sau
	\begin{center}
		\begin{tabular}{|c|c|c|c|c|c|}
			\hline Thời gian sử dụng (giờ) & $[7;9)$ & $[9;11)$ & $[11;13)$ & $[13,15)$  & $[15;17)$ \\
			 \hline Số lần & $2$ & $5$ & $7$ & $6$ & $3$ \\
			 \hline
		\end{tabular}
	\end{center}
	Hãy ước lượng thời gian sử dụng trung bình từ lúc chị An sạc đầy  pin điện thoại cho tới khi hết pin.
	 \loigiai{
		Thời  gian sử dụng trung bình
		$\overline{x}=\dfrac{2\cdot 8+5\cdot 10+7\cdot 12+6\cdot 14+3\cdot 16}{2+5+7+6+3}  \approx 12{,}26$
	 }
\end{ex}

% nguồn: Ok_TuanLuan_Toan11_Lan2
\begin{ex}%[1D5V2-2]%[ma:aaafe540]
	Bảng sau thống kê số ca nhiễm mới SARS-CoV-2 mỗi  ngày trong tháng 12/2021 tại Việt Nam.
	\begin{center}
		\begin{tabular}{|c|c|c|c|c|c|c|c|}
			\hline Ngày  & Số ca & Ngày & Số ca & Ngày & Số ca & Ngày & Số ca \\
			\hline $1$ & $15139$ & $9$ & $15965$ & $17$  & $15685$ & $25$ & $16046$ \\
			\hline $2$ & $14295$ & $10$ & $15474$ & $18$ & $16363$ & $26$ & $15667$  \\
			\hline $3$ & $14254$ & $11$ & $16830$ &  $19$ & $16586$ & $27$ & $15310$ \\
			\hline $4$ & $14598$ & $12$ & $15264$ &  $20$ & $15420$ & $28$ & $14866$ \\
			\hline $5$ & $14927$ & $13$ & $16035$ & $21$ & $16806$ & $29$ & $14299$  \\
			\hline $6$ & $15215$ & $14$ & $15871$ & $22$ & $17044$ & $30$ &  $20454$ \\
			\hline $7$ & $14433$ & $15$ & $16192$ &  $23$ & $16860$ & $31$ & $17004$ \\
			\hline $8$  & $15223$ & $16$ & $15720$ & $24$ & $16633$ & & \\
			 \hline
		\end{tabular}
	\end{center}
	\begin{listEX}
		 \item Xác định số trung bình và tứ phân vị của mẫu số liệu trên.
		\item Hoàn thành bảng tần  số ghép nhóm theo mẫu sau
		\begin{center}
			\begin{tabular}{|c|c|c|c|c|c|}
				\hline Số ca (nghìn) & $[14;15{,}5)$ & $[15{,}5;17)$ & $[17;18{,}5)$ & $[18{,}5;20)$  & $[20;21{,}5)$ \\
				\hline Số ngày & $?$ & $?$ & $?$  & $?$ & $?$ \\
				\hline 
			\end{tabular}
		\end{center}
		\item Hãy ước lượng số trung bình, tứ phân vị và  mốt của mẫu số liệu ở bảng tần số ghép nhóm trên.
	\end{listEX}
	\loigiai{ 
		\begin{listEX}
			\item Số trung bình:  $\overline{x}=\dfrac{14254+14295+\ldots+20454+17004}{31} \approx 15821$.\\
			Tứ phân  vị
			 Xếp mẫu số liệu không giảm ta được
			\begin{center}
				\begin{tabular}{llllllllll}
					$14254$ &  $14295$ & $14299$ & $14433$ & $14598$ & $14866$ & $14927$ & $15139$ & $15215$ & $15223$ \\
					$15264$ & $15310$ & $15420$ & $15474$ & $15668$ & $15685$ & $15720$  & $15871$ & $15965$ & $16035$ \\
					$16046$ & $16192$ & $16363$ & $16586$ & $16633$ & $16806$ & $16830$ & $16860$ & $17004$ & $17044$ \\ 
					$20454$ & & & & &  & & & &
				\end{tabular}
			\end{center}
			Từ  đó ta có
			\begin{itemize}
				\item Tứ phân vị thứ nhất của mẫu số liệu là $15139$. 
				\item  Tứ phân vị thứ hai của mẫu số liệu là $15685$.
				\item Tứ phân vị thứ ba của mẫu số liệu  là $16586$.
			\end{itemize}
			\item Hoàn thành bảng tần số ghép nhóm  theo mẫu sau
			\begin{center}
				\begin{tabular}{|c|c|c|c|c|c|}
					\hline  Số ca (nghìn) & $[14;15{,}5)$ & $[15{,}5;17)$ & $[17;18{,}5)$ & $[18{,}5;20)$ & $[20;21{,}5)$ \\
					\hline Số ngày &  13 & 15 & 2 & 0 & 1 \\
					 \hline
				\end{tabular}
			\end{center}
			 \item
			Số trung bình: $\overline{x}=\dfrac{14{,}75\cdot 13+16{,}25\cdot 15+17{,}75\cdot 2+19{,}25\cdot 0+20{,}75\cdot 1}{31}=\dfrac{1967}{124} \approx 15{,}86$.\\ 
			Tứ  phân vị thứ nhất. Nhóm $[14;15{,}5)$
			\begin{eqnarray*}
				Q_{1}=14+\dfrac{\dfrac{1\cdot  31}{4}-0}{13}(15{,}5-14)=\dfrac{1549}{104} \approx 14{,}89
			\end{eqnarray*}
			Tứ phân vị thứ hai.  Nhóm $[15{,}5;17)$
			\begin{eqnarray*}
				Q_{2}=15{,}5+\dfrac{\dfrac{2\cdot 31}{4}-(0+13)}{15}(17-15{,}5)=\dfrac{63}{4} \approx  15{,}75
			\end{eqnarray*}
			Tứ  phân vị thứ ba. Nhóm $[17;18{,}5)$
			\begin{eqnarray*}
				Q_{3}=15{,}5+\dfrac{\dfrac{3\cdot  31}{4}-(0+13+15)}{2}(18{,}5-17)=\dfrac{215}{16} \approx 13{,}44
			\end{eqnarray*}
			 Mốt $M_{0}$ chứa trong nhóm $[15{,}5;17)$\\
			Do đó: $u_{m}=15{,}5$;  $u_{m+1}=17 \Rightarrow u_{m+1}-u_{m}=17-15{,}5=1{,}5$
			 $n_{m-1}=13$; $n_{m}=15$; $n_{m+1}=2$
			$M_{0}=15{,}5+\dfrac{15-13}{(15-13)+(15-2)}(17-15{,}5)=\dfrac{157}{10}=15{,}7$ 
		\end{listEX}
	} 
\end{ex}

% nguồn: Ok_TuanLuan_Toan11_Lan2
\begin{ex}%[1D5V2-2]%[ma:5dea0684]
	Xét dữ liệu cho trong tình huống  mở đầu.
	\begin{listEX}
		\item Mẫu số liệu về tổng điểm, kí  hiệu là $(T)$, có bao nhiêu giá trị?
		\item Nếu lập bảng tần số cho mẫu số liệu $(T)$ thì có dễ hình dung được bức tranh  tổng thể về kết quả thi không? Vì sao?
		\item Mẫu số liệu  $(T)$ được mô tả dưới dạng bảng thống kê sau
		\begin{center}
			\begin{tabular}{|l|l|l|l|l|l|l|}
				 \hline Tổng điểm & $<6$ & $[6;7)$ & $[7;8)$ & $\ldots$ & $[28;29)$ & $[29;30]$ \\
				\hline Số thí sinh &  $23$ & $69$ & $192$ & $\ldots$ & $216$ & $12$ \\
				\hline 
			\end{tabular}
		\end{center}
		Hãy đọc và giải thích số liệu được biểu diễn  trong bảng thống kê.
	\end{listEX}
	 \loigiai{
		\begin{listEX}
			\item Có 344752 thí sinh dự  thi nên mẫu số liệu về tổng điểm $(T)$ có $344752$ giá trị.
			\item Nếu lập bảng tần số cho mẫu số liệu ($T$) thì không thể hình dung được bức tranh tổng thể về kết quả thi vì tổng điểm thi 3 ba môn  của các thí sinh có rất nhiều giá trị khác nhau dẫn đến bảng tần số sẽ dài dòng và phức tạp.
			\item Từ bảng thống kê trên ta có thể thấy số lượng thí sinh đạt tổng điểm 3 môn trong từng nhóm điểm. Chẳng hạn 
			\begin{itemize}
				\item Số thí sinh có tổng điểm $3$ môn nhỏ hơn $6$ là  $23$ thí sinh;
				\item Số thí sinh có tổng điểm $3$ môn từ $6$ đến dưới $7$ điểm là $69$ thí sinh; 
				\item Số thí sinh có tổng điểm $3$  môn từ $29$ đến $30$ điểm là $12$ thí sinh.
			\end{itemize}
		\end{listEX}
	 }
\end{ex}

% nguồn: Ok_TuanLuan_Toan11_Lan1
\begin{bt}%[1D5B2-2]%[ma:666027f3]
	Cho hàm số  $y=f(x)=x^3-5x+3$ có đồ thị $(C)$
	\begin{enumerate}
		\item Tính  đạo hàm của hàm số trên.
		\item Viết phương trình  tiếp tuyến của đồ thị $(C)$ tại điểm $M(2;1)$.
	\end{enumerate}
	\loigiai{ 
		\begin{enumerate}
			\item  $f'(x)=(x^3-5x+3)'=3x^2-5$.
			\item Gọi $k$ là hệ số góc của  tiếp tuyến tại $M$. Ta có $k=f'(2)=7$.\\
			Vậy phương trình tiếp  tuyến tại $M$ là
			 $$y=7(x-2)+1=7x-13.$$
		\end{enumerate}
	 }
\end{bt}

% nguồn: Ok_TuanLuan_Toan11_Lan1
\begin{bt}%[1D5B2-2]%[ma:b50aae54]
	Viết phương trình tiếp tuyến của đồ thị hàm số  $y=\dfrac{2x-1}{x-3}$ tại điểm có hoành độ $x=4$.
	\loigiai{ 
		Ta có $y'=\dfrac{(2x-1)'(x-3)-(2x-1)(x-3)'}{(x-3)^2}=\dfrac{-5}{(x-3)^2}\Rightarrow  y'(4)=-5$.\\
		Ta có $y(4)=\dfrac{2\cdot 4-1}{4-3}=7$.\\ 
		Phương trình tiếp  tuyến của đồ thị hàm số tại điểm có hoành độ $x=4$ là
		$$y=y'(4)\cdot  (x-4)+y(4)=-5(x-4)+7=-5x+27.$$
	 }
\end{bt}

% nguồn: Ok_TuanLuan_Toan11_Lan1
\begin{bt}%[1D5B2-2]%[ma:d3a16140]
	Cho hàm số $y=x^3-3x^2+2$ có đồ thị  $(C)$. Viết phương trình tiếp tuyến của đồ thị hàm số $(C)$ tại giao điểm với trục tung.
	\loigiai{Cho  $x=0\Rightarrow y=2$.\\
		Ta có giao  điểm của $(C)$ với trục tung là $M(0;2)$.\\
		 $y'=3x^2-6x\Rightarrow y'(0)=0$.\\
	Phương trình tiếp tuyến  là $y=0(x-0)+2\Leftrightarrow y=2$.}
\end{bt}

% nguồn: Ok_TuanLuan_Toan11_Lan1
\begin{bt}%[1D5K2-2]%[ma:af1292da]
	Cho hàm số $y=-x^4-mx^2+m+1$ có đồ thị $(C)$. Tìm tất cả các giá trị của tham số  $m$ sao cho tiếp tuyến của $(C)$ tại $A(1;0)$ và $B(-1;0)$ vuông góc với nhau.
	 \loigiai{
		Ta có $y'=-4x^3-2mx$,  $\forall x\in \mathbb{R}$.\\
		Tiếp tuyến của $(C)$ tại $A(1;0)$ có hệ số góc là  $k_1=y'(1)=-4-2m$.\\
		Tiếp tuyến của $(C)$ tại  $B(-1;0)$ có hệ số góc là $k_2=y'(-1)=4+2m$.\\
		Tiếp tuyến của $(C)$ tại $A$ và $B$ vuông góc với nhau khi và chỉ khi 
		\[k_1\cdot k_2=-1 \Leftrightarrow (2m+4)^2=1  \Leftrightarrow \hoac{&m=-\dfrac{3}{2}\\&m=-\dfrac{5}{2}.}\]
	} 
\end{bt}

% nguồn: Ok_TuanLuan_Toan11_Lan2
\begin{ex}%[1D5B2-2]%[ma:868bc9b7]
	Cho hàm  số $y=f(x)=x^3-5x+3$ có đồ thị $(C)$
	\begin{enumerate}
		\item Tính đạo hàm của hàm số trên. 
		\item Viết phương trình tiếp tuyến của  đồ thị $(C)$ tại điểm $M(2;1)$.
	\end{enumerate}
	\loigiai{ 
		\begin{enumerate}
			\item  $f'(x)=(x^3-5x+3)'=3x^2-5$.
			 \item Gọi $k$ là hệ số góc của tiếp tuyến tại $M$. Ta có $k=f'(2)=7$.\\
			 Vậy phương trình tiếp tuyến tại $M$ là
			 $$y=7(x-2)+1=7x-13.$$
		\end{enumerate}
	 }
\end{ex}

% nguồn: Ok_TuanLuan_Toan11_Lan2
\begin{ex}%[1D5B2-2]%[ma:8116e2f1]
	Viết phương trình tiếp tuyến của đồ thị hàm số $y=\dfrac{2x-1}{x-3}$ tại điểm có hoành độ $x=4$. 
	\loigiai{ 
		Ta có $y'=\dfrac{(2x-1)'(x-3)-(2x-1)(x-3)'}{(x-3)^2}=\dfrac{-5}{(x-3)^2}\Rightarrow  y'(4)=-5$.\\
		Ta có  $y(4)=\dfrac{2\cdot 4-1}{4-3}=7$.\\
		 Phương trình tiếp tuyến của đồ thị hàm số tại điểm có hoành độ $x=4$ là
		$$y=y'(4)\cdot  (x-4)+y(4)=-5(x-4)+7=-5x+27.$$
	} 
\end{ex}

% nguồn: Ok_TuanLuan_Toan11_Lan2
\begin{ex}%[1D5B2-2]%[ma:6d8ddeb5]
	 Cho hàm số $y=x^3-3x^2+2$ có đồ thị $(C)$. Viết phương trình tiếp tuyến của đồ thị hàm số $(C)$ tại giao điểm với trục tung.
	\loigiai{Cho  $x=0\Rightarrow y=2$.\\
		Ta có giao điểm của $(C)$ với  trục tung là $M(0;2)$.\\
		$y'=3x^2-6x\Rightarrow  y'(0)=0$.\\
	Phương trình tiếp tuyến  là $y=0(x-0)+2\Leftrightarrow y=2$.}
\end{ex}

% nguồn: Ok_TuanLuan_Toan11_Lan2
\begin{ex}%[1D5K2-2]%[ma:f4eb0f1a]
	Cho hàm số $y=-x^4-mx^2+m+1$ có đồ thị $(C)$. Tìm tất cả các giá trị của  tham số $m$ sao cho tiếp tuyến của $(C)$ tại $A(1;0)$ và $B(-1;0)$ vuông góc với nhau.
	\loigiai{ 
		Ta có $y'=-4x^3-2mx$, $\forall x\in \mathbb{R}$.\\ 
		Tiếp tuyến của $(C)$  tại $A(1;0)$ có hệ số góc là $k_1=y'(1)=-4-2m$.\\
		Tiếp tuyến của $(C)$ tại  $B(-1;0)$ có hệ số góc là $k_2=y'(-1)=4+2m$.\\
		Tiếp tuyến của $(C)$ tại $A$ và $B$ vuông góc với  nhau khi và chỉ khi
		\[k_1\cdot k_2=-1 \Leftrightarrow  (2m+4)^2=1 \Leftrightarrow \hoac{&m=-\dfrac{3}{2}\\&m=-\dfrac{5}{2}.}\]
	} 
\end{ex}
